\documentclass[11pt]{amsart}
\usepackage[margin=1.05in]{geometry}
\usepackage{amsmath,amssymb,amsthm,mathtools}
\usepackage{enumitem}
\usepackage{hyperref}
\usepackage{tikz-cd}
\usepackage{mathrsfs}
\usepackage{microtype}
\hypersetup{colorlinks=true,linkcolor=blue,citecolor=teal,urlcolor=purple}
\newtheorem{theorem}{Theorem}[section]
\newtheorem{proposition}[theorem]{Proposition}
\newtheorem{lemma}[theorem]{Lemma}
\newtheorem{corollary}[theorem]{Corollary}
\theoremstyle{definition}
\newtheorem{definition}[theorem]{Definition}
\newtheorem{example}[theorem]{Example}
\newtheorem{exercise}[theorem]{Exercise}
\newtheorem{remark}[theorem]{Remark}
\newcommand{\Sch}{\mathbf{Sch}}
\newcommand{\Set}{\mathbf{Set}}
\newcommand{\Gpd}{\mathbf{Gpd}}
\newcommand{\Spec}{\operatorname{Spec}}
\newcommand{\Hom}{\operatorname{Hom}}
\newcommand{\Aut}{\operatorname{Aut}}
\newcommand{\Pic}{\operatorname{Pic}}
\newcommand{\Hilb}{\operatorname{Hilb}}
\newcommand{\Quot}{\operatorname{Quot}}
\newcommand{\Ext}{\operatorname{Ext}}
\newcommand{\Def}{\operatorname{Def}}
\newcommand{\cM}{\mathcal M}
\newcommand{\cX}{\mathcal X}
\newcommand{\cO}{\mathcal O}
\newcommand{\A}{\mathbb A}
\newcommand{\Pj}{\mathbb P}
\title{A Working Introduction to Moduli Theory\\A Hundred-Page Style Manuscript of Definitions, Theorems, Examples, and Exercises}
\author{Generated study notes for the jiapeng repository}
\date{May 2026}
\begin{document}
\begin{abstract}
These notes give a self-contained, paper-style tour of moduli theory.  The aim is not to replace a specialized monograph, but to provide a long-form manuscript that repeatedly practices the central pattern: formulate a moduli problem as a functor or stack, test whether it is represented, understand automorphisms, construct quotients, and compactify families.  Each page-like sheet contains definitions, statements, examples, and exercises.  The exposition is original and synthetic, while the themes follow the standard mathematical literature on moduli spaces, deformation theory, geometric invariant theory, stable curves, sheaves, and stacks.
\end{abstract}
\maketitle
\tableofcontents
\section*{How to read these notes}
A moduli problem asks for a geometric object whose points classify other geometric objects.  The word ``classify'' hides subtle choices: whether points classify objects over algebraically closed fields, whether families are remembered, whether automorphisms are allowed, and whether degenerations are included.  The guiding slogan is that a moduli space is best understood through its functor of points.  These notes use schemes over a fixed base field $k$ unless another base is displayed.  All exercises are intended as prompts; some are computational, some are proof sketches, and some are invitations to build examples.
\clearpage
\section{Sheet 1: Motivation and functorial language: families}
This sheet studies \emph{families} as it appears in the broader chapter on \emph{motivation and functorial language}.  The recurring question is whether a proposed parameter space remembers enough information about families while remaining geometric enough to compute with.
\begin{definition}[Families viewpoint]
A families moduli functor assigns to each test scheme $T$ the set, groupoid, or category of families over $T$ satisfying the chosen condition.  Two families are identified only when the moduli problem says that an isomorphism, not merely a bijection on fibers, has been supplied.
\end{definition}
\begin{proposition}[Families viewpoint]
If a moduli functor for families is represented by a scheme $M$, then natural transformations from the functor to any functor $G$ are the same as elements of $G(M)$.  In practice this converts universal constructions on families into morphisms out of $M$.
\end{proposition}
\begin{example}[Families viewpoint]
For families, take a base $T=\Spec R$ and replace a single object by an $R$-flat family.  Flatness prevents the numerical invariants from jumping on first-order thickenings, so it is the minimal condition under which the fibers can reasonably be compared.
\end{example}
\begin{theorem}[Families viewpoint]
Suppose the objects in the families problem have no non-trivial automorphisms and admit effective descent for the chosen topology.  Then a universal family, if it exists, is unique up to unique isomorphism.  Consequently the representing object is unique up to unique isomorphism.
\end{theorem}
\begin{remark}[Families viewpoint]
The phrase ``the moduli space of families'' is therefore shorthand.  It may mean a fine moduli scheme, a coarse moduli space, an algebraic stack, a GIT quotient, or a compactification.  The correct meaning is determined by what information is retained about families and automorphisms.
\end{remark}
\begin{proof}[Proof sketch of the theorem/proposition on this sheet]
The argument is the standard Yoneda-and-descent argument specialized to families.  A family over a test scheme $T$ corresponds, under representability, to a morphism $T\to M$.  Pulling the universal family back along this morphism recovers the given family.  If two morphisms give isomorphic pullbacks, the functorial equality on all test schemes forces the morphisms to agree.  The hypotheses on automorphisms ensure that the equality is not weakened to an equivalence relation in a groupoid.
\end{proof}
\begin{example}[A test calculation for families]
Let $T=\Spec k[s]$ and let the invariant of interest be denoted $\nu$.  A toy family $X\to T$ is chosen so that $\nu(X_t)=\nu_0$ for $s\ne0$ and the special fiber at $s=0$ records the boundary behavior.  In a concrete course one may take $X$ to be a family of divisors, modules, curves, bundles, or representations, depending on the chapter.  The moduli-theoretic lesson is that the morphism $T\to M$ detects both the general member and the limiting point.
\end{example}
\begin{exercise}
Write the functor of points for the moduli problem on this sheet and identify its value on $\Spec k$ and on $\Spec k[\epsilon]/(\epsilon^2)$.
\end{exercise}
\begin{exercise}
Formulate a small example in which the naive set of isomorphism classes is not Hausdorff, not separated, or not functorial.  Then repair the example by adding stability, rigidification, or a stack structure.
\end{exercise}
\clearpage
\section{Sheet 2: Motivation and functorial language: isomorphism classes}
This sheet studies \emph{isomorphism classes} as it appears in the broader chapter on \emph{motivation and functorial language}.  The recurring question is whether a proposed parameter space remembers enough information about families while remaining geometric enough to compute with.
\begin{definition}[Isomorphism classes viewpoint]
A isomorphism classes moduli functor assigns to each test scheme $T$ the set, groupoid, or category of families over $T$ satisfying the chosen condition.  Two families are identified only when the moduli problem says that an isomorphism, not merely a bijection on fibers, has been supplied.
\end{definition}
\begin{proposition}[Isomorphism classes viewpoint]
If a moduli functor for isomorphism classes is represented by a scheme $M$, then natural transformations from the functor to any functor $G$ are the same as elements of $G(M)$.  In practice this converts universal constructions on families into morphisms out of $M$.
\end{proposition}
\begin{example}[Isomorphism classes viewpoint]
For isomorphism classes, take a base $T=\Spec R$ and replace a single object by an $R$-flat family.  Flatness prevents the numerical invariants from jumping on first-order thickenings, so it is the minimal condition under which the fibers can reasonably be compared.
\end{example}
\begin{theorem}[Isomorphism classes viewpoint]
Suppose the objects in the isomorphism classes problem have no non-trivial automorphisms and admit effective descent for the chosen topology.  Then a universal family, if it exists, is unique up to unique isomorphism.  Consequently the representing object is unique up to unique isomorphism.
\end{theorem}
\begin{remark}[Isomorphism classes viewpoint]
The phrase ``the moduli space of isomorphism classes'' is therefore shorthand.  It may mean a fine moduli scheme, a coarse moduli space, an algebraic stack, a GIT quotient, or a compactification.  The correct meaning is determined by what information is retained about families and automorphisms.
\end{remark}
\begin{proof}[Proof sketch of the theorem/proposition on this sheet]
The argument is the standard Yoneda-and-descent argument specialized to isomorphism classes.  A family over a test scheme $T$ corresponds, under representability, to a morphism $T\to M$.  Pulling the universal family back along this morphism recovers the given family.  If two morphisms give isomorphic pullbacks, the functorial equality on all test schemes forces the morphisms to agree.  The hypotheses on automorphisms ensure that the equality is not weakened to an equivalence relation in a groupoid.
\end{proof}
\begin{example}[A test calculation for isomorphism classes]
Let $T=\Spec k[s]$ and let the invariant of interest be denoted $\nu$.  A toy family $X\to T$ is chosen so that $\nu(X_t)=\nu_0$ for $s\ne0$ and the special fiber at $s=0$ records the boundary behavior.  In a concrete course one may take $X$ to be a family of divisors, modules, curves, bundles, or representations, depending on the chapter.  The moduli-theoretic lesson is that the morphism $T\to M$ detects both the general member and the limiting point.
\end{example}
\begin{exercise}
Find one automorphism group occurring in the example and decide whether it obstructs the existence of a fine moduli scheme.
\end{exercise}
\begin{exercise}
Formulate a small example in which the naive set of isomorphism classes is not Hausdorff, not separated, or not functorial.  Then repair the example by adding stability, rigidification, or a stack structure.
\end{exercise}
\clearpage
\section{Sheet 3: Motivation and functorial language: base change}
This sheet studies \emph{base change} as it appears in the broader chapter on \emph{motivation and functorial language}.  The recurring question is whether a proposed parameter space remembers enough information about families while remaining geometric enough to compute with.
\begin{definition}[Base change viewpoint]
A base change moduli functor assigns to each test scheme $T$ the set, groupoid, or category of families over $T$ satisfying the chosen condition.  Two families are identified only when the moduli problem says that an isomorphism, not merely a bijection on fibers, has been supplied.
\end{definition}
\begin{proposition}[Base change viewpoint]
If a moduli functor for base change is represented by a scheme $M$, then natural transformations from the functor to any functor $G$ are the same as elements of $G(M)$.  In practice this converts universal constructions on families into morphisms out of $M$.
\end{proposition}
\begin{example}[Base change viewpoint]
For base change, take a base $T=\Spec R$ and replace a single object by an $R$-flat family.  Flatness prevents the numerical invariants from jumping on first-order thickenings, so it is the minimal condition under which the fibers can reasonably be compared.
\end{example}
\begin{theorem}[Base change viewpoint]
Suppose the objects in the base change problem have no non-trivial automorphisms and admit effective descent for the chosen topology.  Then a universal family, if it exists, is unique up to unique isomorphism.  Consequently the representing object is unique up to unique isomorphism.
\end{theorem}
\begin{remark}[Base change viewpoint]
The phrase ``the moduli space of base change'' is therefore shorthand.  It may mean a fine moduli scheme, a coarse moduli space, an algebraic stack, a GIT quotient, or a compactification.  The correct meaning is determined by what information is retained about families and automorphisms.
\end{remark}
\begin{proof}[Proof sketch of the theorem/proposition on this sheet]
The argument is the standard Yoneda-and-descent argument specialized to base change.  A family over a test scheme $T$ corresponds, under representability, to a morphism $T\to M$.  Pulling the universal family back along this morphism recovers the given family.  If two morphisms give isomorphic pullbacks, the functorial equality on all test schemes forces the morphisms to agree.  The hypotheses on automorphisms ensure that the equality is not weakened to an equivalence relation in a groupoid.
\end{proof}
\begin{example}[A test calculation for base change]
Let $T=\Spec k[s]$ and let the invariant of interest be denoted $\nu$.  A toy family $X\to T$ is chosen so that $\nu(X_t)=\nu_0$ for $s\ne0$ and the special fiber at $s=0$ records the boundary behavior.  In a concrete course one may take $X$ to be a family of divisors, modules, curves, bundles, or representations, depending on the chapter.  The moduli-theoretic lesson is that the morphism $T\to M$ detects both the general member and the limiting point.
\end{example}
\begin{exercise}
Check a base-change square explicitly: pull the displayed family back along $T\prime\to T$ and verify which invariant is preserved.
\end{exercise}
\begin{exercise}
Formulate a small example in which the naive set of isomorphism classes is not Hausdorff, not separated, or not functorial.  Then repair the example by adding stability, rigidification, or a stack structure.
\end{exercise}
\clearpage
\section{Sheet 4: Motivation and functorial language: universal objects}
This sheet studies \emph{universal objects} as it appears in the broader chapter on \emph{motivation and functorial language}.  The recurring question is whether a proposed parameter space remembers enough information about families while remaining geometric enough to compute with.
\begin{definition}[Universal objects viewpoint]
A universal objects moduli functor assigns to each test scheme $T$ the set, groupoid, or category of families over $T$ satisfying the chosen condition.  Two families are identified only when the moduli problem says that an isomorphism, not merely a bijection on fibers, has been supplied.
\end{definition}
\begin{proposition}[Universal objects viewpoint]
If a moduli functor for universal objects is represented by a scheme $M$, then natural transformations from the functor to any functor $G$ are the same as elements of $G(M)$.  In practice this converts universal constructions on families into morphisms out of $M$.
\end{proposition}
\begin{example}[Universal objects viewpoint]
For universal objects, take a base $T=\Spec R$ and replace a single object by an $R$-flat family.  Flatness prevents the numerical invariants from jumping on first-order thickenings, so it is the minimal condition under which the fibers can reasonably be compared.
\end{example}
\begin{theorem}[Universal objects viewpoint]
Suppose the objects in the universal objects problem have no non-trivial automorphisms and admit effective descent for the chosen topology.  Then a universal family, if it exists, is unique up to unique isomorphism.  Consequently the representing object is unique up to unique isomorphism.
\end{theorem}
\begin{remark}[Universal objects viewpoint]
The phrase ``the moduli space of universal objects'' is therefore shorthand.  It may mean a fine moduli scheme, a coarse moduli space, an algebraic stack, a GIT quotient, or a compactification.  The correct meaning is determined by what information is retained about families and automorphisms.
\end{remark}
\begin{proof}[Proof sketch of the theorem/proposition on this sheet]
The argument is the standard Yoneda-and-descent argument specialized to universal objects.  A family over a test scheme $T$ corresponds, under representability, to a morphism $T\to M$.  Pulling the universal family back along this morphism recovers the given family.  If two morphisms give isomorphic pullbacks, the functorial equality on all test schemes forces the morphisms to agree.  The hypotheses on automorphisms ensure that the equality is not weakened to an equivalence relation in a groupoid.
\end{proof}
\begin{example}[A test calculation for universal objects]
Let $T=\Spec k[s]$ and let the invariant of interest be denoted $\nu$.  A toy family $X\to T$ is chosen so that $\nu(X_t)=\nu_0$ for $s\ne0$ and the special fiber at $s=0$ records the boundary behavior.  In a concrete course one may take $X$ to be a family of divisors, modules, curves, bundles, or representations, depending on the chapter.  The moduli-theoretic lesson is that the morphism $T\to M$ detects both the general member and the limiting point.
\end{example}
\begin{exercise}
Construct a degeneration over $\Spec k[[t]]$ whose generic fiber satisfies the condition but whose special fiber forces a compactification.
\end{exercise}
\begin{exercise}
Formulate a small example in which the naive set of isomorphism classes is not Hausdorff, not separated, or not functorial.  Then repair the example by adding stability, rigidification, or a stack structure.
\end{exercise}
\clearpage
\section{Sheet 5: Motivation and functorial language: fine and coarse moduli}
This sheet studies \emph{fine and coarse moduli} as it appears in the broader chapter on \emph{motivation and functorial language}.  The recurring question is whether a proposed parameter space remembers enough information about families while remaining geometric enough to compute with.
\begin{definition}[Fine and coarse moduli viewpoint]
A fine and coarse moduli moduli functor assigns to each test scheme $T$ the set, groupoid, or category of families over $T$ satisfying the chosen condition.  Two families are identified only when the moduli problem says that an isomorphism, not merely a bijection on fibers, has been supplied.
\end{definition}
\begin{proposition}[Fine and coarse moduli viewpoint]
If a moduli functor for fine and coarse moduli is represented by a scheme $M$, then natural transformations from the functor to any functor $G$ are the same as elements of $G(M)$.  In practice this converts universal constructions on families into morphisms out of $M$.
\end{proposition}
\begin{example}[Fine and coarse moduli viewpoint]
For fine and coarse moduli, take a base $T=\Spec R$ and replace a single object by an $R$-flat family.  Flatness prevents the numerical invariants from jumping on first-order thickenings, so it is the minimal condition under which the fibers can reasonably be compared.
\end{example}
\begin{theorem}[Fine and coarse moduli viewpoint]
Suppose the objects in the fine and coarse moduli problem have no non-trivial automorphisms and admit effective descent for the chosen topology.  Then a universal family, if it exists, is unique up to unique isomorphism.  Consequently the representing object is unique up to unique isomorphism.
\end{theorem}
\begin{remark}[Fine and coarse moduli viewpoint]
The phrase ``the moduli space of fine and coarse moduli'' is therefore shorthand.  It may mean a fine moduli scheme, a coarse moduli space, an algebraic stack, a GIT quotient, or a compactification.  The correct meaning is determined by what information is retained about families and automorphisms.
\end{remark}
\begin{proof}[Proof sketch of the theorem/proposition on this sheet]
The argument is the standard Yoneda-and-descent argument specialized to fine and coarse moduli.  A family over a test scheme $T$ corresponds, under representability, to a morphism $T\to M$.  Pulling the universal family back along this morphism recovers the given family.  If two morphisms give isomorphic pullbacks, the functorial equality on all test schemes forces the morphisms to agree.  The hypotheses on automorphisms ensure that the equality is not weakened to an equivalence relation in a groupoid.
\end{proof}
\begin{example}[A test calculation for fine and coarse moduli]
Let $T=\Spec k[s]$ and let the invariant of interest be denoted $\nu$.  A toy family $X\to T$ is chosen so that $\nu(X_t)=\nu_0$ for $s\ne0$ and the special fiber at $s=0$ records the boundary behavior.  In a concrete course one may take $X$ to be a family of divisors, modules, curves, bundles, or representations, depending on the chapter.  The moduli-theoretic lesson is that the morphism $T\to M$ detects both the general member and the limiting point.
\end{example}
\begin{exercise}
Compare the functor with the corresponding groupoid-valued functor.  Which information is lost after passing to isomorphism classes?
\end{exercise}
\begin{exercise}
Formulate a small example in which the naive set of isomorphism classes is not Hausdorff, not separated, or not functorial.  Then repair the example by adding stability, rigidification, or a stack structure.
\end{exercise}
\clearpage
\section{Sheet 6: Representable functors: Yoneda lemma}
This sheet studies \emph{Yoneda lemma} as it appears in the broader chapter on \emph{representable functors}.  The recurring question is whether a proposed parameter space remembers enough information about families while remaining geometric enough to compute with.
\begin{definition}[Yoneda lemma viewpoint]
A Yoneda lemma moduli functor assigns to each test scheme $T$ the set, groupoid, or category of families over $T$ satisfying the chosen condition.  Two families are identified only when the moduli problem says that an isomorphism, not merely a bijection on fibers, has been supplied.
\end{definition}
\begin{proposition}[Yoneda lemma viewpoint]
If a moduli functor for Yoneda lemma is represented by a scheme $M$, then natural transformations from the functor to any functor $G$ are the same as elements of $G(M)$.  In practice this converts universal constructions on families into morphisms out of $M$.
\end{proposition}
\begin{example}[Yoneda lemma viewpoint]
For Yoneda lemma, take a base $T=\Spec R$ and replace a single object by an $R$-flat family.  Flatness prevents the numerical invariants from jumping on first-order thickenings, so it is the minimal condition under which the fibers can reasonably be compared.
\end{example}
\begin{theorem}[Yoneda lemma viewpoint]
Suppose the objects in the Yoneda lemma problem have no non-trivial automorphisms and admit effective descent for the chosen topology.  Then a universal family, if it exists, is unique up to unique isomorphism.  Consequently the representing object is unique up to unique isomorphism.
\end{theorem}
\begin{remark}[Yoneda lemma viewpoint]
The phrase ``the moduli space of Yoneda lemma'' is therefore shorthand.  It may mean a fine moduli scheme, a coarse moduli space, an algebraic stack, a GIT quotient, or a compactification.  The correct meaning is determined by what information is retained about families and automorphisms.
\end{remark}
\begin{proof}[Proof sketch of the theorem/proposition on this sheet]
The argument is the standard Yoneda-and-descent argument specialized to Yoneda lemma.  A family over a test scheme $T$ corresponds, under representability, to a morphism $T\to M$.  Pulling the universal family back along this morphism recovers the given family.  If two morphisms give isomorphic pullbacks, the functorial equality on all test schemes forces the morphisms to agree.  The hypotheses on automorphisms ensure that the equality is not weakened to an equivalence relation in a groupoid.
\end{proof}
\begin{example}[A test calculation for Yoneda lemma]
Let $T=\Spec k[s]$ and let the invariant of interest be denoted $\nu$.  A toy family $X\to T$ is chosen so that $\nu(X_t)=\nu_0$ for $s\ne0$ and the special fiber at $s=0$ records the boundary behavior.  In a concrete course one may take $X$ to be a family of divisors, modules, curves, bundles, or representations, depending on the chapter.  The moduli-theoretic lesson is that the morphism $T\to M$ detects both the general member and the limiting point.
\end{example}
\begin{exercise}
Write the functor of points for the moduli problem on this sheet and identify its value on $\Spec k$ and on $\Spec k[\epsilon]/(\epsilon^2)$.
\end{exercise}
\begin{exercise}
Formulate a small example in which the naive set of isomorphism classes is not Hausdorff, not separated, or not functorial.  Then repair the example by adding stability, rigidification, or a stack structure.
\end{exercise}
\clearpage
\section{Sheet 7: Representable functors: affine schemes}
This sheet studies \emph{affine schemes} as it appears in the broader chapter on \emph{representable functors}.  The recurring question is whether a proposed parameter space remembers enough information about families while remaining geometric enough to compute with.
\begin{definition}[Affine schemes viewpoint]
A affine schemes moduli functor assigns to each test scheme $T$ the set, groupoid, or category of families over $T$ satisfying the chosen condition.  Two families are identified only when the moduli problem says that an isomorphism, not merely a bijection on fibers, has been supplied.
\end{definition}
\begin{proposition}[Affine schemes viewpoint]
If a moduli functor for affine schemes is represented by a scheme $M$, then natural transformations from the functor to any functor $G$ are the same as elements of $G(M)$.  In practice this converts universal constructions on families into morphisms out of $M$.
\end{proposition}
\begin{example}[Affine schemes viewpoint]
For affine schemes, take a base $T=\Spec R$ and replace a single object by an $R$-flat family.  Flatness prevents the numerical invariants from jumping on first-order thickenings, so it is the minimal condition under which the fibers can reasonably be compared.
\end{example}
\begin{theorem}[Affine schemes viewpoint]
Suppose the objects in the affine schemes problem have no non-trivial automorphisms and admit effective descent for the chosen topology.  Then a universal family, if it exists, is unique up to unique isomorphism.  Consequently the representing object is unique up to unique isomorphism.
\end{theorem}
\begin{remark}[Affine schemes viewpoint]
The phrase ``the moduli space of affine schemes'' is therefore shorthand.  It may mean a fine moduli scheme, a coarse moduli space, an algebraic stack, a GIT quotient, or a compactification.  The correct meaning is determined by what information is retained about families and automorphisms.
\end{remark}
\begin{proof}[Proof sketch of the theorem/proposition on this sheet]
The argument is the standard Yoneda-and-descent argument specialized to affine schemes.  A family over a test scheme $T$ corresponds, under representability, to a morphism $T\to M$.  Pulling the universal family back along this morphism recovers the given family.  If two morphisms give isomorphic pullbacks, the functorial equality on all test schemes forces the morphisms to agree.  The hypotheses on automorphisms ensure that the equality is not weakened to an equivalence relation in a groupoid.
\end{proof}
\begin{example}[A test calculation for affine schemes]
Let $T=\Spec k[s]$ and let the invariant of interest be denoted $\nu$.  A toy family $X\to T$ is chosen so that $\nu(X_t)=\nu_0$ for $s\ne0$ and the special fiber at $s=0$ records the boundary behavior.  In a concrete course one may take $X$ to be a family of divisors, modules, curves, bundles, or representations, depending on the chapter.  The moduli-theoretic lesson is that the morphism $T\to M$ detects both the general member and the limiting point.
\end{example}
\begin{exercise}
Find one automorphism group occurring in the example and decide whether it obstructs the existence of a fine moduli scheme.
\end{exercise}
\begin{exercise}
Formulate a small example in which the naive set of isomorphism classes is not Hausdorff, not separated, or not functorial.  Then repair the example by adding stability, rigidification, or a stack structure.
\end{exercise}
\clearpage
\section{Sheet 8: Representable functors: Grassmannians}
This sheet studies \emph{Grassmannians} as it appears in the broader chapter on \emph{representable functors}.  The recurring question is whether a proposed parameter space remembers enough information about families while remaining geometric enough to compute with.
\begin{definition}[Grassmannians viewpoint]
A Grassmannians moduli functor assigns to each test scheme $T$ the set, groupoid, or category of families over $T$ satisfying the chosen condition.  Two families are identified only when the moduli problem says that an isomorphism, not merely a bijection on fibers, has been supplied.
\end{definition}
\begin{proposition}[Grassmannians viewpoint]
If a moduli functor for Grassmannians is represented by a scheme $M$, then natural transformations from the functor to any functor $G$ are the same as elements of $G(M)$.  In practice this converts universal constructions on families into morphisms out of $M$.
\end{proposition}
\begin{example}[Grassmannians viewpoint]
For Grassmannians, take a base $T=\Spec R$ and replace a single object by an $R$-flat family.  Flatness prevents the numerical invariants from jumping on first-order thickenings, so it is the minimal condition under which the fibers can reasonably be compared.
\end{example}
\begin{theorem}[Grassmannians viewpoint]
Suppose the objects in the Grassmannians problem have no non-trivial automorphisms and admit effective descent for the chosen topology.  Then a universal family, if it exists, is unique up to unique isomorphism.  Consequently the representing object is unique up to unique isomorphism.
\end{theorem}
\begin{remark}[Grassmannians viewpoint]
The phrase ``the moduli space of Grassmannians'' is therefore shorthand.  It may mean a fine moduli scheme, a coarse moduli space, an algebraic stack, a GIT quotient, or a compactification.  The correct meaning is determined by what information is retained about families and automorphisms.
\end{remark}
\begin{proof}[Proof sketch of the theorem/proposition on this sheet]
The argument is the standard Yoneda-and-descent argument specialized to Grassmannians.  A family over a test scheme $T$ corresponds, under representability, to a morphism $T\to M$.  Pulling the universal family back along this morphism recovers the given family.  If two morphisms give isomorphic pullbacks, the functorial equality on all test schemes forces the morphisms to agree.  The hypotheses on automorphisms ensure that the equality is not weakened to an equivalence relation in a groupoid.
\end{proof}
\begin{example}[A test calculation for Grassmannians]
Let $T=\Spec k[s]$ and let the invariant of interest be denoted $\nu$.  A toy family $X\to T$ is chosen so that $\nu(X_t)=\nu_0$ for $s\ne0$ and the special fiber at $s=0$ records the boundary behavior.  In a concrete course one may take $X$ to be a family of divisors, modules, curves, bundles, or representations, depending on the chapter.  The moduli-theoretic lesson is that the morphism $T\to M$ detects both the general member and the limiting point.
\end{example}
\begin{exercise}
Check a base-change square explicitly: pull the displayed family back along $T\prime\to T$ and verify which invariant is preserved.
\end{exercise}
\begin{exercise}
Formulate a small example in which the naive set of isomorphism classes is not Hausdorff, not separated, or not functorial.  Then repair the example by adding stability, rigidification, or a stack structure.
\end{exercise}
\clearpage
\section{Sheet 9: Representable functors: Hilbert schemes}
This sheet studies \emph{Hilbert schemes} as it appears in the broader chapter on \emph{representable functors}.  The recurring question is whether a proposed parameter space remembers enough information about families while remaining geometric enough to compute with.
\begin{definition}[Hilbert schemes viewpoint]
A Hilbert schemes moduli functor assigns to each test scheme $T$ the set, groupoid, or category of families over $T$ satisfying the chosen condition.  Two families are identified only when the moduli problem says that an isomorphism, not merely a bijection on fibers, has been supplied.
\end{definition}
\begin{proposition}[Hilbert schemes viewpoint]
If a moduli functor for Hilbert schemes is represented by a scheme $M$, then natural transformations from the functor to any functor $G$ are the same as elements of $G(M)$.  In practice this converts universal constructions on families into morphisms out of $M$.
\end{proposition}
\begin{example}[Hilbert schemes viewpoint]
For Hilbert schemes, take a base $T=\Spec R$ and replace a single object by an $R$-flat family.  Flatness prevents the numerical invariants from jumping on first-order thickenings, so it is the minimal condition under which the fibers can reasonably be compared.
\end{example}
\begin{theorem}[Hilbert schemes viewpoint]
Suppose the objects in the Hilbert schemes problem have no non-trivial automorphisms and admit effective descent for the chosen topology.  Then a universal family, if it exists, is unique up to unique isomorphism.  Consequently the representing object is unique up to unique isomorphism.
\end{theorem}
\begin{remark}[Hilbert schemes viewpoint]
The phrase ``the moduli space of Hilbert schemes'' is therefore shorthand.  It may mean a fine moduli scheme, a coarse moduli space, an algebraic stack, a GIT quotient, or a compactification.  The correct meaning is determined by what information is retained about families and automorphisms.
\end{remark}
\begin{proof}[Proof sketch of the theorem/proposition on this sheet]
The argument is the standard Yoneda-and-descent argument specialized to Hilbert schemes.  A family over a test scheme $T$ corresponds, under representability, to a morphism $T\to M$.  Pulling the universal family back along this morphism recovers the given family.  If two morphisms give isomorphic pullbacks, the functorial equality on all test schemes forces the morphisms to agree.  The hypotheses on automorphisms ensure that the equality is not weakened to an equivalence relation in a groupoid.
\end{proof}
\begin{example}[A test calculation for Hilbert schemes]
Let $T=\Spec k[s]$ and let the invariant of interest be denoted $\nu$.  A toy family $X\to T$ is chosen so that $\nu(X_t)=\nu_0$ for $s\ne0$ and the special fiber at $s=0$ records the boundary behavior.  In a concrete course one may take $X$ to be a family of divisors, modules, curves, bundles, or representations, depending on the chapter.  The moduli-theoretic lesson is that the morphism $T\to M$ detects both the general member and the limiting point.
\end{example}
\begin{exercise}
Construct a degeneration over $\Spec k[[t]]$ whose generic fiber satisfies the condition but whose special fiber forces a compactification.
\end{exercise}
\begin{exercise}
Formulate a small example in which the naive set of isomorphism classes is not Hausdorff, not separated, or not functorial.  Then repair the example by adding stability, rigidification, or a stack structure.
\end{exercise}
\clearpage
\section{Sheet 10: Representable functors: relative Hom}
This sheet studies \emph{relative Hom} as it appears in the broader chapter on \emph{representable functors}.  The recurring question is whether a proposed parameter space remembers enough information about families while remaining geometric enough to compute with.
\begin{definition}[Relative hom viewpoint]
A relative Hom moduli functor assigns to each test scheme $T$ the set, groupoid, or category of families over $T$ satisfying the chosen condition.  Two families are identified only when the moduli problem says that an isomorphism, not merely a bijection on fibers, has been supplied.
\end{definition}
\begin{proposition}[Relative hom viewpoint]
If a moduli functor for relative Hom is represented by a scheme $M$, then natural transformations from the functor to any functor $G$ are the same as elements of $G(M)$.  In practice this converts universal constructions on families into morphisms out of $M$.
\end{proposition}
\begin{example}[Relative hom viewpoint]
For relative Hom, take a base $T=\Spec R$ and replace a single object by an $R$-flat family.  Flatness prevents the numerical invariants from jumping on first-order thickenings, so it is the minimal condition under which the fibers can reasonably be compared.
\end{example}
\begin{theorem}[Relative hom viewpoint]
Suppose the objects in the relative Hom problem have no non-trivial automorphisms and admit effective descent for the chosen topology.  Then a universal family, if it exists, is unique up to unique isomorphism.  Consequently the representing object is unique up to unique isomorphism.
\end{theorem}
\begin{remark}[Relative hom viewpoint]
The phrase ``the moduli space of relative Hom'' is therefore shorthand.  It may mean a fine moduli scheme, a coarse moduli space, an algebraic stack, a GIT quotient, or a compactification.  The correct meaning is determined by what information is retained about families and automorphisms.
\end{remark}
\begin{proof}[Proof sketch of the theorem/proposition on this sheet]
The argument is the standard Yoneda-and-descent argument specialized to relative Hom.  A family over a test scheme $T$ corresponds, under representability, to a morphism $T\to M$.  Pulling the universal family back along this morphism recovers the given family.  If two morphisms give isomorphic pullbacks, the functorial equality on all test schemes forces the morphisms to agree.  The hypotheses on automorphisms ensure that the equality is not weakened to an equivalence relation in a groupoid.
\end{proof}
\begin{example}[A test calculation for relative Hom]
Let $T=\Spec k[s]$ and let the invariant of interest be denoted $\nu$.  A toy family $X\to T$ is chosen so that $\nu(X_t)=\nu_0$ for $s\ne0$ and the special fiber at $s=0$ records the boundary behavior.  In a concrete course one may take $X$ to be a family of divisors, modules, curves, bundles, or representations, depending on the chapter.  The moduli-theoretic lesson is that the morphism $T\to M$ detects both the general member and the limiting point.
\end{example}
\begin{exercise}
Compare the functor with the corresponding groupoid-valued functor.  Which information is lost after passing to isomorphism classes?
\end{exercise}
\begin{exercise}
Formulate a small example in which the naive set of isomorphism classes is not Hausdorff, not separated, or not functorial.  Then repair the example by adding stability, rigidification, or a stack structure.
\end{exercise}
\clearpage
\section{Sheet 11: Equivalence relations and quotients: group actions}
This sheet studies \emph{group actions} as it appears in the broader chapter on \emph{equivalence relations and quotients}.  The recurring question is whether a proposed parameter space remembers enough information about families while remaining geometric enough to compute with.
\begin{definition}[Group actions viewpoint]
A group actions moduli functor assigns to each test scheme $T$ the set, groupoid, or category of families over $T$ satisfying the chosen condition.  Two families are identified only when the moduli problem says that an isomorphism, not merely a bijection on fibers, has been supplied.
\end{definition}
\begin{proposition}[Group actions viewpoint]
If a moduli functor for group actions is represented by a scheme $M$, then natural transformations from the functor to any functor $G$ are the same as elements of $G(M)$.  In practice this converts universal constructions on families into morphisms out of $M$.
\end{proposition}
\begin{example}[Group actions viewpoint]
For group actions, take a base $T=\Spec R$ and replace a single object by an $R$-flat family.  Flatness prevents the numerical invariants from jumping on first-order thickenings, so it is the minimal condition under which the fibers can reasonably be compared.
\end{example}
\begin{theorem}[Group actions viewpoint]
Suppose the objects in the group actions problem have no non-trivial automorphisms and admit effective descent for the chosen topology.  Then a universal family, if it exists, is unique up to unique isomorphism.  Consequently the representing object is unique up to unique isomorphism.
\end{theorem}
\begin{remark}[Group actions viewpoint]
The phrase ``the moduli space of group actions'' is therefore shorthand.  It may mean a fine moduli scheme, a coarse moduli space, an algebraic stack, a GIT quotient, or a compactification.  The correct meaning is determined by what information is retained about families and automorphisms.
\end{remark}
\begin{proof}[Proof sketch of the theorem/proposition on this sheet]
The argument is the standard Yoneda-and-descent argument specialized to group actions.  A family over a test scheme $T$ corresponds, under representability, to a morphism $T\to M$.  Pulling the universal family back along this morphism recovers the given family.  If two morphisms give isomorphic pullbacks, the functorial equality on all test schemes forces the morphisms to agree.  The hypotheses on automorphisms ensure that the equality is not weakened to an equivalence relation in a groupoid.
\end{proof}
\begin{example}[A test calculation for group actions]
Let $T=\Spec k[s]$ and let the invariant of interest be denoted $\nu$.  A toy family $X\to T$ is chosen so that $\nu(X_t)=\nu_0$ for $s\ne0$ and the special fiber at $s=0$ records the boundary behavior.  In a concrete course one may take $X$ to be a family of divisors, modules, curves, bundles, or representations, depending on the chapter.  The moduli-theoretic lesson is that the morphism $T\to M$ detects both the general member and the limiting point.
\end{example}
\begin{exercise}
Write the functor of points for the moduli problem on this sheet and identify its value on $\Spec k$ and on $\Spec k[\epsilon]/(\epsilon^2)$.
\end{exercise}
\begin{exercise}
Formulate a small example in which the naive set of isomorphism classes is not Hausdorff, not separated, or not functorial.  Then repair the example by adding stability, rigidification, or a stack structure.
\end{exercise}
\clearpage
\section{Sheet 12: Equivalence relations and quotients: orbits}
This sheet studies \emph{orbits} as it appears in the broader chapter on \emph{equivalence relations and quotients}.  The recurring question is whether a proposed parameter space remembers enough information about families while remaining geometric enough to compute with.
\begin{definition}[Orbits viewpoint]
A orbits moduli functor assigns to each test scheme $T$ the set, groupoid, or category of families over $T$ satisfying the chosen condition.  Two families are identified only when the moduli problem says that an isomorphism, not merely a bijection on fibers, has been supplied.
\end{definition}
\begin{proposition}[Orbits viewpoint]
If a moduli functor for orbits is represented by a scheme $M$, then natural transformations from the functor to any functor $G$ are the same as elements of $G(M)$.  In practice this converts universal constructions on families into morphisms out of $M$.
\end{proposition}
\begin{example}[Orbits viewpoint]
For orbits, take a base $T=\Spec R$ and replace a single object by an $R$-flat family.  Flatness prevents the numerical invariants from jumping on first-order thickenings, so it is the minimal condition under which the fibers can reasonably be compared.
\end{example}
\begin{theorem}[Orbits viewpoint]
Suppose the objects in the orbits problem have no non-trivial automorphisms and admit effective descent for the chosen topology.  Then a universal family, if it exists, is unique up to unique isomorphism.  Consequently the representing object is unique up to unique isomorphism.
\end{theorem}
\begin{remark}[Orbits viewpoint]
The phrase ``the moduli space of orbits'' is therefore shorthand.  It may mean a fine moduli scheme, a coarse moduli space, an algebraic stack, a GIT quotient, or a compactification.  The correct meaning is determined by what information is retained about families and automorphisms.
\end{remark}
\begin{proof}[Proof sketch of the theorem/proposition on this sheet]
The argument is the standard Yoneda-and-descent argument specialized to orbits.  A family over a test scheme $T$ corresponds, under representability, to a morphism $T\to M$.  Pulling the universal family back along this morphism recovers the given family.  If two morphisms give isomorphic pullbacks, the functorial equality on all test schemes forces the morphisms to agree.  The hypotheses on automorphisms ensure that the equality is not weakened to an equivalence relation in a groupoid.
\end{proof}
\begin{example}[A test calculation for orbits]
Let $T=\Spec k[s]$ and let the invariant of interest be denoted $\nu$.  A toy family $X\to T$ is chosen so that $\nu(X_t)=\nu_0$ for $s\ne0$ and the special fiber at $s=0$ records the boundary behavior.  In a concrete course one may take $X$ to be a family of divisors, modules, curves, bundles, or representations, depending on the chapter.  The moduli-theoretic lesson is that the morphism $T\to M$ detects both the general member and the limiting point.
\end{example}
\begin{exercise}
Find one automorphism group occurring in the example and decide whether it obstructs the existence of a fine moduli scheme.
\end{exercise}
\begin{exercise}
Formulate a small example in which the naive set of isomorphism classes is not Hausdorff, not separated, or not functorial.  Then repair the example by adding stability, rigidification, or a stack structure.
\end{exercise}
\clearpage
\section{Sheet 13: Equivalence relations and quotients: stabilizers}
This sheet studies \emph{stabilizers} as it appears in the broader chapter on \emph{equivalence relations and quotients}.  The recurring question is whether a proposed parameter space remembers enough information about families while remaining geometric enough to compute with.
\begin{definition}[Stabilizers viewpoint]
A stabilizers moduli functor assigns to each test scheme $T$ the set, groupoid, or category of families over $T$ satisfying the chosen condition.  Two families are identified only when the moduli problem says that an isomorphism, not merely a bijection on fibers, has been supplied.
\end{definition}
\begin{proposition}[Stabilizers viewpoint]
If a moduli functor for stabilizers is represented by a scheme $M$, then natural transformations from the functor to any functor $G$ are the same as elements of $G(M)$.  In practice this converts universal constructions on families into morphisms out of $M$.
\end{proposition}
\begin{example}[Stabilizers viewpoint]
For stabilizers, take a base $T=\Spec R$ and replace a single object by an $R$-flat family.  Flatness prevents the numerical invariants from jumping on first-order thickenings, so it is the minimal condition under which the fibers can reasonably be compared.
\end{example}
\begin{theorem}[Stabilizers viewpoint]
Suppose the objects in the stabilizers problem have no non-trivial automorphisms and admit effective descent for the chosen topology.  Then a universal family, if it exists, is unique up to unique isomorphism.  Consequently the representing object is unique up to unique isomorphism.
\end{theorem}
\begin{remark}[Stabilizers viewpoint]
The phrase ``the moduli space of stabilizers'' is therefore shorthand.  It may mean a fine moduli scheme, a coarse moduli space, an algebraic stack, a GIT quotient, or a compactification.  The correct meaning is determined by what information is retained about families and automorphisms.
\end{remark}
\begin{proof}[Proof sketch of the theorem/proposition on this sheet]
The argument is the standard Yoneda-and-descent argument specialized to stabilizers.  A family over a test scheme $T$ corresponds, under representability, to a morphism $T\to M$.  Pulling the universal family back along this morphism recovers the given family.  If two morphisms give isomorphic pullbacks, the functorial equality on all test schemes forces the morphisms to agree.  The hypotheses on automorphisms ensure that the equality is not weakened to an equivalence relation in a groupoid.
\end{proof}
\begin{example}[A test calculation for stabilizers]
Let $T=\Spec k[s]$ and let the invariant of interest be denoted $\nu$.  A toy family $X\to T$ is chosen so that $\nu(X_t)=\nu_0$ for $s\ne0$ and the special fiber at $s=0$ records the boundary behavior.  In a concrete course one may take $X$ to be a family of divisors, modules, curves, bundles, or representations, depending on the chapter.  The moduli-theoretic lesson is that the morphism $T\to M$ detects both the general member and the limiting point.
\end{example}
\begin{exercise}
Check a base-change square explicitly: pull the displayed family back along $T\prime\to T$ and verify which invariant is preserved.
\end{exercise}
\begin{exercise}
Formulate a small example in which the naive set of isomorphism classes is not Hausdorff, not separated, or not functorial.  Then repair the example by adding stability, rigidification, or a stack structure.
\end{exercise}
\clearpage
\section{Sheet 14: Equivalence relations and quotients: categorical quotients}
This sheet studies \emph{categorical quotients} as it appears in the broader chapter on \emph{equivalence relations and quotients}.  The recurring question is whether a proposed parameter space remembers enough information about families while remaining geometric enough to compute with.
\begin{definition}[Categorical quotients viewpoint]
A categorical quotients moduli functor assigns to each test scheme $T$ the set, groupoid, or category of families over $T$ satisfying the chosen condition.  Two families are identified only when the moduli problem says that an isomorphism, not merely a bijection on fibers, has been supplied.
\end{definition}
\begin{proposition}[Categorical quotients viewpoint]
If a moduli functor for categorical quotients is represented by a scheme $M$, then natural transformations from the functor to any functor $G$ are the same as elements of $G(M)$.  In practice this converts universal constructions on families into morphisms out of $M$.
\end{proposition}
\begin{example}[Categorical quotients viewpoint]
For categorical quotients, take a base $T=\Spec R$ and replace a single object by an $R$-flat family.  Flatness prevents the numerical invariants from jumping on first-order thickenings, so it is the minimal condition under which the fibers can reasonably be compared.
\end{example}
\begin{theorem}[Categorical quotients viewpoint]
Suppose the objects in the categorical quotients problem have no non-trivial automorphisms and admit effective descent for the chosen topology.  Then a universal family, if it exists, is unique up to unique isomorphism.  Consequently the representing object is unique up to unique isomorphism.
\end{theorem}
\begin{remark}[Categorical quotients viewpoint]
The phrase ``the moduli space of categorical quotients'' is therefore shorthand.  It may mean a fine moduli scheme, a coarse moduli space, an algebraic stack, a GIT quotient, or a compactification.  The correct meaning is determined by what information is retained about families and automorphisms.
\end{remark}
\begin{proof}[Proof sketch of the theorem/proposition on this sheet]
The argument is the standard Yoneda-and-descent argument specialized to categorical quotients.  A family over a test scheme $T$ corresponds, under representability, to a morphism $T\to M$.  Pulling the universal family back along this morphism recovers the given family.  If two morphisms give isomorphic pullbacks, the functorial equality on all test schemes forces the morphisms to agree.  The hypotheses on automorphisms ensure that the equality is not weakened to an equivalence relation in a groupoid.
\end{proof}
\begin{example}[A test calculation for categorical quotients]
Let $T=\Spec k[s]$ and let the invariant of interest be denoted $\nu$.  A toy family $X\to T$ is chosen so that $\nu(X_t)=\nu_0$ for $s\ne0$ and the special fiber at $s=0$ records the boundary behavior.  In a concrete course one may take $X$ to be a family of divisors, modules, curves, bundles, or representations, depending on the chapter.  The moduli-theoretic lesson is that the morphism $T\to M$ detects both the general member and the limiting point.
\end{example}
\begin{exercise}
Construct a degeneration over $\Spec k[[t]]$ whose generic fiber satisfies the condition but whose special fiber forces a compactification.
\end{exercise}
\begin{exercise}
Formulate a small example in which the naive set of isomorphism classes is not Hausdorff, not separated, or not functorial.  Then repair the example by adding stability, rigidification, or a stack structure.
\end{exercise}
\clearpage
\section{Sheet 15: Equivalence relations and quotients: geometric quotients}
This sheet studies \emph{geometric quotients} as it appears in the broader chapter on \emph{equivalence relations and quotients}.  The recurring question is whether a proposed parameter space remembers enough information about families while remaining geometric enough to compute with.
\begin{definition}[Geometric quotients viewpoint]
A geometric quotients moduli functor assigns to each test scheme $T$ the set, groupoid, or category of families over $T$ satisfying the chosen condition.  Two families are identified only when the moduli problem says that an isomorphism, not merely a bijection on fibers, has been supplied.
\end{definition}
\begin{proposition}[Geometric quotients viewpoint]
If a moduli functor for geometric quotients is represented by a scheme $M$, then natural transformations from the functor to any functor $G$ are the same as elements of $G(M)$.  In practice this converts universal constructions on families into morphisms out of $M$.
\end{proposition}
\begin{example}[Geometric quotients viewpoint]
For geometric quotients, take a base $T=\Spec R$ and replace a single object by an $R$-flat family.  Flatness prevents the numerical invariants from jumping on first-order thickenings, so it is the minimal condition under which the fibers can reasonably be compared.
\end{example}
\begin{theorem}[Geometric quotients viewpoint]
Suppose the objects in the geometric quotients problem have no non-trivial automorphisms and admit effective descent for the chosen topology.  Then a universal family, if it exists, is unique up to unique isomorphism.  Consequently the representing object is unique up to unique isomorphism.
\end{theorem}
\begin{remark}[Geometric quotients viewpoint]
The phrase ``the moduli space of geometric quotients'' is therefore shorthand.  It may mean a fine moduli scheme, a coarse moduli space, an algebraic stack, a GIT quotient, or a compactification.  The correct meaning is determined by what information is retained about families and automorphisms.
\end{remark}
\begin{proof}[Proof sketch of the theorem/proposition on this sheet]
The argument is the standard Yoneda-and-descent argument specialized to geometric quotients.  A family over a test scheme $T$ corresponds, under representability, to a morphism $T\to M$.  Pulling the universal family back along this morphism recovers the given family.  If two morphisms give isomorphic pullbacks, the functorial equality on all test schemes forces the morphisms to agree.  The hypotheses on automorphisms ensure that the equality is not weakened to an equivalence relation in a groupoid.
\end{proof}
\begin{example}[A test calculation for geometric quotients]
Let $T=\Spec k[s]$ and let the invariant of interest be denoted $\nu$.  A toy family $X\to T$ is chosen so that $\nu(X_t)=\nu_0$ for $s\ne0$ and the special fiber at $s=0$ records the boundary behavior.  In a concrete course one may take $X$ to be a family of divisors, modules, curves, bundles, or representations, depending on the chapter.  The moduli-theoretic lesson is that the morphism $T\to M$ detects both the general member and the limiting point.
\end{example}
\begin{exercise}
Compare the functor with the corresponding groupoid-valued functor.  Which information is lost after passing to isomorphism classes?
\end{exercise}
\begin{exercise}
Formulate a small example in which the naive set of isomorphism classes is not Hausdorff, not separated, or not functorial.  Then repair the example by adding stability, rigidification, or a stack structure.
\end{exercise}
\clearpage
\section{Sheet 16: Deformation theory: dual numbers}
This sheet studies \emph{dual numbers} as it appears in the broader chapter on \emph{deformation theory}.  The recurring question is whether a proposed parameter space remembers enough information about families while remaining geometric enough to compute with.
\begin{definition}[Dual numbers viewpoint]
A dual numbers moduli functor assigns to each test scheme $T$ the set, groupoid, or category of families over $T$ satisfying the chosen condition.  Two families are identified only when the moduli problem says that an isomorphism, not merely a bijection on fibers, has been supplied.
\end{definition}
\begin{proposition}[Dual numbers viewpoint]
If a moduli functor for dual numbers is represented by a scheme $M$, then natural transformations from the functor to any functor $G$ are the same as elements of $G(M)$.  In practice this converts universal constructions on families into morphisms out of $M$.
\end{proposition}
\begin{example}[Dual numbers viewpoint]
For dual numbers, take a base $T=\Spec R$ and replace a single object by an $R$-flat family.  Flatness prevents the numerical invariants from jumping on first-order thickenings, so it is the minimal condition under which the fibers can reasonably be compared.
\end{example}
\begin{theorem}[Dual numbers viewpoint]
Suppose the objects in the dual numbers problem have no non-trivial automorphisms and admit effective descent for the chosen topology.  Then a universal family, if it exists, is unique up to unique isomorphism.  Consequently the representing object is unique up to unique isomorphism.
\end{theorem}
\begin{remark}[Dual numbers viewpoint]
The phrase ``the moduli space of dual numbers'' is therefore shorthand.  It may mean a fine moduli scheme, a coarse moduli space, an algebraic stack, a GIT quotient, or a compactification.  The correct meaning is determined by what information is retained about families and automorphisms.
\end{remark}
\begin{proof}[Proof sketch of the theorem/proposition on this sheet]
The argument is the standard Yoneda-and-descent argument specialized to dual numbers.  A family over a test scheme $T$ corresponds, under representability, to a morphism $T\to M$.  Pulling the universal family back along this morphism recovers the given family.  If two morphisms give isomorphic pullbacks, the functorial equality on all test schemes forces the morphisms to agree.  The hypotheses on automorphisms ensure that the equality is not weakened to an equivalence relation in a groupoid.
\end{proof}
\begin{example}[A test calculation for dual numbers]
Let $T=\Spec k[s]$ and let the invariant of interest be denoted $\nu$.  A toy family $X\to T$ is chosen so that $\nu(X_t)=\nu_0$ for $s\ne0$ and the special fiber at $s=0$ records the boundary behavior.  In a concrete course one may take $X$ to be a family of divisors, modules, curves, bundles, or representations, depending on the chapter.  The moduli-theoretic lesson is that the morphism $T\to M$ detects both the general member and the limiting point.
\end{example}
\begin{exercise}
Write the functor of points for the moduli problem on this sheet and identify its value on $\Spec k$ and on $\Spec k[\epsilon]/(\epsilon^2)$.
\end{exercise}
\begin{exercise}
Formulate a small example in which the naive set of isomorphism classes is not Hausdorff, not separated, or not functorial.  Then repair the example by adding stability, rigidification, or a stack structure.
\end{exercise}
\clearpage
\section{Sheet 17: Deformation theory: tangent spaces}
This sheet studies \emph{tangent spaces} as it appears in the broader chapter on \emph{deformation theory}.  The recurring question is whether a proposed parameter space remembers enough information about families while remaining geometric enough to compute with.
\begin{definition}[Tangent spaces viewpoint]
A tangent spaces moduli functor assigns to each test scheme $T$ the set, groupoid, or category of families over $T$ satisfying the chosen condition.  Two families are identified only when the moduli problem says that an isomorphism, not merely a bijection on fibers, has been supplied.
\end{definition}
\begin{proposition}[Tangent spaces viewpoint]
If a moduli functor for tangent spaces is represented by a scheme $M$, then natural transformations from the functor to any functor $G$ are the same as elements of $G(M)$.  In practice this converts universal constructions on families into morphisms out of $M$.
\end{proposition}
\begin{example}[Tangent spaces viewpoint]
For tangent spaces, take a base $T=\Spec R$ and replace a single object by an $R$-flat family.  Flatness prevents the numerical invariants from jumping on first-order thickenings, so it is the minimal condition under which the fibers can reasonably be compared.
\end{example}
\begin{theorem}[Tangent spaces viewpoint]
Suppose the objects in the tangent spaces problem have no non-trivial automorphisms and admit effective descent for the chosen topology.  Then a universal family, if it exists, is unique up to unique isomorphism.  Consequently the representing object is unique up to unique isomorphism.
\end{theorem}
\begin{remark}[Tangent spaces viewpoint]
The phrase ``the moduli space of tangent spaces'' is therefore shorthand.  It may mean a fine moduli scheme, a coarse moduli space, an algebraic stack, a GIT quotient, or a compactification.  The correct meaning is determined by what information is retained about families and automorphisms.
\end{remark}
\begin{proof}[Proof sketch of the theorem/proposition on this sheet]
The argument is the standard Yoneda-and-descent argument specialized to tangent spaces.  A family over a test scheme $T$ corresponds, under representability, to a morphism $T\to M$.  Pulling the universal family back along this morphism recovers the given family.  If two morphisms give isomorphic pullbacks, the functorial equality on all test schemes forces the morphisms to agree.  The hypotheses on automorphisms ensure that the equality is not weakened to an equivalence relation in a groupoid.
\end{proof}
\begin{example}[A test calculation for tangent spaces]
Let $T=\Spec k[s]$ and let the invariant of interest be denoted $\nu$.  A toy family $X\to T$ is chosen so that $\nu(X_t)=\nu_0$ for $s\ne0$ and the special fiber at $s=0$ records the boundary behavior.  In a concrete course one may take $X$ to be a family of divisors, modules, curves, bundles, or representations, depending on the chapter.  The moduli-theoretic lesson is that the morphism $T\to M$ detects both the general member and the limiting point.
\end{example}
\begin{exercise}
Find one automorphism group occurring in the example and decide whether it obstructs the existence of a fine moduli scheme.
\end{exercise}
\begin{exercise}
Formulate a small example in which the naive set of isomorphism classes is not Hausdorff, not separated, or not functorial.  Then repair the example by adding stability, rigidification, or a stack structure.
\end{exercise}
\clearpage
\section{Sheet 18: Deformation theory: obstructions}
This sheet studies \emph{obstructions} as it appears in the broader chapter on \emph{deformation theory}.  The recurring question is whether a proposed parameter space remembers enough information about families while remaining geometric enough to compute with.
\begin{definition}[Obstructions viewpoint]
A obstructions moduli functor assigns to each test scheme $T$ the set, groupoid, or category of families over $T$ satisfying the chosen condition.  Two families are identified only when the moduli problem says that an isomorphism, not merely a bijection on fibers, has been supplied.
\end{definition}
\begin{proposition}[Obstructions viewpoint]
If a moduli functor for obstructions is represented by a scheme $M$, then natural transformations from the functor to any functor $G$ are the same as elements of $G(M)$.  In practice this converts universal constructions on families into morphisms out of $M$.
\end{proposition}
\begin{example}[Obstructions viewpoint]
For obstructions, take a base $T=\Spec R$ and replace a single object by an $R$-flat family.  Flatness prevents the numerical invariants from jumping on first-order thickenings, so it is the minimal condition under which the fibers can reasonably be compared.
\end{example}
\begin{theorem}[Obstructions viewpoint]
Suppose the objects in the obstructions problem have no non-trivial automorphisms and admit effective descent for the chosen topology.  Then a universal family, if it exists, is unique up to unique isomorphism.  Consequently the representing object is unique up to unique isomorphism.
\end{theorem}
\begin{remark}[Obstructions viewpoint]
The phrase ``the moduli space of obstructions'' is therefore shorthand.  It may mean a fine moduli scheme, a coarse moduli space, an algebraic stack, a GIT quotient, or a compactification.  The correct meaning is determined by what information is retained about families and automorphisms.
\end{remark}
\begin{proof}[Proof sketch of the theorem/proposition on this sheet]
The argument is the standard Yoneda-and-descent argument specialized to obstructions.  A family over a test scheme $T$ corresponds, under representability, to a morphism $T\to M$.  Pulling the universal family back along this morphism recovers the given family.  If two morphisms give isomorphic pullbacks, the functorial equality on all test schemes forces the morphisms to agree.  The hypotheses on automorphisms ensure that the equality is not weakened to an equivalence relation in a groupoid.
\end{proof}
\begin{example}[A test calculation for obstructions]
Let $T=\Spec k[s]$ and let the invariant of interest be denoted $\nu$.  A toy family $X\to T$ is chosen so that $\nu(X_t)=\nu_0$ for $s\ne0$ and the special fiber at $s=0$ records the boundary behavior.  In a concrete course one may take $X$ to be a family of divisors, modules, curves, bundles, or representations, depending on the chapter.  The moduli-theoretic lesson is that the morphism $T\to M$ detects both the general member and the limiting point.
\end{example}
\begin{exercise}
Check a base-change square explicitly: pull the displayed family back along $T\prime\to T$ and verify which invariant is preserved.
\end{exercise}
\begin{exercise}
Formulate a small example in which the naive set of isomorphism classes is not Hausdorff, not separated, or not functorial.  Then repair the example by adding stability, rigidification, or a stack structure.
\end{exercise}
\clearpage
\section{Sheet 19: Deformation theory: Schlessinger conditions}
This sheet studies \emph{Schlessinger conditions} as it appears in the broader chapter on \emph{deformation theory}.  The recurring question is whether a proposed parameter space remembers enough information about families while remaining geometric enough to compute with.
\begin{definition}[Schlessinger conditions viewpoint]
A Schlessinger conditions moduli functor assigns to each test scheme $T$ the set, groupoid, or category of families over $T$ satisfying the chosen condition.  Two families are identified only when the moduli problem says that an isomorphism, not merely a bijection on fibers, has been supplied.
\end{definition}
\begin{proposition}[Schlessinger conditions viewpoint]
If a moduli functor for Schlessinger conditions is represented by a scheme $M$, then natural transformations from the functor to any functor $G$ are the same as elements of $G(M)$.  In practice this converts universal constructions on families into morphisms out of $M$.
\end{proposition}
\begin{example}[Schlessinger conditions viewpoint]
For Schlessinger conditions, take a base $T=\Spec R$ and replace a single object by an $R$-flat family.  Flatness prevents the numerical invariants from jumping on first-order thickenings, so it is the minimal condition under which the fibers can reasonably be compared.
\end{example}
\begin{theorem}[Schlessinger conditions viewpoint]
Suppose the objects in the Schlessinger conditions problem have no non-trivial automorphisms and admit effective descent for the chosen topology.  Then a universal family, if it exists, is unique up to unique isomorphism.  Consequently the representing object is unique up to unique isomorphism.
\end{theorem}
\begin{remark}[Schlessinger conditions viewpoint]
The phrase ``the moduli space of Schlessinger conditions'' is therefore shorthand.  It may mean a fine moduli scheme, a coarse moduli space, an algebraic stack, a GIT quotient, or a compactification.  The correct meaning is determined by what information is retained about families and automorphisms.
\end{remark}
\begin{proof}[Proof sketch of the theorem/proposition on this sheet]
The argument is the standard Yoneda-and-descent argument specialized to Schlessinger conditions.  A family over a test scheme $T$ corresponds, under representability, to a morphism $T\to M$.  Pulling the universal family back along this morphism recovers the given family.  If two morphisms give isomorphic pullbacks, the functorial equality on all test schemes forces the morphisms to agree.  The hypotheses on automorphisms ensure that the equality is not weakened to an equivalence relation in a groupoid.
\end{proof}
\begin{example}[A test calculation for Schlessinger conditions]
Let $T=\Spec k[s]$ and let the invariant of interest be denoted $\nu$.  A toy family $X\to T$ is chosen so that $\nu(X_t)=\nu_0$ for $s\ne0$ and the special fiber at $s=0$ records the boundary behavior.  In a concrete course one may take $X$ to be a family of divisors, modules, curves, bundles, or representations, depending on the chapter.  The moduli-theoretic lesson is that the morphism $T\to M$ detects both the general member and the limiting point.
\end{example}
\begin{exercise}
Construct a degeneration over $\Spec k[[t]]$ whose generic fiber satisfies the condition but whose special fiber forces a compactification.
\end{exercise}
\begin{exercise}
Formulate a small example in which the naive set of isomorphism classes is not Hausdorff, not separated, or not functorial.  Then repair the example by adding stability, rigidification, or a stack structure.
\end{exercise}
\clearpage
\section{Sheet 20: Deformation theory: formal moduli}
This sheet studies \emph{formal moduli} as it appears in the broader chapter on \emph{deformation theory}.  The recurring question is whether a proposed parameter space remembers enough information about families while remaining geometric enough to compute with.
\begin{definition}[Formal moduli viewpoint]
A formal moduli moduli functor assigns to each test scheme $T$ the set, groupoid, or category of families over $T$ satisfying the chosen condition.  Two families are identified only when the moduli problem says that an isomorphism, not merely a bijection on fibers, has been supplied.
\end{definition}
\begin{proposition}[Formal moduli viewpoint]
If a moduli functor for formal moduli is represented by a scheme $M$, then natural transformations from the functor to any functor $G$ are the same as elements of $G(M)$.  In practice this converts universal constructions on families into morphisms out of $M$.
\end{proposition}
\begin{example}[Formal moduli viewpoint]
For formal moduli, take a base $T=\Spec R$ and replace a single object by an $R$-flat family.  Flatness prevents the numerical invariants from jumping on first-order thickenings, so it is the minimal condition under which the fibers can reasonably be compared.
\end{example}
\begin{theorem}[Formal moduli viewpoint]
Suppose the objects in the formal moduli problem have no non-trivial automorphisms and admit effective descent for the chosen topology.  Then a universal family, if it exists, is unique up to unique isomorphism.  Consequently the representing object is unique up to unique isomorphism.
\end{theorem}
\begin{remark}[Formal moduli viewpoint]
The phrase ``the moduli space of formal moduli'' is therefore shorthand.  It may mean a fine moduli scheme, a coarse moduli space, an algebraic stack, a GIT quotient, or a compactification.  The correct meaning is determined by what information is retained about families and automorphisms.
\end{remark}
\begin{proof}[Proof sketch of the theorem/proposition on this sheet]
The argument is the standard Yoneda-and-descent argument specialized to formal moduli.  A family over a test scheme $T$ corresponds, under representability, to a morphism $T\to M$.  Pulling the universal family back along this morphism recovers the given family.  If two morphisms give isomorphic pullbacks, the functorial equality on all test schemes forces the morphisms to agree.  The hypotheses on automorphisms ensure that the equality is not weakened to an equivalence relation in a groupoid.
\end{proof}
\begin{example}[A test calculation for formal moduli]
Let $T=\Spec k[s]$ and let the invariant of interest be denoted $\nu$.  A toy family $X\to T$ is chosen so that $\nu(X_t)=\nu_0$ for $s\ne0$ and the special fiber at $s=0$ records the boundary behavior.  In a concrete course one may take $X$ to be a family of divisors, modules, curves, bundles, or representations, depending on the chapter.  The moduli-theoretic lesson is that the morphism $T\to M$ detects both the general member and the limiting point.
\end{example}
\begin{exercise}
Compare the functor with the corresponding groupoid-valued functor.  Which information is lost after passing to isomorphism classes?
\end{exercise}
\begin{exercise}
Formulate a small example in which the naive set of isomorphism classes is not Hausdorff, not separated, or not functorial.  Then repair the example by adding stability, rigidification, or a stack structure.
\end{exercise}
\clearpage
\section{Sheet 21: Moduli of vector bundles: rank and degree}
This sheet studies \emph{rank and degree} as it appears in the broader chapter on \emph{moduli of vector bundles}.  The recurring question is whether a proposed parameter space remembers enough information about families while remaining geometric enough to compute with.
\begin{definition}[Rank and degree viewpoint]
A rank and degree moduli functor assigns to each test scheme $T$ the set, groupoid, or category of families over $T$ satisfying the chosen condition.  Two families are identified only when the moduli problem says that an isomorphism, not merely a bijection on fibers, has been supplied.
\end{definition}
\begin{proposition}[Rank and degree viewpoint]
If a moduli functor for rank and degree is represented by a scheme $M$, then natural transformations from the functor to any functor $G$ are the same as elements of $G(M)$.  In practice this converts universal constructions on families into morphisms out of $M$.
\end{proposition}
\begin{example}[Rank and degree viewpoint]
For rank and degree, take a base $T=\Spec R$ and replace a single object by an $R$-flat family.  Flatness prevents the numerical invariants from jumping on first-order thickenings, so it is the minimal condition under which the fibers can reasonably be compared.
\end{example}
\begin{theorem}[Rank and degree viewpoint]
Suppose the objects in the rank and degree problem have no non-trivial automorphisms and admit effective descent for the chosen topology.  Then a universal family, if it exists, is unique up to unique isomorphism.  Consequently the representing object is unique up to unique isomorphism.
\end{theorem}
\begin{remark}[Rank and degree viewpoint]
The phrase ``the moduli space of rank and degree'' is therefore shorthand.  It may mean a fine moduli scheme, a coarse moduli space, an algebraic stack, a GIT quotient, or a compactification.  The correct meaning is determined by what information is retained about families and automorphisms.
\end{remark}
\begin{proof}[Proof sketch of the theorem/proposition on this sheet]
The argument is the standard Yoneda-and-descent argument specialized to rank and degree.  A family over a test scheme $T$ corresponds, under representability, to a morphism $T\to M$.  Pulling the universal family back along this morphism recovers the given family.  If two morphisms give isomorphic pullbacks, the functorial equality on all test schemes forces the morphisms to agree.  The hypotheses on automorphisms ensure that the equality is not weakened to an equivalence relation in a groupoid.
\end{proof}
\begin{example}[A test calculation for rank and degree]
Let $T=\Spec k[s]$ and let the invariant of interest be denoted $\nu$.  A toy family $X\to T$ is chosen so that $\nu(X_t)=\nu_0$ for $s\ne0$ and the special fiber at $s=0$ records the boundary behavior.  In a concrete course one may take $X$ to be a family of divisors, modules, curves, bundles, or representations, depending on the chapter.  The moduli-theoretic lesson is that the morphism $T\to M$ detects both the general member and the limiting point.
\end{example}
\begin{exercise}
Write the functor of points for the moduli problem on this sheet and identify its value on $\Spec k$ and on $\Spec k[\epsilon]/(\epsilon^2)$.
\end{exercise}
\begin{exercise}
Formulate a small example in which the naive set of isomorphism classes is not Hausdorff, not separated, or not functorial.  Then repair the example by adding stability, rigidification, or a stack structure.
\end{exercise}
\clearpage
\section{Sheet 22: Moduli of vector bundles: families of bundles}
This sheet studies \emph{families of bundles} as it appears in the broader chapter on \emph{moduli of vector bundles}.  The recurring question is whether a proposed parameter space remembers enough information about families while remaining geometric enough to compute with.
\begin{definition}[Families of bundles viewpoint]
A families of bundles moduli functor assigns to each test scheme $T$ the set, groupoid, or category of families over $T$ satisfying the chosen condition.  Two families are identified only when the moduli problem says that an isomorphism, not merely a bijection on fibers, has been supplied.
\end{definition}
\begin{proposition}[Families of bundles viewpoint]
If a moduli functor for families of bundles is represented by a scheme $M$, then natural transformations from the functor to any functor $G$ are the same as elements of $G(M)$.  In practice this converts universal constructions on families into morphisms out of $M$.
\end{proposition}
\begin{example}[Families of bundles viewpoint]
For families of bundles, take a base $T=\Spec R$ and replace a single object by an $R$-flat family.  Flatness prevents the numerical invariants from jumping on first-order thickenings, so it is the minimal condition under which the fibers can reasonably be compared.
\end{example}
\begin{theorem}[Families of bundles viewpoint]
Suppose the objects in the families of bundles problem have no non-trivial automorphisms and admit effective descent for the chosen topology.  Then a universal family, if it exists, is unique up to unique isomorphism.  Consequently the representing object is unique up to unique isomorphism.
\end{theorem}
\begin{remark}[Families of bundles viewpoint]
The phrase ``the moduli space of families of bundles'' is therefore shorthand.  It may mean a fine moduli scheme, a coarse moduli space, an algebraic stack, a GIT quotient, or a compactification.  The correct meaning is determined by what information is retained about families and automorphisms.
\end{remark}
\begin{proof}[Proof sketch of the theorem/proposition on this sheet]
The argument is the standard Yoneda-and-descent argument specialized to families of bundles.  A family over a test scheme $T$ corresponds, under representability, to a morphism $T\to M$.  Pulling the universal family back along this morphism recovers the given family.  If two morphisms give isomorphic pullbacks, the functorial equality on all test schemes forces the morphisms to agree.  The hypotheses on automorphisms ensure that the equality is not weakened to an equivalence relation in a groupoid.
\end{proof}
\begin{example}[A test calculation for families of bundles]
Let $T=\Spec k[s]$ and let the invariant of interest be denoted $\nu$.  A toy family $X\to T$ is chosen so that $\nu(X_t)=\nu_0$ for $s\ne0$ and the special fiber at $s=0$ records the boundary behavior.  In a concrete course one may take $X$ to be a family of divisors, modules, curves, bundles, or representations, depending on the chapter.  The moduli-theoretic lesson is that the morphism $T\to M$ detects both the general member and the limiting point.
\end{example}
\begin{exercise}
Find one automorphism group occurring in the example and decide whether it obstructs the existence of a fine moduli scheme.
\end{exercise}
\begin{exercise}
Formulate a small example in which the naive set of isomorphism classes is not Hausdorff, not separated, or not functorial.  Then repair the example by adding stability, rigidification, or a stack structure.
\end{exercise}
\clearpage
\section{Sheet 23: Moduli of vector bundles: slope stability}
This sheet studies \emph{slope stability} as it appears in the broader chapter on \emph{moduli of vector bundles}.  The recurring question is whether a proposed parameter space remembers enough information about families while remaining geometric enough to compute with.
\begin{definition}[Slope stability viewpoint]
A slope stability moduli functor assigns to each test scheme $T$ the set, groupoid, or category of families over $T$ satisfying the chosen condition.  Two families are identified only when the moduli problem says that an isomorphism, not merely a bijection on fibers, has been supplied.
\end{definition}
\begin{proposition}[Slope stability viewpoint]
If a moduli functor for slope stability is represented by a scheme $M$, then natural transformations from the functor to any functor $G$ are the same as elements of $G(M)$.  In practice this converts universal constructions on families into morphisms out of $M$.
\end{proposition}
\begin{example}[Slope stability viewpoint]
For slope stability, take a base $T=\Spec R$ and replace a single object by an $R$-flat family.  Flatness prevents the numerical invariants from jumping on first-order thickenings, so it is the minimal condition under which the fibers can reasonably be compared.
\end{example}
\begin{theorem}[Slope stability viewpoint]
Suppose the objects in the slope stability problem have no non-trivial automorphisms and admit effective descent for the chosen topology.  Then a universal family, if it exists, is unique up to unique isomorphism.  Consequently the representing object is unique up to unique isomorphism.
\end{theorem}
\begin{remark}[Slope stability viewpoint]
The phrase ``the moduli space of slope stability'' is therefore shorthand.  It may mean a fine moduli scheme, a coarse moduli space, an algebraic stack, a GIT quotient, or a compactification.  The correct meaning is determined by what information is retained about families and automorphisms.
\end{remark}
\begin{proof}[Proof sketch of the theorem/proposition on this sheet]
The argument is the standard Yoneda-and-descent argument specialized to slope stability.  A family over a test scheme $T$ corresponds, under representability, to a morphism $T\to M$.  Pulling the universal family back along this morphism recovers the given family.  If two morphisms give isomorphic pullbacks, the functorial equality on all test schemes forces the morphisms to agree.  The hypotheses on automorphisms ensure that the equality is not weakened to an equivalence relation in a groupoid.
\end{proof}
\begin{example}[A test calculation for slope stability]
Let $T=\Spec k[s]$ and let the invariant of interest be denoted $\nu$.  A toy family $X\to T$ is chosen so that $\nu(X_t)=\nu_0$ for $s\ne0$ and the special fiber at $s=0$ records the boundary behavior.  In a concrete course one may take $X$ to be a family of divisors, modules, curves, bundles, or representations, depending on the chapter.  The moduli-theoretic lesson is that the morphism $T\to M$ detects both the general member and the limiting point.
\end{example}
\begin{exercise}
Check a base-change square explicitly: pull the displayed family back along $T\prime\to T$ and verify which invariant is preserved.
\end{exercise}
\begin{exercise}
Formulate a small example in which the naive set of isomorphism classes is not Hausdorff, not separated, or not functorial.  Then repair the example by adding stability, rigidification, or a stack structure.
\end{exercise}
\clearpage
\section{Sheet 24: Moduli of vector bundles: Jordan--Holder filtrations}
This sheet studies \emph{Jordan--Holder filtrations} as it appears in the broader chapter on \emph{moduli of vector bundles}.  The recurring question is whether a proposed parameter space remembers enough information about families while remaining geometric enough to compute with.
\begin{definition}[Jordan--holder filtrations viewpoint]
A Jordan--Holder filtrations moduli functor assigns to each test scheme $T$ the set, groupoid, or category of families over $T$ satisfying the chosen condition.  Two families are identified only when the moduli problem says that an isomorphism, not merely a bijection on fibers, has been supplied.
\end{definition}
\begin{proposition}[Jordan--holder filtrations viewpoint]
If a moduli functor for Jordan--Holder filtrations is represented by a scheme $M$, then natural transformations from the functor to any functor $G$ are the same as elements of $G(M)$.  In practice this converts universal constructions on families into morphisms out of $M$.
\end{proposition}
\begin{example}[Jordan--holder filtrations viewpoint]
For Jordan--Holder filtrations, take a base $T=\Spec R$ and replace a single object by an $R$-flat family.  Flatness prevents the numerical invariants from jumping on first-order thickenings, so it is the minimal condition under which the fibers can reasonably be compared.
\end{example}
\begin{theorem}[Jordan--holder filtrations viewpoint]
Suppose the objects in the Jordan--Holder filtrations problem have no non-trivial automorphisms and admit effective descent for the chosen topology.  Then a universal family, if it exists, is unique up to unique isomorphism.  Consequently the representing object is unique up to unique isomorphism.
\end{theorem}
\begin{remark}[Jordan--holder filtrations viewpoint]
The phrase ``the moduli space of Jordan--Holder filtrations'' is therefore shorthand.  It may mean a fine moduli scheme, a coarse moduli space, an algebraic stack, a GIT quotient, or a compactification.  The correct meaning is determined by what information is retained about families and automorphisms.
\end{remark}
\begin{proof}[Proof sketch of the theorem/proposition on this sheet]
The argument is the standard Yoneda-and-descent argument specialized to Jordan--Holder filtrations.  A family over a test scheme $T$ corresponds, under representability, to a morphism $T\to M$.  Pulling the universal family back along this morphism recovers the given family.  If two morphisms give isomorphic pullbacks, the functorial equality on all test schemes forces the morphisms to agree.  The hypotheses on automorphisms ensure that the equality is not weakened to an equivalence relation in a groupoid.
\end{proof}
\begin{example}[A test calculation for Jordan--Holder filtrations]
Let $T=\Spec k[s]$ and let the invariant of interest be denoted $\nu$.  A toy family $X\to T$ is chosen so that $\nu(X_t)=\nu_0$ for $s\ne0$ and the special fiber at $s=0$ records the boundary behavior.  In a concrete course one may take $X$ to be a family of divisors, modules, curves, bundles, or representations, depending on the chapter.  The moduli-theoretic lesson is that the morphism $T\to M$ detects both the general member and the limiting point.
\end{example}
\begin{exercise}
Construct a degeneration over $\Spec k[[t]]$ whose generic fiber satisfies the condition but whose special fiber forces a compactification.
\end{exercise}
\begin{exercise}
Formulate a small example in which the naive set of isomorphism classes is not Hausdorff, not separated, or not functorial.  Then repair the example by adding stability, rigidification, or a stack structure.
\end{exercise}
\clearpage
\section{Sheet 25: Moduli of vector bundles: S-equivalence}
This sheet studies \emph{S-equivalence} as it appears in the broader chapter on \emph{moduli of vector bundles}.  The recurring question is whether a proposed parameter space remembers enough information about families while remaining geometric enough to compute with.
\begin{definition}[S-equivalence viewpoint]
A S-equivalence moduli functor assigns to each test scheme $T$ the set, groupoid, or category of families over $T$ satisfying the chosen condition.  Two families are identified only when the moduli problem says that an isomorphism, not merely a bijection on fibers, has been supplied.
\end{definition}
\begin{proposition}[S-equivalence viewpoint]
If a moduli functor for S-equivalence is represented by a scheme $M$, then natural transformations from the functor to any functor $G$ are the same as elements of $G(M)$.  In practice this converts universal constructions on families into morphisms out of $M$.
\end{proposition}
\begin{example}[S-equivalence viewpoint]
For S-equivalence, take a base $T=\Spec R$ and replace a single object by an $R$-flat family.  Flatness prevents the numerical invariants from jumping on first-order thickenings, so it is the minimal condition under which the fibers can reasonably be compared.
\end{example}
\begin{theorem}[S-equivalence viewpoint]
Suppose the objects in the S-equivalence problem have no non-trivial automorphisms and admit effective descent for the chosen topology.  Then a universal family, if it exists, is unique up to unique isomorphism.  Consequently the representing object is unique up to unique isomorphism.
\end{theorem}
\begin{remark}[S-equivalence viewpoint]
The phrase ``the moduli space of S-equivalence'' is therefore shorthand.  It may mean a fine moduli scheme, a coarse moduli space, an algebraic stack, a GIT quotient, or a compactification.  The correct meaning is determined by what information is retained about families and automorphisms.
\end{remark}
\begin{proof}[Proof sketch of the theorem/proposition on this sheet]
The argument is the standard Yoneda-and-descent argument specialized to S-equivalence.  A family over a test scheme $T$ corresponds, under representability, to a morphism $T\to M$.  Pulling the universal family back along this morphism recovers the given family.  If two morphisms give isomorphic pullbacks, the functorial equality on all test schemes forces the morphisms to agree.  The hypotheses on automorphisms ensure that the equality is not weakened to an equivalence relation in a groupoid.
\end{proof}
\begin{example}[A test calculation for S-equivalence]
Let $T=\Spec k[s]$ and let the invariant of interest be denoted $\nu$.  A toy family $X\to T$ is chosen so that $\nu(X_t)=\nu_0$ for $s\ne0$ and the special fiber at $s=0$ records the boundary behavior.  In a concrete course one may take $X$ to be a family of divisors, modules, curves, bundles, or representations, depending on the chapter.  The moduli-theoretic lesson is that the morphism $T\to M$ detects both the general member and the limiting point.
\end{example}
\begin{exercise}
Compare the functor with the corresponding groupoid-valued functor.  Which information is lost after passing to isomorphism classes?
\end{exercise}
\begin{exercise}
Formulate a small example in which the naive set of isomorphism classes is not Hausdorff, not separated, or not functorial.  Then repair the example by adding stability, rigidification, or a stack structure.
\end{exercise}
\clearpage
\section{Sheet 26: Curves and stable curves: prestable curves}
This sheet studies \emph{prestable curves} as it appears in the broader chapter on \emph{curves and stable curves}.  The recurring question is whether a proposed parameter space remembers enough information about families while remaining geometric enough to compute with.
\begin{definition}[Prestable curves viewpoint]
A prestable curves moduli functor assigns to each test scheme $T$ the set, groupoid, or category of families over $T$ satisfying the chosen condition.  Two families are identified only when the moduli problem says that an isomorphism, not merely a bijection on fibers, has been supplied.
\end{definition}
\begin{proposition}[Prestable curves viewpoint]
If a moduli functor for prestable curves is represented by a scheme $M$, then natural transformations from the functor to any functor $G$ are the same as elements of $G(M)$.  In practice this converts universal constructions on families into morphisms out of $M$.
\end{proposition}
\begin{example}[Prestable curves viewpoint]
For prestable curves, take a base $T=\Spec R$ and replace a single object by an $R$-flat family.  Flatness prevents the numerical invariants from jumping on first-order thickenings, so it is the minimal condition under which the fibers can reasonably be compared.
\end{example}
\begin{theorem}[Prestable curves viewpoint]
Suppose the objects in the prestable curves problem have no non-trivial automorphisms and admit effective descent for the chosen topology.  Then a universal family, if it exists, is unique up to unique isomorphism.  Consequently the representing object is unique up to unique isomorphism.
\end{theorem}
\begin{remark}[Prestable curves viewpoint]
The phrase ``the moduli space of prestable curves'' is therefore shorthand.  It may mean a fine moduli scheme, a coarse moduli space, an algebraic stack, a GIT quotient, or a compactification.  The correct meaning is determined by what information is retained about families and automorphisms.
\end{remark}
\begin{proof}[Proof sketch of the theorem/proposition on this sheet]
The argument is the standard Yoneda-and-descent argument specialized to prestable curves.  A family over a test scheme $T$ corresponds, under representability, to a morphism $T\to M$.  Pulling the universal family back along this morphism recovers the given family.  If two morphisms give isomorphic pullbacks, the functorial equality on all test schemes forces the morphisms to agree.  The hypotheses on automorphisms ensure that the equality is not weakened to an equivalence relation in a groupoid.
\end{proof}
\begin{example}[A test calculation for prestable curves]
Let $T=\Spec k[s]$ and let the invariant of interest be denoted $\nu$.  A toy family $X\to T$ is chosen so that $\nu(X_t)=\nu_0$ for $s\ne0$ and the special fiber at $s=0$ records the boundary behavior.  In a concrete course one may take $X$ to be a family of divisors, modules, curves, bundles, or representations, depending on the chapter.  The moduli-theoretic lesson is that the morphism $T\to M$ detects both the general member and the limiting point.
\end{example}
\begin{exercise}
Write the functor of points for the moduli problem on this sheet and identify its value on $\Spec k$ and on $\Spec k[\epsilon]/(\epsilon^2)$.
\end{exercise}
\begin{exercise}
Formulate a small example in which the naive set of isomorphism classes is not Hausdorff, not separated, or not functorial.  Then repair the example by adding stability, rigidification, or a stack structure.
\end{exercise}
\clearpage
\section{Sheet 27: Curves and stable curves: dual graphs}
This sheet studies \emph{dual graphs} as it appears in the broader chapter on \emph{curves and stable curves}.  The recurring question is whether a proposed parameter space remembers enough information about families while remaining geometric enough to compute with.
\begin{definition}[Dual graphs viewpoint]
A dual graphs moduli functor assigns to each test scheme $T$ the set, groupoid, or category of families over $T$ satisfying the chosen condition.  Two families are identified only when the moduli problem says that an isomorphism, not merely a bijection on fibers, has been supplied.
\end{definition}
\begin{proposition}[Dual graphs viewpoint]
If a moduli functor for dual graphs is represented by a scheme $M$, then natural transformations from the functor to any functor $G$ are the same as elements of $G(M)$.  In practice this converts universal constructions on families into morphisms out of $M$.
\end{proposition}
\begin{example}[Dual graphs viewpoint]
For dual graphs, take a base $T=\Spec R$ and replace a single object by an $R$-flat family.  Flatness prevents the numerical invariants from jumping on first-order thickenings, so it is the minimal condition under which the fibers can reasonably be compared.
\end{example}
\begin{theorem}[Dual graphs viewpoint]
Suppose the objects in the dual graphs problem have no non-trivial automorphisms and admit effective descent for the chosen topology.  Then a universal family, if it exists, is unique up to unique isomorphism.  Consequently the representing object is unique up to unique isomorphism.
\end{theorem}
\begin{remark}[Dual graphs viewpoint]
The phrase ``the moduli space of dual graphs'' is therefore shorthand.  It may mean a fine moduli scheme, a coarse moduli space, an algebraic stack, a GIT quotient, or a compactification.  The correct meaning is determined by what information is retained about families and automorphisms.
\end{remark}
\begin{proof}[Proof sketch of the theorem/proposition on this sheet]
The argument is the standard Yoneda-and-descent argument specialized to dual graphs.  A family over a test scheme $T$ corresponds, under representability, to a morphism $T\to M$.  Pulling the universal family back along this morphism recovers the given family.  If two morphisms give isomorphic pullbacks, the functorial equality on all test schemes forces the morphisms to agree.  The hypotheses on automorphisms ensure that the equality is not weakened to an equivalence relation in a groupoid.
\end{proof}
\begin{example}[A test calculation for dual graphs]
Let $T=\Spec k[s]$ and let the invariant of interest be denoted $\nu$.  A toy family $X\to T$ is chosen so that $\nu(X_t)=\nu_0$ for $s\ne0$ and the special fiber at $s=0$ records the boundary behavior.  In a concrete course one may take $X$ to be a family of divisors, modules, curves, bundles, or representations, depending on the chapter.  The moduli-theoretic lesson is that the morphism $T\to M$ detects both the general member and the limiting point.
\end{example}
\begin{exercise}
Find one automorphism group occurring in the example and decide whether it obstructs the existence of a fine moduli scheme.
\end{exercise}
\begin{exercise}
Formulate a small example in which the naive set of isomorphism classes is not Hausdorff, not separated, or not functorial.  Then repair the example by adding stability, rigidification, or a stack structure.
\end{exercise}
\clearpage
\section{Sheet 28: Curves and stable curves: stable pointed curves}
This sheet studies \emph{stable pointed curves} as it appears in the broader chapter on \emph{curves and stable curves}.  The recurring question is whether a proposed parameter space remembers enough information about families while remaining geometric enough to compute with.
\begin{definition}[Stable pointed curves viewpoint]
A stable pointed curves moduli functor assigns to each test scheme $T$ the set, groupoid, or category of families over $T$ satisfying the chosen condition.  Two families are identified only when the moduli problem says that an isomorphism, not merely a bijection on fibers, has been supplied.
\end{definition}
\begin{proposition}[Stable pointed curves viewpoint]
If a moduli functor for stable pointed curves is represented by a scheme $M$, then natural transformations from the functor to any functor $G$ are the same as elements of $G(M)$.  In practice this converts universal constructions on families into morphisms out of $M$.
\end{proposition}
\begin{example}[Stable pointed curves viewpoint]
For stable pointed curves, take a base $T=\Spec R$ and replace a single object by an $R$-flat family.  Flatness prevents the numerical invariants from jumping on first-order thickenings, so it is the minimal condition under which the fibers can reasonably be compared.
\end{example}
\begin{theorem}[Stable pointed curves viewpoint]
Suppose the objects in the stable pointed curves problem have no non-trivial automorphisms and admit effective descent for the chosen topology.  Then a universal family, if it exists, is unique up to unique isomorphism.  Consequently the representing object is unique up to unique isomorphism.
\end{theorem}
\begin{remark}[Stable pointed curves viewpoint]
The phrase ``the moduli space of stable pointed curves'' is therefore shorthand.  It may mean a fine moduli scheme, a coarse moduli space, an algebraic stack, a GIT quotient, or a compactification.  The correct meaning is determined by what information is retained about families and automorphisms.
\end{remark}
\begin{proof}[Proof sketch of the theorem/proposition on this sheet]
The argument is the standard Yoneda-and-descent argument specialized to stable pointed curves.  A family over a test scheme $T$ corresponds, under representability, to a morphism $T\to M$.  Pulling the universal family back along this morphism recovers the given family.  If two morphisms give isomorphic pullbacks, the functorial equality on all test schemes forces the morphisms to agree.  The hypotheses on automorphisms ensure that the equality is not weakened to an equivalence relation in a groupoid.
\end{proof}
\begin{example}[A test calculation for stable pointed curves]
Let $T=\Spec k[s]$ and let the invariant of interest be denoted $\nu$.  A toy family $X\to T$ is chosen so that $\nu(X_t)=\nu_0$ for $s\ne0$ and the special fiber at $s=0$ records the boundary behavior.  In a concrete course one may take $X$ to be a family of divisors, modules, curves, bundles, or representations, depending on the chapter.  The moduli-theoretic lesson is that the morphism $T\to M$ detects both the general member and the limiting point.
\end{example}
\begin{exercise}
Check a base-change square explicitly: pull the displayed family back along $T\prime\to T$ and verify which invariant is preserved.
\end{exercise}
\begin{exercise}
Formulate a small example in which the naive set of isomorphism classes is not Hausdorff, not separated, or not functorial.  Then repair the example by adding stability, rigidification, or a stack structure.
\end{exercise}
\clearpage
\section{Sheet 29: Curves and stable curves: clutching maps}
This sheet studies \emph{clutching maps} as it appears in the broader chapter on \emph{curves and stable curves}.  The recurring question is whether a proposed parameter space remembers enough information about families while remaining geometric enough to compute with.
\begin{definition}[Clutching maps viewpoint]
A clutching maps moduli functor assigns to each test scheme $T$ the set, groupoid, or category of families over $T$ satisfying the chosen condition.  Two families are identified only when the moduli problem says that an isomorphism, not merely a bijection on fibers, has been supplied.
\end{definition}
\begin{proposition}[Clutching maps viewpoint]
If a moduli functor for clutching maps is represented by a scheme $M$, then natural transformations from the functor to any functor $G$ are the same as elements of $G(M)$.  In practice this converts universal constructions on families into morphisms out of $M$.
\end{proposition}
\begin{example}[Clutching maps viewpoint]
For clutching maps, take a base $T=\Spec R$ and replace a single object by an $R$-flat family.  Flatness prevents the numerical invariants from jumping on first-order thickenings, so it is the minimal condition under which the fibers can reasonably be compared.
\end{example}
\begin{theorem}[Clutching maps viewpoint]
Suppose the objects in the clutching maps problem have no non-trivial automorphisms and admit effective descent for the chosen topology.  Then a universal family, if it exists, is unique up to unique isomorphism.  Consequently the representing object is unique up to unique isomorphism.
\end{theorem}
\begin{remark}[Clutching maps viewpoint]
The phrase ``the moduli space of clutching maps'' is therefore shorthand.  It may mean a fine moduli scheme, a coarse moduli space, an algebraic stack, a GIT quotient, or a compactification.  The correct meaning is determined by what information is retained about families and automorphisms.
\end{remark}
\begin{proof}[Proof sketch of the theorem/proposition on this sheet]
The argument is the standard Yoneda-and-descent argument specialized to clutching maps.  A family over a test scheme $T$ corresponds, under representability, to a morphism $T\to M$.  Pulling the universal family back along this morphism recovers the given family.  If two morphisms give isomorphic pullbacks, the functorial equality on all test schemes forces the morphisms to agree.  The hypotheses on automorphisms ensure that the equality is not weakened to an equivalence relation in a groupoid.
\end{proof}
\begin{example}[A test calculation for clutching maps]
Let $T=\Spec k[s]$ and let the invariant of interest be denoted $\nu$.  A toy family $X\to T$ is chosen so that $\nu(X_t)=\nu_0$ for $s\ne0$ and the special fiber at $s=0$ records the boundary behavior.  In a concrete course one may take $X$ to be a family of divisors, modules, curves, bundles, or representations, depending on the chapter.  The moduli-theoretic lesson is that the morphism $T\to M$ detects both the general member and the limiting point.
\end{example}
\begin{exercise}
Construct a degeneration over $\Spec k[[t]]$ whose generic fiber satisfies the condition but whose special fiber forces a compactification.
\end{exercise}
\begin{exercise}
Formulate a small example in which the naive set of isomorphism classes is not Hausdorff, not separated, or not functorial.  Then repair the example by adding stability, rigidification, or a stack structure.
\end{exercise}
\clearpage
\section{Sheet 30: Curves and stable curves: boundary strata}
This sheet studies \emph{boundary strata} as it appears in the broader chapter on \emph{curves and stable curves}.  The recurring question is whether a proposed parameter space remembers enough information about families while remaining geometric enough to compute with.
\begin{definition}[Boundary strata viewpoint]
A boundary strata moduli functor assigns to each test scheme $T$ the set, groupoid, or category of families over $T$ satisfying the chosen condition.  Two families are identified only when the moduli problem says that an isomorphism, not merely a bijection on fibers, has been supplied.
\end{definition}
\begin{proposition}[Boundary strata viewpoint]
If a moduli functor for boundary strata is represented by a scheme $M$, then natural transformations from the functor to any functor $G$ are the same as elements of $G(M)$.  In practice this converts universal constructions on families into morphisms out of $M$.
\end{proposition}
\begin{example}[Boundary strata viewpoint]
For boundary strata, take a base $T=\Spec R$ and replace a single object by an $R$-flat family.  Flatness prevents the numerical invariants from jumping on first-order thickenings, so it is the minimal condition under which the fibers can reasonably be compared.
\end{example}
\begin{theorem}[Boundary strata viewpoint]
Suppose the objects in the boundary strata problem have no non-trivial automorphisms and admit effective descent for the chosen topology.  Then a universal family, if it exists, is unique up to unique isomorphism.  Consequently the representing object is unique up to unique isomorphism.
\end{theorem}
\begin{remark}[Boundary strata viewpoint]
The phrase ``the moduli space of boundary strata'' is therefore shorthand.  It may mean a fine moduli scheme, a coarse moduli space, an algebraic stack, a GIT quotient, or a compactification.  The correct meaning is determined by what information is retained about families and automorphisms.
\end{remark}
\begin{proof}[Proof sketch of the theorem/proposition on this sheet]
The argument is the standard Yoneda-and-descent argument specialized to boundary strata.  A family over a test scheme $T$ corresponds, under representability, to a morphism $T\to M$.  Pulling the universal family back along this morphism recovers the given family.  If two morphisms give isomorphic pullbacks, the functorial equality on all test schemes forces the morphisms to agree.  The hypotheses on automorphisms ensure that the equality is not weakened to an equivalence relation in a groupoid.
\end{proof}
\begin{example}[A test calculation for boundary strata]
Let $T=\Spec k[s]$ and let the invariant of interest be denoted $\nu$.  A toy family $X\to T$ is chosen so that $\nu(X_t)=\nu_0$ for $s\ne0$ and the special fiber at $s=0$ records the boundary behavior.  In a concrete course one may take $X$ to be a family of divisors, modules, curves, bundles, or representations, depending on the chapter.  The moduli-theoretic lesson is that the morphism $T\to M$ detects both the general member and the limiting point.
\end{example}
\begin{exercise}
Compare the functor with the corresponding groupoid-valued functor.  Which information is lost after passing to isomorphism classes?
\end{exercise}
\begin{exercise}
Formulate a small example in which the naive set of isomorphism classes is not Hausdorff, not separated, or not functorial.  Then repair the example by adding stability, rigidification, or a stack structure.
\end{exercise}
\clearpage
\section{Sheet 31: Line bundles and Picard functors: relative Picard functor}
This sheet studies \emph{relative Picard functor} as it appears in the broader chapter on \emph{line bundles and picard functors}.  The recurring question is whether a proposed parameter space remembers enough information about families while remaining geometric enough to compute with.
\begin{definition}[Relative picard functor viewpoint]
A relative Picard functor moduli functor assigns to each test scheme $T$ the set, groupoid, or category of families over $T$ satisfying the chosen condition.  Two families are identified only when the moduli problem says that an isomorphism, not merely a bijection on fibers, has been supplied.
\end{definition}
\begin{proposition}[Relative picard functor viewpoint]
If a moduli functor for relative Picard functor is represented by a scheme $M$, then natural transformations from the functor to any functor $G$ are the same as elements of $G(M)$.  In practice this converts universal constructions on families into morphisms out of $M$.
\end{proposition}
\begin{example}[Relative picard functor viewpoint]
For relative Picard functor, take a base $T=\Spec R$ and replace a single object by an $R$-flat family.  Flatness prevents the numerical invariants from jumping on first-order thickenings, so it is the minimal condition under which the fibers can reasonably be compared.
\end{example}
\begin{theorem}[Relative picard functor viewpoint]
Suppose the objects in the relative Picard functor problem have no non-trivial automorphisms and admit effective descent for the chosen topology.  Then a universal family, if it exists, is unique up to unique isomorphism.  Consequently the representing object is unique up to unique isomorphism.
\end{theorem}
\begin{remark}[Relative picard functor viewpoint]
The phrase ``the moduli space of relative Picard functor'' is therefore shorthand.  It may mean a fine moduli scheme, a coarse moduli space, an algebraic stack, a GIT quotient, or a compactification.  The correct meaning is determined by what information is retained about families and automorphisms.
\end{remark}
\begin{proof}[Proof sketch of the theorem/proposition on this sheet]
The argument is the standard Yoneda-and-descent argument specialized to relative Picard functor.  A family over a test scheme $T$ corresponds, under representability, to a morphism $T\to M$.  Pulling the universal family back along this morphism recovers the given family.  If two morphisms give isomorphic pullbacks, the functorial equality on all test schemes forces the morphisms to agree.  The hypotheses on automorphisms ensure that the equality is not weakened to an equivalence relation in a groupoid.
\end{proof}
\begin{example}[A test calculation for relative Picard functor]
Let $T=\Spec k[s]$ and let the invariant of interest be denoted $\nu$.  A toy family $X\to T$ is chosen so that $\nu(X_t)=\nu_0$ for $s\ne0$ and the special fiber at $s=0$ records the boundary behavior.  In a concrete course one may take $X$ to be a family of divisors, modules, curves, bundles, or representations, depending on the chapter.  The moduli-theoretic lesson is that the morphism $T\to M$ detects both the general member and the limiting point.
\end{example}
\begin{exercise}
Write the functor of points for the moduli problem on this sheet and identify its value on $\Spec k$ and on $\Spec k[\epsilon]/(\epsilon^2)$.
\end{exercise}
\begin{exercise}
Formulate a small example in which the naive set of isomorphism classes is not Hausdorff, not separated, or not functorial.  Then repair the example by adding stability, rigidification, or a stack structure.
\end{exercise}
\clearpage
\section{Sheet 32: Line bundles and Picard functors: Jacobians}
This sheet studies \emph{Jacobians} as it appears in the broader chapter on \emph{line bundles and picard functors}.  The recurring question is whether a proposed parameter space remembers enough information about families while remaining geometric enough to compute with.
\begin{definition}[Jacobians viewpoint]
A Jacobians moduli functor assigns to each test scheme $T$ the set, groupoid, or category of families over $T$ satisfying the chosen condition.  Two families are identified only when the moduli problem says that an isomorphism, not merely a bijection on fibers, has been supplied.
\end{definition}
\begin{proposition}[Jacobians viewpoint]
If a moduli functor for Jacobians is represented by a scheme $M$, then natural transformations from the functor to any functor $G$ are the same as elements of $G(M)$.  In practice this converts universal constructions on families into morphisms out of $M$.
\end{proposition}
\begin{example}[Jacobians viewpoint]
For Jacobians, take a base $T=\Spec R$ and replace a single object by an $R$-flat family.  Flatness prevents the numerical invariants from jumping on first-order thickenings, so it is the minimal condition under which the fibers can reasonably be compared.
\end{example}
\begin{theorem}[Jacobians viewpoint]
Suppose the objects in the Jacobians problem have no non-trivial automorphisms and admit effective descent for the chosen topology.  Then a universal family, if it exists, is unique up to unique isomorphism.  Consequently the representing object is unique up to unique isomorphism.
\end{theorem}
\begin{remark}[Jacobians viewpoint]
The phrase ``the moduli space of Jacobians'' is therefore shorthand.  It may mean a fine moduli scheme, a coarse moduli space, an algebraic stack, a GIT quotient, or a compactification.  The correct meaning is determined by what information is retained about families and automorphisms.
\end{remark}
\begin{proof}[Proof sketch of the theorem/proposition on this sheet]
The argument is the standard Yoneda-and-descent argument specialized to Jacobians.  A family over a test scheme $T$ corresponds, under representability, to a morphism $T\to M$.  Pulling the universal family back along this morphism recovers the given family.  If two morphisms give isomorphic pullbacks, the functorial equality on all test schemes forces the morphisms to agree.  The hypotheses on automorphisms ensure that the equality is not weakened to an equivalence relation in a groupoid.
\end{proof}
\begin{example}[A test calculation for Jacobians]
Let $T=\Spec k[s]$ and let the invariant of interest be denoted $\nu$.  A toy family $X\to T$ is chosen so that $\nu(X_t)=\nu_0$ for $s\ne0$ and the special fiber at $s=0$ records the boundary behavior.  In a concrete course one may take $X$ to be a family of divisors, modules, curves, bundles, or representations, depending on the chapter.  The moduli-theoretic lesson is that the morphism $T\to M$ detects both the general member and the limiting point.
\end{example}
\begin{exercise}
Find one automorphism group occurring in the example and decide whether it obstructs the existence of a fine moduli scheme.
\end{exercise}
\begin{exercise}
Formulate a small example in which the naive set of isomorphism classes is not Hausdorff, not separated, or not functorial.  Then repair the example by adding stability, rigidification, or a stack structure.
\end{exercise}
\clearpage
\section{Sheet 33: Line bundles and Picard functors: Poincare bundles}
This sheet studies \emph{Poincare bundles} as it appears in the broader chapter on \emph{line bundles and picard functors}.  The recurring question is whether a proposed parameter space remembers enough information about families while remaining geometric enough to compute with.
\begin{definition}[Poincare bundles viewpoint]
A Poincare bundles moduli functor assigns to each test scheme $T$ the set, groupoid, or category of families over $T$ satisfying the chosen condition.  Two families are identified only when the moduli problem says that an isomorphism, not merely a bijection on fibers, has been supplied.
\end{definition}
\begin{proposition}[Poincare bundles viewpoint]
If a moduli functor for Poincare bundles is represented by a scheme $M$, then natural transformations from the functor to any functor $G$ are the same as elements of $G(M)$.  In practice this converts universal constructions on families into morphisms out of $M$.
\end{proposition}
\begin{example}[Poincare bundles viewpoint]
For Poincare bundles, take a base $T=\Spec R$ and replace a single object by an $R$-flat family.  Flatness prevents the numerical invariants from jumping on first-order thickenings, so it is the minimal condition under which the fibers can reasonably be compared.
\end{example}
\begin{theorem}[Poincare bundles viewpoint]
Suppose the objects in the Poincare bundles problem have no non-trivial automorphisms and admit effective descent for the chosen topology.  Then a universal family, if it exists, is unique up to unique isomorphism.  Consequently the representing object is unique up to unique isomorphism.
\end{theorem}
\begin{remark}[Poincare bundles viewpoint]
The phrase ``the moduli space of Poincare bundles'' is therefore shorthand.  It may mean a fine moduli scheme, a coarse moduli space, an algebraic stack, a GIT quotient, or a compactification.  The correct meaning is determined by what information is retained about families and automorphisms.
\end{remark}
\begin{proof}[Proof sketch of the theorem/proposition on this sheet]
The argument is the standard Yoneda-and-descent argument specialized to Poincare bundles.  A family over a test scheme $T$ corresponds, under representability, to a morphism $T\to M$.  Pulling the universal family back along this morphism recovers the given family.  If two morphisms give isomorphic pullbacks, the functorial equality on all test schemes forces the morphisms to agree.  The hypotheses on automorphisms ensure that the equality is not weakened to an equivalence relation in a groupoid.
\end{proof}
\begin{example}[A test calculation for Poincare bundles]
Let $T=\Spec k[s]$ and let the invariant of interest be denoted $\nu$.  A toy family $X\to T$ is chosen so that $\nu(X_t)=\nu_0$ for $s\ne0$ and the special fiber at $s=0$ records the boundary behavior.  In a concrete course one may take $X$ to be a family of divisors, modules, curves, bundles, or representations, depending on the chapter.  The moduli-theoretic lesson is that the morphism $T\to M$ detects both the general member and the limiting point.
\end{example}
\begin{exercise}
Check a base-change square explicitly: pull the displayed family back along $T\prime\to T$ and verify which invariant is preserved.
\end{exercise}
\begin{exercise}
Formulate a small example in which the naive set of isomorphism classes is not Hausdorff, not separated, or not functorial.  Then repair the example by adding stability, rigidification, or a stack structure.
\end{exercise}
\clearpage
\section{Sheet 34: Line bundles and Picard functors: Neron--Severi groups}
This sheet studies \emph{Neron--Severi groups} as it appears in the broader chapter on \emph{line bundles and picard functors}.  The recurring question is whether a proposed parameter space remembers enough information about families while remaining geometric enough to compute with.
\begin{definition}[Neron--severi groups viewpoint]
A Neron--Severi groups moduli functor assigns to each test scheme $T$ the set, groupoid, or category of families over $T$ satisfying the chosen condition.  Two families are identified only when the moduli problem says that an isomorphism, not merely a bijection on fibers, has been supplied.
\end{definition}
\begin{proposition}[Neron--severi groups viewpoint]
If a moduli functor for Neron--Severi groups is represented by a scheme $M$, then natural transformations from the functor to any functor $G$ are the same as elements of $G(M)$.  In practice this converts universal constructions on families into morphisms out of $M$.
\end{proposition}
\begin{example}[Neron--severi groups viewpoint]
For Neron--Severi groups, take a base $T=\Spec R$ and replace a single object by an $R$-flat family.  Flatness prevents the numerical invariants from jumping on first-order thickenings, so it is the minimal condition under which the fibers can reasonably be compared.
\end{example}
\begin{theorem}[Neron--severi groups viewpoint]
Suppose the objects in the Neron--Severi groups problem have no non-trivial automorphisms and admit effective descent for the chosen topology.  Then a universal family, if it exists, is unique up to unique isomorphism.  Consequently the representing object is unique up to unique isomorphism.
\end{theorem}
\begin{remark}[Neron--severi groups viewpoint]
The phrase ``the moduli space of Neron--Severi groups'' is therefore shorthand.  It may mean a fine moduli scheme, a coarse moduli space, an algebraic stack, a GIT quotient, or a compactification.  The correct meaning is determined by what information is retained about families and automorphisms.
\end{remark}
\begin{proof}[Proof sketch of the theorem/proposition on this sheet]
The argument is the standard Yoneda-and-descent argument specialized to Neron--Severi groups.  A family over a test scheme $T$ corresponds, under representability, to a morphism $T\to M$.  Pulling the universal family back along this morphism recovers the given family.  If two morphisms give isomorphic pullbacks, the functorial equality on all test schemes forces the morphisms to agree.  The hypotheses on automorphisms ensure that the equality is not weakened to an equivalence relation in a groupoid.
\end{proof}
\begin{example}[A test calculation for Neron--Severi groups]
Let $T=\Spec k[s]$ and let the invariant of interest be denoted $\nu$.  A toy family $X\to T$ is chosen so that $\nu(X_t)=\nu_0$ for $s\ne0$ and the special fiber at $s=0$ records the boundary behavior.  In a concrete course one may take $X$ to be a family of divisors, modules, curves, bundles, or representations, depending on the chapter.  The moduli-theoretic lesson is that the morphism $T\to M$ detects both the general member and the limiting point.
\end{example}
\begin{exercise}
Construct a degeneration over $\Spec k[[t]]$ whose generic fiber satisfies the condition but whose special fiber forces a compactification.
\end{exercise}
\begin{exercise}
Formulate a small example in which the naive set of isomorphism classes is not Hausdorff, not separated, or not functorial.  Then repair the example by adding stability, rigidification, or a stack structure.
\end{exercise}
\clearpage
\section{Sheet 35: Line bundles and Picard functors: Abel maps}
This sheet studies \emph{Abel maps} as it appears in the broader chapter on \emph{line bundles and picard functors}.  The recurring question is whether a proposed parameter space remembers enough information about families while remaining geometric enough to compute with.
\begin{definition}[Abel maps viewpoint]
A Abel maps moduli functor assigns to each test scheme $T$ the set, groupoid, or category of families over $T$ satisfying the chosen condition.  Two families are identified only when the moduli problem says that an isomorphism, not merely a bijection on fibers, has been supplied.
\end{definition}
\begin{proposition}[Abel maps viewpoint]
If a moduli functor for Abel maps is represented by a scheme $M$, then natural transformations from the functor to any functor $G$ are the same as elements of $G(M)$.  In practice this converts universal constructions on families into morphisms out of $M$.
\end{proposition}
\begin{example}[Abel maps viewpoint]
For Abel maps, take a base $T=\Spec R$ and replace a single object by an $R$-flat family.  Flatness prevents the numerical invariants from jumping on first-order thickenings, so it is the minimal condition under which the fibers can reasonably be compared.
\end{example}
\begin{theorem}[Abel maps viewpoint]
Suppose the objects in the Abel maps problem have no non-trivial automorphisms and admit effective descent for the chosen topology.  Then a universal family, if it exists, is unique up to unique isomorphism.  Consequently the representing object is unique up to unique isomorphism.
\end{theorem}
\begin{remark}[Abel maps viewpoint]
The phrase ``the moduli space of Abel maps'' is therefore shorthand.  It may mean a fine moduli scheme, a coarse moduli space, an algebraic stack, a GIT quotient, or a compactification.  The correct meaning is determined by what information is retained about families and automorphisms.
\end{remark}
\begin{proof}[Proof sketch of the theorem/proposition on this sheet]
The argument is the standard Yoneda-and-descent argument specialized to Abel maps.  A family over a test scheme $T$ corresponds, under representability, to a morphism $T\to M$.  Pulling the universal family back along this morphism recovers the given family.  If two morphisms give isomorphic pullbacks, the functorial equality on all test schemes forces the morphisms to agree.  The hypotheses on automorphisms ensure that the equality is not weakened to an equivalence relation in a groupoid.
\end{proof}
\begin{example}[A test calculation for Abel maps]
Let $T=\Spec k[s]$ and let the invariant of interest be denoted $\nu$.  A toy family $X\to T$ is chosen so that $\nu(X_t)=\nu_0$ for $s\ne0$ and the special fiber at $s=0$ records the boundary behavior.  In a concrete course one may take $X$ to be a family of divisors, modules, curves, bundles, or representations, depending on the chapter.  The moduli-theoretic lesson is that the morphism $T\to M$ detects both the general member and the limiting point.
\end{example}
\begin{exercise}
Compare the functor with the corresponding groupoid-valued functor.  Which information is lost after passing to isomorphism classes?
\end{exercise}
\begin{exercise}
Formulate a small example in which the naive set of isomorphism classes is not Hausdorff, not separated, or not functorial.  Then repair the example by adding stability, rigidification, or a stack structure.
\end{exercise}
\clearpage
\section{Sheet 36: Stable maps: Kontsevich stability}
This sheet studies \emph{Kontsevich stability} as it appears in the broader chapter on \emph{stable maps}.  The recurring question is whether a proposed parameter space remembers enough information about families while remaining geometric enough to compute with.
\begin{definition}[Kontsevich stability viewpoint]
A Kontsevich stability moduli functor assigns to each test scheme $T$ the set, groupoid, or category of families over $T$ satisfying the chosen condition.  Two families are identified only when the moduli problem says that an isomorphism, not merely a bijection on fibers, has been supplied.
\end{definition}
\begin{proposition}[Kontsevich stability viewpoint]
If a moduli functor for Kontsevich stability is represented by a scheme $M$, then natural transformations from the functor to any functor $G$ are the same as elements of $G(M)$.  In practice this converts universal constructions on families into morphisms out of $M$.
\end{proposition}
\begin{example}[Kontsevich stability viewpoint]
For Kontsevich stability, take a base $T=\Spec R$ and replace a single object by an $R$-flat family.  Flatness prevents the numerical invariants from jumping on first-order thickenings, so it is the minimal condition under which the fibers can reasonably be compared.
\end{example}
\begin{theorem}[Kontsevich stability viewpoint]
Suppose the objects in the Kontsevich stability problem have no non-trivial automorphisms and admit effective descent for the chosen topology.  Then a universal family, if it exists, is unique up to unique isomorphism.  Consequently the representing object is unique up to unique isomorphism.
\end{theorem}
\begin{remark}[Kontsevich stability viewpoint]
The phrase ``the moduli space of Kontsevich stability'' is therefore shorthand.  It may mean a fine moduli scheme, a coarse moduli space, an algebraic stack, a GIT quotient, or a compactification.  The correct meaning is determined by what information is retained about families and automorphisms.
\end{remark}
\begin{proof}[Proof sketch of the theorem/proposition on this sheet]
The argument is the standard Yoneda-and-descent argument specialized to Kontsevich stability.  A family over a test scheme $T$ corresponds, under representability, to a morphism $T\to M$.  Pulling the universal family back along this morphism recovers the given family.  If two morphisms give isomorphic pullbacks, the functorial equality on all test schemes forces the morphisms to agree.  The hypotheses on automorphisms ensure that the equality is not weakened to an equivalence relation in a groupoid.
\end{proof}
\begin{example}[A test calculation for Kontsevich stability]
Let $T=\Spec k[s]$ and let the invariant of interest be denoted $\nu$.  A toy family $X\to T$ is chosen so that $\nu(X_t)=\nu_0$ for $s\ne0$ and the special fiber at $s=0$ records the boundary behavior.  In a concrete course one may take $X$ to be a family of divisors, modules, curves, bundles, or representations, depending on the chapter.  The moduli-theoretic lesson is that the morphism $T\to M$ detects both the general member and the limiting point.
\end{example}
\begin{exercise}
Write the functor of points for the moduli problem on this sheet and identify its value on $\Spec k$ and on $\Spec k[\epsilon]/(\epsilon^2)$.
\end{exercise}
\begin{exercise}
Formulate a small example in which the naive set of isomorphism classes is not Hausdorff, not separated, or not functorial.  Then repair the example by adding stability, rigidification, or a stack structure.
\end{exercise}
\clearpage
\section{Sheet 37: Stable maps: evaluation maps}
This sheet studies \emph{evaluation maps} as it appears in the broader chapter on \emph{stable maps}.  The recurring question is whether a proposed parameter space remembers enough information about families while remaining geometric enough to compute with.
\begin{definition}[Evaluation maps viewpoint]
A evaluation maps moduli functor assigns to each test scheme $T$ the set, groupoid, or category of families over $T$ satisfying the chosen condition.  Two families are identified only when the moduli problem says that an isomorphism, not merely a bijection on fibers, has been supplied.
\end{definition}
\begin{proposition}[Evaluation maps viewpoint]
If a moduli functor for evaluation maps is represented by a scheme $M$, then natural transformations from the functor to any functor $G$ are the same as elements of $G(M)$.  In practice this converts universal constructions on families into morphisms out of $M$.
\end{proposition}
\begin{example}[Evaluation maps viewpoint]
For evaluation maps, take a base $T=\Spec R$ and replace a single object by an $R$-flat family.  Flatness prevents the numerical invariants from jumping on first-order thickenings, so it is the minimal condition under which the fibers can reasonably be compared.
\end{example}
\begin{theorem}[Evaluation maps viewpoint]
Suppose the objects in the evaluation maps problem have no non-trivial automorphisms and admit effective descent for the chosen topology.  Then a universal family, if it exists, is unique up to unique isomorphism.  Consequently the representing object is unique up to unique isomorphism.
\end{theorem}
\begin{remark}[Evaluation maps viewpoint]
The phrase ``the moduli space of evaluation maps'' is therefore shorthand.  It may mean a fine moduli scheme, a coarse moduli space, an algebraic stack, a GIT quotient, or a compactification.  The correct meaning is determined by what information is retained about families and automorphisms.
\end{remark}
\begin{proof}[Proof sketch of the theorem/proposition on this sheet]
The argument is the standard Yoneda-and-descent argument specialized to evaluation maps.  A family over a test scheme $T$ corresponds, under representability, to a morphism $T\to M$.  Pulling the universal family back along this morphism recovers the given family.  If two morphisms give isomorphic pullbacks, the functorial equality on all test schemes forces the morphisms to agree.  The hypotheses on automorphisms ensure that the equality is not weakened to an equivalence relation in a groupoid.
\end{proof}
\begin{example}[A test calculation for evaluation maps]
Let $T=\Spec k[s]$ and let the invariant of interest be denoted $\nu$.  A toy family $X\to T$ is chosen so that $\nu(X_t)=\nu_0$ for $s\ne0$ and the special fiber at $s=0$ records the boundary behavior.  In a concrete course one may take $X$ to be a family of divisors, modules, curves, bundles, or representations, depending on the chapter.  The moduli-theoretic lesson is that the morphism $T\to M$ detects both the general member and the limiting point.
\end{example}
\begin{exercise}
Find one automorphism group occurring in the example and decide whether it obstructs the existence of a fine moduli scheme.
\end{exercise}
\begin{exercise}
Formulate a small example in which the naive set of isomorphism classes is not Hausdorff, not separated, or not functorial.  Then repair the example by adding stability, rigidification, or a stack structure.
\end{exercise}
\clearpage
\section{Sheet 38: Stable maps: virtual dimension}
This sheet studies \emph{virtual dimension} as it appears in the broader chapter on \emph{stable maps}.  The recurring question is whether a proposed parameter space remembers enough information about families while remaining geometric enough to compute with.
\begin{definition}[Virtual dimension viewpoint]
A virtual dimension moduli functor assigns to each test scheme $T$ the set, groupoid, or category of families over $T$ satisfying the chosen condition.  Two families are identified only when the moduli problem says that an isomorphism, not merely a bijection on fibers, has been supplied.
\end{definition}
\begin{proposition}[Virtual dimension viewpoint]
If a moduli functor for virtual dimension is represented by a scheme $M$, then natural transformations from the functor to any functor $G$ are the same as elements of $G(M)$.  In practice this converts universal constructions on families into morphisms out of $M$.
\end{proposition}
\begin{example}[Virtual dimension viewpoint]
For virtual dimension, take a base $T=\Spec R$ and replace a single object by an $R$-flat family.  Flatness prevents the numerical invariants from jumping on first-order thickenings, so it is the minimal condition under which the fibers can reasonably be compared.
\end{example}
\begin{theorem}[Virtual dimension viewpoint]
Suppose the objects in the virtual dimension problem have no non-trivial automorphisms and admit effective descent for the chosen topology.  Then a universal family, if it exists, is unique up to unique isomorphism.  Consequently the representing object is unique up to unique isomorphism.
\end{theorem}
\begin{remark}[Virtual dimension viewpoint]
The phrase ``the moduli space of virtual dimension'' is therefore shorthand.  It may mean a fine moduli scheme, a coarse moduli space, an algebraic stack, a GIT quotient, or a compactification.  The correct meaning is determined by what information is retained about families and automorphisms.
\end{remark}
\begin{proof}[Proof sketch of the theorem/proposition on this sheet]
The argument is the standard Yoneda-and-descent argument specialized to virtual dimension.  A family over a test scheme $T$ corresponds, under representability, to a morphism $T\to M$.  Pulling the universal family back along this morphism recovers the given family.  If two morphisms give isomorphic pullbacks, the functorial equality on all test schemes forces the morphisms to agree.  The hypotheses on automorphisms ensure that the equality is not weakened to an equivalence relation in a groupoid.
\end{proof}
\begin{example}[A test calculation for virtual dimension]
Let $T=\Spec k[s]$ and let the invariant of interest be denoted $\nu$.  A toy family $X\to T$ is chosen so that $\nu(X_t)=\nu_0$ for $s\ne0$ and the special fiber at $s=0$ records the boundary behavior.  In a concrete course one may take $X$ to be a family of divisors, modules, curves, bundles, or representations, depending on the chapter.  The moduli-theoretic lesson is that the morphism $T\to M$ detects both the general member and the limiting point.
\end{example}
\begin{exercise}
Check a base-change square explicitly: pull the displayed family back along $T\prime\to T$ and verify which invariant is preserved.
\end{exercise}
\begin{exercise}
Formulate a small example in which the naive set of isomorphism classes is not Hausdorff, not separated, or not functorial.  Then repair the example by adding stability, rigidification, or a stack structure.
\end{exercise}
\clearpage
\section{Sheet 39: Stable maps: Gromov--Witten invariants}
This sheet studies \emph{Gromov--Witten invariants} as it appears in the broader chapter on \emph{stable maps}.  The recurring question is whether a proposed parameter space remembers enough information about families while remaining geometric enough to compute with.
\begin{definition}[Gromov--witten invariants viewpoint]
A Gromov--Witten invariants moduli functor assigns to each test scheme $T$ the set, groupoid, or category of families over $T$ satisfying the chosen condition.  Two families are identified only when the moduli problem says that an isomorphism, not merely a bijection on fibers, has been supplied.
\end{definition}
\begin{proposition}[Gromov--witten invariants viewpoint]
If a moduli functor for Gromov--Witten invariants is represented by a scheme $M$, then natural transformations from the functor to any functor $G$ are the same as elements of $G(M)$.  In practice this converts universal constructions on families into morphisms out of $M$.
\end{proposition}
\begin{example}[Gromov--witten invariants viewpoint]
For Gromov--Witten invariants, take a base $T=\Spec R$ and replace a single object by an $R$-flat family.  Flatness prevents the numerical invariants from jumping on first-order thickenings, so it is the minimal condition under which the fibers can reasonably be compared.
\end{example}
\begin{theorem}[Gromov--witten invariants viewpoint]
Suppose the objects in the Gromov--Witten invariants problem have no non-trivial automorphisms and admit effective descent for the chosen topology.  Then a universal family, if it exists, is unique up to unique isomorphism.  Consequently the representing object is unique up to unique isomorphism.
\end{theorem}
\begin{remark}[Gromov--witten invariants viewpoint]
The phrase ``the moduli space of Gromov--Witten invariants'' is therefore shorthand.  It may mean a fine moduli scheme, a coarse moduli space, an algebraic stack, a GIT quotient, or a compactification.  The correct meaning is determined by what information is retained about families and automorphisms.
\end{remark}
\begin{proof}[Proof sketch of the theorem/proposition on this sheet]
The argument is the standard Yoneda-and-descent argument specialized to Gromov--Witten invariants.  A family over a test scheme $T$ corresponds, under representability, to a morphism $T\to M$.  Pulling the universal family back along this morphism recovers the given family.  If two morphisms give isomorphic pullbacks, the functorial equality on all test schemes forces the morphisms to agree.  The hypotheses on automorphisms ensure that the equality is not weakened to an equivalence relation in a groupoid.
\end{proof}
\begin{example}[A test calculation for Gromov--Witten invariants]
Let $T=\Spec k[s]$ and let the invariant of interest be denoted $\nu$.  A toy family $X\to T$ is chosen so that $\nu(X_t)=\nu_0$ for $s\ne0$ and the special fiber at $s=0$ records the boundary behavior.  In a concrete course one may take $X$ to be a family of divisors, modules, curves, bundles, or representations, depending on the chapter.  The moduli-theoretic lesson is that the morphism $T\to M$ detects both the general member and the limiting point.
\end{example}
\begin{exercise}
Construct a degeneration over $\Spec k[[t]]$ whose generic fiber satisfies the condition but whose special fiber forces a compactification.
\end{exercise}
\begin{exercise}
Formulate a small example in which the naive set of isomorphism classes is not Hausdorff, not separated, or not functorial.  Then repair the example by adding stability, rigidification, or a stack structure.
\end{exercise}
\clearpage
\section{Sheet 40: Stable maps: forgetful morphisms}
This sheet studies \emph{forgetful morphisms} as it appears in the broader chapter on \emph{stable maps}.  The recurring question is whether a proposed parameter space remembers enough information about families while remaining geometric enough to compute with.
\begin{definition}[Forgetful morphisms viewpoint]
A forgetful morphisms moduli functor assigns to each test scheme $T$ the set, groupoid, or category of families over $T$ satisfying the chosen condition.  Two families are identified only when the moduli problem says that an isomorphism, not merely a bijection on fibers, has been supplied.
\end{definition}
\begin{proposition}[Forgetful morphisms viewpoint]
If a moduli functor for forgetful morphisms is represented by a scheme $M$, then natural transformations from the functor to any functor $G$ are the same as elements of $G(M)$.  In practice this converts universal constructions on families into morphisms out of $M$.
\end{proposition}
\begin{example}[Forgetful morphisms viewpoint]
For forgetful morphisms, take a base $T=\Spec R$ and replace a single object by an $R$-flat family.  Flatness prevents the numerical invariants from jumping on first-order thickenings, so it is the minimal condition under which the fibers can reasonably be compared.
\end{example}
\begin{theorem}[Forgetful morphisms viewpoint]
Suppose the objects in the forgetful morphisms problem have no non-trivial automorphisms and admit effective descent for the chosen topology.  Then a universal family, if it exists, is unique up to unique isomorphism.  Consequently the representing object is unique up to unique isomorphism.
\end{theorem}
\begin{remark}[Forgetful morphisms viewpoint]
The phrase ``the moduli space of forgetful morphisms'' is therefore shorthand.  It may mean a fine moduli scheme, a coarse moduli space, an algebraic stack, a GIT quotient, or a compactification.  The correct meaning is determined by what information is retained about families and automorphisms.
\end{remark}
\begin{proof}[Proof sketch of the theorem/proposition on this sheet]
The argument is the standard Yoneda-and-descent argument specialized to forgetful morphisms.  A family over a test scheme $T$ corresponds, under representability, to a morphism $T\to M$.  Pulling the universal family back along this morphism recovers the given family.  If two morphisms give isomorphic pullbacks, the functorial equality on all test schemes forces the morphisms to agree.  The hypotheses on automorphisms ensure that the equality is not weakened to an equivalence relation in a groupoid.
\end{proof}
\begin{example}[A test calculation for forgetful morphisms]
Let $T=\Spec k[s]$ and let the invariant of interest be denoted $\nu$.  A toy family $X\to T$ is chosen so that $\nu(X_t)=\nu_0$ for $s\ne0$ and the special fiber at $s=0$ records the boundary behavior.  In a concrete course one may take $X$ to be a family of divisors, modules, curves, bundles, or representations, depending on the chapter.  The moduli-theoretic lesson is that the morphism $T\to M$ detects both the general member and the limiting point.
\end{example}
\begin{exercise}
Compare the functor with the corresponding groupoid-valued functor.  Which information is lost after passing to isomorphism classes?
\end{exercise}
\begin{exercise}
Formulate a small example in which the naive set of isomorphism classes is not Hausdorff, not separated, or not functorial.  Then repair the example by adding stability, rigidification, or a stack structure.
\end{exercise}
\clearpage
\section{Sheet 41: Stacks: fibered categories}
This sheet studies \emph{fibered categories} as it appears in the broader chapter on \emph{stacks}.  The recurring question is whether a proposed parameter space remembers enough information about families while remaining geometric enough to compute with.
\begin{definition}[Fibered categories viewpoint]
A fibered categories moduli functor assigns to each test scheme $T$ the set, groupoid, or category of families over $T$ satisfying the chosen condition.  Two families are identified only when the moduli problem says that an isomorphism, not merely a bijection on fibers, has been supplied.
\end{definition}
\begin{proposition}[Fibered categories viewpoint]
If a moduli functor for fibered categories is represented by a scheme $M$, then natural transformations from the functor to any functor $G$ are the same as elements of $G(M)$.  In practice this converts universal constructions on families into morphisms out of $M$.
\end{proposition}
\begin{example}[Fibered categories viewpoint]
For fibered categories, take a base $T=\Spec R$ and replace a single object by an $R$-flat family.  Flatness prevents the numerical invariants from jumping on first-order thickenings, so it is the minimal condition under which the fibers can reasonably be compared.
\end{example}
\begin{theorem}[Fibered categories viewpoint]
Suppose the objects in the fibered categories problem have no non-trivial automorphisms and admit effective descent for the chosen topology.  Then a universal family, if it exists, is unique up to unique isomorphism.  Consequently the representing object is unique up to unique isomorphism.
\end{theorem}
\begin{remark}[Fibered categories viewpoint]
The phrase ``the moduli space of fibered categories'' is therefore shorthand.  It may mean a fine moduli scheme, a coarse moduli space, an algebraic stack, a GIT quotient, or a compactification.  The correct meaning is determined by what information is retained about families and automorphisms.
\end{remark}
\begin{proof}[Proof sketch of the theorem/proposition on this sheet]
The argument is the standard Yoneda-and-descent argument specialized to fibered categories.  A family over a test scheme $T$ corresponds, under representability, to a morphism $T\to M$.  Pulling the universal family back along this morphism recovers the given family.  If two morphisms give isomorphic pullbacks, the functorial equality on all test schemes forces the morphisms to agree.  The hypotheses on automorphisms ensure that the equality is not weakened to an equivalence relation in a groupoid.
\end{proof}
\begin{example}[A test calculation for fibered categories]
Let $T=\Spec k[s]$ and let the invariant of interest be denoted $\nu$.  A toy family $X\to T$ is chosen so that $\nu(X_t)=\nu_0$ for $s\ne0$ and the special fiber at $s=0$ records the boundary behavior.  In a concrete course one may take $X$ to be a family of divisors, modules, curves, bundles, or representations, depending on the chapter.  The moduli-theoretic lesson is that the morphism $T\to M$ detects both the general member and the limiting point.
\end{example}
\begin{exercise}
Write the functor of points for the moduli problem on this sheet and identify its value on $\Spec k$ and on $\Spec k[\epsilon]/(\epsilon^2)$.
\end{exercise}
\begin{exercise}
Formulate a small example in which the naive set of isomorphism classes is not Hausdorff, not separated, or not functorial.  Then repair the example by adding stability, rigidification, or a stack structure.
\end{exercise}
\clearpage
\section{Sheet 42: Stacks: descent}
This sheet studies \emph{descent} as it appears in the broader chapter on \emph{stacks}.  The recurring question is whether a proposed parameter space remembers enough information about families while remaining geometric enough to compute with.
\begin{definition}[Descent viewpoint]
A descent moduli functor assigns to each test scheme $T$ the set, groupoid, or category of families over $T$ satisfying the chosen condition.  Two families are identified only when the moduli problem says that an isomorphism, not merely a bijection on fibers, has been supplied.
\end{definition}
\begin{proposition}[Descent viewpoint]
If a moduli functor for descent is represented by a scheme $M$, then natural transformations from the functor to any functor $G$ are the same as elements of $G(M)$.  In practice this converts universal constructions on families into morphisms out of $M$.
\end{proposition}
\begin{example}[Descent viewpoint]
For descent, take a base $T=\Spec R$ and replace a single object by an $R$-flat family.  Flatness prevents the numerical invariants from jumping on first-order thickenings, so it is the minimal condition under which the fibers can reasonably be compared.
\end{example}
\begin{theorem}[Descent viewpoint]
Suppose the objects in the descent problem have no non-trivial automorphisms and admit effective descent for the chosen topology.  Then a universal family, if it exists, is unique up to unique isomorphism.  Consequently the representing object is unique up to unique isomorphism.
\end{theorem}
\begin{remark}[Descent viewpoint]
The phrase ``the moduli space of descent'' is therefore shorthand.  It may mean a fine moduli scheme, a coarse moduli space, an algebraic stack, a GIT quotient, or a compactification.  The correct meaning is determined by what information is retained about families and automorphisms.
\end{remark}
\begin{proof}[Proof sketch of the theorem/proposition on this sheet]
The argument is the standard Yoneda-and-descent argument specialized to descent.  A family over a test scheme $T$ corresponds, under representability, to a morphism $T\to M$.  Pulling the universal family back along this morphism recovers the given family.  If two morphisms give isomorphic pullbacks, the functorial equality on all test schemes forces the morphisms to agree.  The hypotheses on automorphisms ensure that the equality is not weakened to an equivalence relation in a groupoid.
\end{proof}
\begin{example}[A test calculation for descent]
Let $T=\Spec k[s]$ and let the invariant of interest be denoted $\nu$.  A toy family $X\to T$ is chosen so that $\nu(X_t)=\nu_0$ for $s\ne0$ and the special fiber at $s=0$ records the boundary behavior.  In a concrete course one may take $X$ to be a family of divisors, modules, curves, bundles, or representations, depending on the chapter.  The moduli-theoretic lesson is that the morphism $T\to M$ detects both the general member and the limiting point.
\end{example}
\begin{exercise}
Find one automorphism group occurring in the example and decide whether it obstructs the existence of a fine moduli scheme.
\end{exercise}
\begin{exercise}
Formulate a small example in which the naive set of isomorphism classes is not Hausdorff, not separated, or not functorial.  Then repair the example by adding stability, rigidification, or a stack structure.
\end{exercise}
\clearpage
\section{Sheet 43: Stacks: algebraic stacks}
This sheet studies \emph{algebraic stacks} as it appears in the broader chapter on \emph{stacks}.  The recurring question is whether a proposed parameter space remembers enough information about families while remaining geometric enough to compute with.
\begin{definition}[Algebraic stacks viewpoint]
A algebraic stacks moduli functor assigns to each test scheme $T$ the set, groupoid, or category of families over $T$ satisfying the chosen condition.  Two families are identified only when the moduli problem says that an isomorphism, not merely a bijection on fibers, has been supplied.
\end{definition}
\begin{proposition}[Algebraic stacks viewpoint]
If a moduli functor for algebraic stacks is represented by a scheme $M$, then natural transformations from the functor to any functor $G$ are the same as elements of $G(M)$.  In practice this converts universal constructions on families into morphisms out of $M$.
\end{proposition}
\begin{example}[Algebraic stacks viewpoint]
For algebraic stacks, take a base $T=\Spec R$ and replace a single object by an $R$-flat family.  Flatness prevents the numerical invariants from jumping on first-order thickenings, so it is the minimal condition under which the fibers can reasonably be compared.
\end{example}
\begin{theorem}[Algebraic stacks viewpoint]
Suppose the objects in the algebraic stacks problem have no non-trivial automorphisms and admit effective descent for the chosen topology.  Then a universal family, if it exists, is unique up to unique isomorphism.  Consequently the representing object is unique up to unique isomorphism.
\end{theorem}
\begin{remark}[Algebraic stacks viewpoint]
The phrase ``the moduli space of algebraic stacks'' is therefore shorthand.  It may mean a fine moduli scheme, a coarse moduli space, an algebraic stack, a GIT quotient, or a compactification.  The correct meaning is determined by what information is retained about families and automorphisms.
\end{remark}
\begin{proof}[Proof sketch of the theorem/proposition on this sheet]
The argument is the standard Yoneda-and-descent argument specialized to algebraic stacks.  A family over a test scheme $T$ corresponds, under representability, to a morphism $T\to M$.  Pulling the universal family back along this morphism recovers the given family.  If two morphisms give isomorphic pullbacks, the functorial equality on all test schemes forces the morphisms to agree.  The hypotheses on automorphisms ensure that the equality is not weakened to an equivalence relation in a groupoid.
\end{proof}
\begin{example}[A test calculation for algebraic stacks]
Let $T=\Spec k[s]$ and let the invariant of interest be denoted $\nu$.  A toy family $X\to T$ is chosen so that $\nu(X_t)=\nu_0$ for $s\ne0$ and the special fiber at $s=0$ records the boundary behavior.  In a concrete course one may take $X$ to be a family of divisors, modules, curves, bundles, or representations, depending on the chapter.  The moduli-theoretic lesson is that the morphism $T\to M$ detects both the general member and the limiting point.
\end{example}
\begin{exercise}
Check a base-change square explicitly: pull the displayed family back along $T\prime\to T$ and verify which invariant is preserved.
\end{exercise}
\begin{exercise}
Formulate a small example in which the naive set of isomorphism classes is not Hausdorff, not separated, or not functorial.  Then repair the example by adding stability, rigidification, or a stack structure.
\end{exercise}
\clearpage
\section{Sheet 44: Stacks: inertia}
This sheet studies \emph{inertia} as it appears in the broader chapter on \emph{stacks}.  The recurring question is whether a proposed parameter space remembers enough information about families while remaining geometric enough to compute with.
\begin{definition}[Inertia viewpoint]
A inertia moduli functor assigns to each test scheme $T$ the set, groupoid, or category of families over $T$ satisfying the chosen condition.  Two families are identified only when the moduli problem says that an isomorphism, not merely a bijection on fibers, has been supplied.
\end{definition}
\begin{proposition}[Inertia viewpoint]
If a moduli functor for inertia is represented by a scheme $M$, then natural transformations from the functor to any functor $G$ are the same as elements of $G(M)$.  In practice this converts universal constructions on families into morphisms out of $M$.
\end{proposition}
\begin{example}[Inertia viewpoint]
For inertia, take a base $T=\Spec R$ and replace a single object by an $R$-flat family.  Flatness prevents the numerical invariants from jumping on first-order thickenings, so it is the minimal condition under which the fibers can reasonably be compared.
\end{example}
\begin{theorem}[Inertia viewpoint]
Suppose the objects in the inertia problem have no non-trivial automorphisms and admit effective descent for the chosen topology.  Then a universal family, if it exists, is unique up to unique isomorphism.  Consequently the representing object is unique up to unique isomorphism.
\end{theorem}
\begin{remark}[Inertia viewpoint]
The phrase ``the moduli space of inertia'' is therefore shorthand.  It may mean a fine moduli scheme, a coarse moduli space, an algebraic stack, a GIT quotient, or a compactification.  The correct meaning is determined by what information is retained about families and automorphisms.
\end{remark}
\begin{proof}[Proof sketch of the theorem/proposition on this sheet]
The argument is the standard Yoneda-and-descent argument specialized to inertia.  A family over a test scheme $T$ corresponds, under representability, to a morphism $T\to M$.  Pulling the universal family back along this morphism recovers the given family.  If two morphisms give isomorphic pullbacks, the functorial equality on all test schemes forces the morphisms to agree.  The hypotheses on automorphisms ensure that the equality is not weakened to an equivalence relation in a groupoid.
\end{proof}
\begin{example}[A test calculation for inertia]
Let $T=\Spec k[s]$ and let the invariant of interest be denoted $\nu$.  A toy family $X\to T$ is chosen so that $\nu(X_t)=\nu_0$ for $s\ne0$ and the special fiber at $s=0$ records the boundary behavior.  In a concrete course one may take $X$ to be a family of divisors, modules, curves, bundles, or representations, depending on the chapter.  The moduli-theoretic lesson is that the morphism $T\to M$ detects both the general member and the limiting point.
\end{example}
\begin{exercise}
Construct a degeneration over $\Spec k[[t]]$ whose generic fiber satisfies the condition but whose special fiber forces a compactification.
\end{exercise}
\begin{exercise}
Formulate a small example in which the naive set of isomorphism classes is not Hausdorff, not separated, or not functorial.  Then repair the example by adding stability, rigidification, or a stack structure.
\end{exercise}
\clearpage
\section{Sheet 45: Stacks: Deligne--Mumford stacks}
This sheet studies \emph{Deligne--Mumford stacks} as it appears in the broader chapter on \emph{stacks}.  The recurring question is whether a proposed parameter space remembers enough information about families while remaining geometric enough to compute with.
\begin{definition}[Deligne--mumford stacks viewpoint]
A Deligne--Mumford stacks moduli functor assigns to each test scheme $T$ the set, groupoid, or category of families over $T$ satisfying the chosen condition.  Two families are identified only when the moduli problem says that an isomorphism, not merely a bijection on fibers, has been supplied.
\end{definition}
\begin{proposition}[Deligne--mumford stacks viewpoint]
If a moduli functor for Deligne--Mumford stacks is represented by a scheme $M$, then natural transformations from the functor to any functor $G$ are the same as elements of $G(M)$.  In practice this converts universal constructions on families into morphisms out of $M$.
\end{proposition}
\begin{example}[Deligne--mumford stacks viewpoint]
For Deligne--Mumford stacks, take a base $T=\Spec R$ and replace a single object by an $R$-flat family.  Flatness prevents the numerical invariants from jumping on first-order thickenings, so it is the minimal condition under which the fibers can reasonably be compared.
\end{example}
\begin{theorem}[Deligne--mumford stacks viewpoint]
Suppose the objects in the Deligne--Mumford stacks problem have no non-trivial automorphisms and admit effective descent for the chosen topology.  Then a universal family, if it exists, is unique up to unique isomorphism.  Consequently the representing object is unique up to unique isomorphism.
\end{theorem}
\begin{remark}[Deligne--mumford stacks viewpoint]
The phrase ``the moduli space of Deligne--Mumford stacks'' is therefore shorthand.  It may mean a fine moduli scheme, a coarse moduli space, an algebraic stack, a GIT quotient, or a compactification.  The correct meaning is determined by what information is retained about families and automorphisms.
\end{remark}
\begin{proof}[Proof sketch of the theorem/proposition on this sheet]
The argument is the standard Yoneda-and-descent argument specialized to Deligne--Mumford stacks.  A family over a test scheme $T$ corresponds, under representability, to a morphism $T\to M$.  Pulling the universal family back along this morphism recovers the given family.  If two morphisms give isomorphic pullbacks, the functorial equality on all test schemes forces the morphisms to agree.  The hypotheses on automorphisms ensure that the equality is not weakened to an equivalence relation in a groupoid.
\end{proof}
\begin{example}[A test calculation for Deligne--Mumford stacks]
Let $T=\Spec k[s]$ and let the invariant of interest be denoted $\nu$.  A toy family $X\to T$ is chosen so that $\nu(X_t)=\nu_0$ for $s\ne0$ and the special fiber at $s=0$ records the boundary behavior.  In a concrete course one may take $X$ to be a family of divisors, modules, curves, bundles, or representations, depending on the chapter.  The moduli-theoretic lesson is that the morphism $T\to M$ detects both the general member and the limiting point.
\end{example}
\begin{exercise}
Compare the functor with the corresponding groupoid-valued functor.  Which information is lost after passing to isomorphism classes?
\end{exercise}
\begin{exercise}
Formulate a small example in which the naive set of isomorphism classes is not Hausdorff, not separated, or not functorial.  Then repair the example by adding stability, rigidification, or a stack structure.
\end{exercise}
\clearpage
\section{Sheet 46: Geometric invariant theory: linearizations}
This sheet studies \emph{linearizations} as it appears in the broader chapter on \emph{geometric invariant theory}.  The recurring question is whether a proposed parameter space remembers enough information about families while remaining geometric enough to compute with.
\begin{definition}[Linearizations viewpoint]
A linearizations moduli functor assigns to each test scheme $T$ the set, groupoid, or category of families over $T$ satisfying the chosen condition.  Two families are identified only when the moduli problem says that an isomorphism, not merely a bijection on fibers, has been supplied.
\end{definition}
\begin{proposition}[Linearizations viewpoint]
If a moduli functor for linearizations is represented by a scheme $M$, then natural transformations from the functor to any functor $G$ are the same as elements of $G(M)$.  In practice this converts universal constructions on families into morphisms out of $M$.
\end{proposition}
\begin{example}[Linearizations viewpoint]
For linearizations, take a base $T=\Spec R$ and replace a single object by an $R$-flat family.  Flatness prevents the numerical invariants from jumping on first-order thickenings, so it is the minimal condition under which the fibers can reasonably be compared.
\end{example}
\begin{theorem}[Linearizations viewpoint]
Suppose the objects in the linearizations problem have no non-trivial automorphisms and admit effective descent for the chosen topology.  Then a universal family, if it exists, is unique up to unique isomorphism.  Consequently the representing object is unique up to unique isomorphism.
\end{theorem}
\begin{remark}[Linearizations viewpoint]
The phrase ``the moduli space of linearizations'' is therefore shorthand.  It may mean a fine moduli scheme, a coarse moduli space, an algebraic stack, a GIT quotient, or a compactification.  The correct meaning is determined by what information is retained about families and automorphisms.
\end{remark}
\begin{proof}[Proof sketch of the theorem/proposition on this sheet]
The argument is the standard Yoneda-and-descent argument specialized to linearizations.  A family over a test scheme $T$ corresponds, under representability, to a morphism $T\to M$.  Pulling the universal family back along this morphism recovers the given family.  If two morphisms give isomorphic pullbacks, the functorial equality on all test schemes forces the morphisms to agree.  The hypotheses on automorphisms ensure that the equality is not weakened to an equivalence relation in a groupoid.
\end{proof}
\begin{example}[A test calculation for linearizations]
Let $T=\Spec k[s]$ and let the invariant of interest be denoted $\nu$.  A toy family $X\to T$ is chosen so that $\nu(X_t)=\nu_0$ for $s\ne0$ and the special fiber at $s=0$ records the boundary behavior.  In a concrete course one may take $X$ to be a family of divisors, modules, curves, bundles, or representations, depending on the chapter.  The moduli-theoretic lesson is that the morphism $T\to M$ detects both the general member and the limiting point.
\end{example}
\begin{exercise}
Write the functor of points for the moduli problem on this sheet and identify its value on $\Spec k$ and on $\Spec k[\epsilon]/(\epsilon^2)$.
\end{exercise}
\begin{exercise}
Formulate a small example in which the naive set of isomorphism classes is not Hausdorff, not separated, or not functorial.  Then repair the example by adding stability, rigidification, or a stack structure.
\end{exercise}
\clearpage
\section{Sheet 47: Geometric invariant theory: Hilbert--Mumford criterion}
This sheet studies \emph{Hilbert--Mumford criterion} as it appears in the broader chapter on \emph{geometric invariant theory}.  The recurring question is whether a proposed parameter space remembers enough information about families while remaining geometric enough to compute with.
\begin{definition}[Hilbert--mumford criterion viewpoint]
A Hilbert--Mumford criterion moduli functor assigns to each test scheme $T$ the set, groupoid, or category of families over $T$ satisfying the chosen condition.  Two families are identified only when the moduli problem says that an isomorphism, not merely a bijection on fibers, has been supplied.
\end{definition}
\begin{proposition}[Hilbert--mumford criterion viewpoint]
If a moduli functor for Hilbert--Mumford criterion is represented by a scheme $M$, then natural transformations from the functor to any functor $G$ are the same as elements of $G(M)$.  In practice this converts universal constructions on families into morphisms out of $M$.
\end{proposition}
\begin{example}[Hilbert--mumford criterion viewpoint]
For Hilbert--Mumford criterion, take a base $T=\Spec R$ and replace a single object by an $R$-flat family.  Flatness prevents the numerical invariants from jumping on first-order thickenings, so it is the minimal condition under which the fibers can reasonably be compared.
\end{example}
\begin{theorem}[Hilbert--mumford criterion viewpoint]
Suppose the objects in the Hilbert--Mumford criterion problem have no non-trivial automorphisms and admit effective descent for the chosen topology.  Then a universal family, if it exists, is unique up to unique isomorphism.  Consequently the representing object is unique up to unique isomorphism.
\end{theorem}
\begin{remark}[Hilbert--mumford criterion viewpoint]
The phrase ``the moduli space of Hilbert--Mumford criterion'' is therefore shorthand.  It may mean a fine moduli scheme, a coarse moduli space, an algebraic stack, a GIT quotient, or a compactification.  The correct meaning is determined by what information is retained about families and automorphisms.
\end{remark}
\begin{proof}[Proof sketch of the theorem/proposition on this sheet]
The argument is the standard Yoneda-and-descent argument specialized to Hilbert--Mumford criterion.  A family over a test scheme $T$ corresponds, under representability, to a morphism $T\to M$.  Pulling the universal family back along this morphism recovers the given family.  If two morphisms give isomorphic pullbacks, the functorial equality on all test schemes forces the morphisms to agree.  The hypotheses on automorphisms ensure that the equality is not weakened to an equivalence relation in a groupoid.
\end{proof}
\begin{example}[A test calculation for Hilbert--Mumford criterion]
Let $T=\Spec k[s]$ and let the invariant of interest be denoted $\nu$.  A toy family $X\to T$ is chosen so that $\nu(X_t)=\nu_0$ for $s\ne0$ and the special fiber at $s=0$ records the boundary behavior.  In a concrete course one may take $X$ to be a family of divisors, modules, curves, bundles, or representations, depending on the chapter.  The moduli-theoretic lesson is that the morphism $T\to M$ detects both the general member and the limiting point.
\end{example}
\begin{exercise}
Find one automorphism group occurring in the example and decide whether it obstructs the existence of a fine moduli scheme.
\end{exercise}
\begin{exercise}
Formulate a small example in which the naive set of isomorphism classes is not Hausdorff, not separated, or not functorial.  Then repair the example by adding stability, rigidification, or a stack structure.
\end{exercise}
\clearpage
\section{Sheet 48: Geometric invariant theory: semistable loci}
This sheet studies \emph{semistable loci} as it appears in the broader chapter on \emph{geometric invariant theory}.  The recurring question is whether a proposed parameter space remembers enough information about families while remaining geometric enough to compute with.
\begin{definition}[Semistable loci viewpoint]
A semistable loci moduli functor assigns to each test scheme $T$ the set, groupoid, or category of families over $T$ satisfying the chosen condition.  Two families are identified only when the moduli problem says that an isomorphism, not merely a bijection on fibers, has been supplied.
\end{definition}
\begin{proposition}[Semistable loci viewpoint]
If a moduli functor for semistable loci is represented by a scheme $M$, then natural transformations from the functor to any functor $G$ are the same as elements of $G(M)$.  In practice this converts universal constructions on families into morphisms out of $M$.
\end{proposition}
\begin{example}[Semistable loci viewpoint]
For semistable loci, take a base $T=\Spec R$ and replace a single object by an $R$-flat family.  Flatness prevents the numerical invariants from jumping on first-order thickenings, so it is the minimal condition under which the fibers can reasonably be compared.
\end{example}
\begin{theorem}[Semistable loci viewpoint]
Suppose the objects in the semistable loci problem have no non-trivial automorphisms and admit effective descent for the chosen topology.  Then a universal family, if it exists, is unique up to unique isomorphism.  Consequently the representing object is unique up to unique isomorphism.
\end{theorem}
\begin{remark}[Semistable loci viewpoint]
The phrase ``the moduli space of semistable loci'' is therefore shorthand.  It may mean a fine moduli scheme, a coarse moduli space, an algebraic stack, a GIT quotient, or a compactification.  The correct meaning is determined by what information is retained about families and automorphisms.
\end{remark}
\begin{proof}[Proof sketch of the theorem/proposition on this sheet]
The argument is the standard Yoneda-and-descent argument specialized to semistable loci.  A family over a test scheme $T$ corresponds, under representability, to a morphism $T\to M$.  Pulling the universal family back along this morphism recovers the given family.  If two morphisms give isomorphic pullbacks, the functorial equality on all test schemes forces the morphisms to agree.  The hypotheses on automorphisms ensure that the equality is not weakened to an equivalence relation in a groupoid.
\end{proof}
\begin{example}[A test calculation for semistable loci]
Let $T=\Spec k[s]$ and let the invariant of interest be denoted $\nu$.  A toy family $X\to T$ is chosen so that $\nu(X_t)=\nu_0$ for $s\ne0$ and the special fiber at $s=0$ records the boundary behavior.  In a concrete course one may take $X$ to be a family of divisors, modules, curves, bundles, or representations, depending on the chapter.  The moduli-theoretic lesson is that the morphism $T\to M$ detects both the general member and the limiting point.
\end{example}
\begin{exercise}
Check a base-change square explicitly: pull the displayed family back along $T\prime\to T$ and verify which invariant is preserved.
\end{exercise}
\begin{exercise}
Formulate a small example in which the naive set of isomorphism classes is not Hausdorff, not separated, or not functorial.  Then repair the example by adding stability, rigidification, or a stack structure.
\end{exercise}
\clearpage
\section{Sheet 49: Geometric invariant theory: GIT quotients}
This sheet studies \emph{GIT quotients} as it appears in the broader chapter on \emph{geometric invariant theory}.  The recurring question is whether a proposed parameter space remembers enough information about families while remaining geometric enough to compute with.
\begin{definition}[Git quotients viewpoint]
A GIT quotients moduli functor assigns to each test scheme $T$ the set, groupoid, or category of families over $T$ satisfying the chosen condition.  Two families are identified only when the moduli problem says that an isomorphism, not merely a bijection on fibers, has been supplied.
\end{definition}
\begin{proposition}[Git quotients viewpoint]
If a moduli functor for GIT quotients is represented by a scheme $M$, then natural transformations from the functor to any functor $G$ are the same as elements of $G(M)$.  In practice this converts universal constructions on families into morphisms out of $M$.
\end{proposition}
\begin{example}[Git quotients viewpoint]
For GIT quotients, take a base $T=\Spec R$ and replace a single object by an $R$-flat family.  Flatness prevents the numerical invariants from jumping on first-order thickenings, so it is the minimal condition under which the fibers can reasonably be compared.
\end{example}
\begin{theorem}[Git quotients viewpoint]
Suppose the objects in the GIT quotients problem have no non-trivial automorphisms and admit effective descent for the chosen topology.  Then a universal family, if it exists, is unique up to unique isomorphism.  Consequently the representing object is unique up to unique isomorphism.
\end{theorem}
\begin{remark}[Git quotients viewpoint]
The phrase ``the moduli space of GIT quotients'' is therefore shorthand.  It may mean a fine moduli scheme, a coarse moduli space, an algebraic stack, a GIT quotient, or a compactification.  The correct meaning is determined by what information is retained about families and automorphisms.
\end{remark}
\begin{proof}[Proof sketch of the theorem/proposition on this sheet]
The argument is the standard Yoneda-and-descent argument specialized to GIT quotients.  A family over a test scheme $T$ corresponds, under representability, to a morphism $T\to M$.  Pulling the universal family back along this morphism recovers the given family.  If two morphisms give isomorphic pullbacks, the functorial equality on all test schemes forces the morphisms to agree.  The hypotheses on automorphisms ensure that the equality is not weakened to an equivalence relation in a groupoid.
\end{proof}
\begin{example}[A test calculation for GIT quotients]
Let $T=\Spec k[s]$ and let the invariant of interest be denoted $\nu$.  A toy family $X\to T$ is chosen so that $\nu(X_t)=\nu_0$ for $s\ne0$ and the special fiber at $s=0$ records the boundary behavior.  In a concrete course one may take $X$ to be a family of divisors, modules, curves, bundles, or representations, depending on the chapter.  The moduli-theoretic lesson is that the morphism $T\to M$ detects both the general member and the limiting point.
\end{example}
\begin{exercise}
Construct a degeneration over $\Spec k[[t]]$ whose generic fiber satisfies the condition but whose special fiber forces a compactification.
\end{exercise}
\begin{exercise}
Formulate a small example in which the naive set of isomorphism classes is not Hausdorff, not separated, or not functorial.  Then repair the example by adding stability, rigidification, or a stack structure.
\end{exercise}
\clearpage
\section{Sheet 50: Geometric invariant theory: variation of GIT}
This sheet studies \emph{variation of GIT} as it appears in the broader chapter on \emph{geometric invariant theory}.  The recurring question is whether a proposed parameter space remembers enough information about families while remaining geometric enough to compute with.
\begin{definition}[Variation of git viewpoint]
A variation of GIT moduli functor assigns to each test scheme $T$ the set, groupoid, or category of families over $T$ satisfying the chosen condition.  Two families are identified only when the moduli problem says that an isomorphism, not merely a bijection on fibers, has been supplied.
\end{definition}
\begin{proposition}[Variation of git viewpoint]
If a moduli functor for variation of GIT is represented by a scheme $M$, then natural transformations from the functor to any functor $G$ are the same as elements of $G(M)$.  In practice this converts universal constructions on families into morphisms out of $M$.
\end{proposition}
\begin{example}[Variation of git viewpoint]
For variation of GIT, take a base $T=\Spec R$ and replace a single object by an $R$-flat family.  Flatness prevents the numerical invariants from jumping on first-order thickenings, so it is the minimal condition under which the fibers can reasonably be compared.
\end{example}
\begin{theorem}[Variation of git viewpoint]
Suppose the objects in the variation of GIT problem have no non-trivial automorphisms and admit effective descent for the chosen topology.  Then a universal family, if it exists, is unique up to unique isomorphism.  Consequently the representing object is unique up to unique isomorphism.
\end{theorem}
\begin{remark}[Variation of git viewpoint]
The phrase ``the moduli space of variation of GIT'' is therefore shorthand.  It may mean a fine moduli scheme, a coarse moduli space, an algebraic stack, a GIT quotient, or a compactification.  The correct meaning is determined by what information is retained about families and automorphisms.
\end{remark}
\begin{proof}[Proof sketch of the theorem/proposition on this sheet]
The argument is the standard Yoneda-and-descent argument specialized to variation of GIT.  A family over a test scheme $T$ corresponds, under representability, to a morphism $T\to M$.  Pulling the universal family back along this morphism recovers the given family.  If two morphisms give isomorphic pullbacks, the functorial equality on all test schemes forces the morphisms to agree.  The hypotheses on automorphisms ensure that the equality is not weakened to an equivalence relation in a groupoid.
\end{proof}
\begin{example}[A test calculation for variation of GIT]
Let $T=\Spec k[s]$ and let the invariant of interest be denoted $\nu$.  A toy family $X\to T$ is chosen so that $\nu(X_t)=\nu_0$ for $s\ne0$ and the special fiber at $s=0$ records the boundary behavior.  In a concrete course one may take $X$ to be a family of divisors, modules, curves, bundles, or representations, depending on the chapter.  The moduli-theoretic lesson is that the morphism $T\to M$ detects both the general member and the limiting point.
\end{example}
\begin{exercise}
Compare the functor with the corresponding groupoid-valued functor.  Which information is lost after passing to isomorphism classes?
\end{exercise}
\begin{exercise}
Formulate a small example in which the naive set of isomorphism classes is not Hausdorff, not separated, or not functorial.  Then repair the example by adding stability, rigidification, or a stack structure.
\end{exercise}
\clearpage
\section{Sheet 51: Moduli of sheaves: flat families}
This sheet studies \emph{flat families} as it appears in the broader chapter on \emph{moduli of sheaves}.  The recurring question is whether a proposed parameter space remembers enough information about families while remaining geometric enough to compute with.
\begin{definition}[Flat families viewpoint]
A flat families moduli functor assigns to each test scheme $T$ the set, groupoid, or category of families over $T$ satisfying the chosen condition.  Two families are identified only when the moduli problem says that an isomorphism, not merely a bijection on fibers, has been supplied.
\end{definition}
\begin{proposition}[Flat families viewpoint]
If a moduli functor for flat families is represented by a scheme $M$, then natural transformations from the functor to any functor $G$ are the same as elements of $G(M)$.  In practice this converts universal constructions on families into morphisms out of $M$.
\end{proposition}
\begin{example}[Flat families viewpoint]
For flat families, take a base $T=\Spec R$ and replace a single object by an $R$-flat family.  Flatness prevents the numerical invariants from jumping on first-order thickenings, so it is the minimal condition under which the fibers can reasonably be compared.
\end{example}
\begin{theorem}[Flat families viewpoint]
Suppose the objects in the flat families problem have no non-trivial automorphisms and admit effective descent for the chosen topology.  Then a universal family, if it exists, is unique up to unique isomorphism.  Consequently the representing object is unique up to unique isomorphism.
\end{theorem}
\begin{remark}[Flat families viewpoint]
The phrase ``the moduli space of flat families'' is therefore shorthand.  It may mean a fine moduli scheme, a coarse moduli space, an algebraic stack, a GIT quotient, or a compactification.  The correct meaning is determined by what information is retained about families and automorphisms.
\end{remark}
\begin{proof}[Proof sketch of the theorem/proposition on this sheet]
The argument is the standard Yoneda-and-descent argument specialized to flat families.  A family over a test scheme $T$ corresponds, under representability, to a morphism $T\to M$.  Pulling the universal family back along this morphism recovers the given family.  If two morphisms give isomorphic pullbacks, the functorial equality on all test schemes forces the morphisms to agree.  The hypotheses on automorphisms ensure that the equality is not weakened to an equivalence relation in a groupoid.
\end{proof}
\begin{example}[A test calculation for flat families]
Let $T=\Spec k[s]$ and let the invariant of interest be denoted $\nu$.  A toy family $X\to T$ is chosen so that $\nu(X_t)=\nu_0$ for $s\ne0$ and the special fiber at $s=0$ records the boundary behavior.  In a concrete course one may take $X$ to be a family of divisors, modules, curves, bundles, or representations, depending on the chapter.  The moduli-theoretic lesson is that the morphism $T\to M$ detects both the general member and the limiting point.
\end{example}
\begin{exercise}
Write the functor of points for the moduli problem on this sheet and identify its value on $\Spec k$ and on $\Spec k[\epsilon]/(\epsilon^2)$.
\end{exercise}
\begin{exercise}
Formulate a small example in which the naive set of isomorphism classes is not Hausdorff, not separated, or not functorial.  Then repair the example by adding stability, rigidification, or a stack structure.
\end{exercise}
\clearpage
\section{Sheet 52: Moduli of sheaves: purity}
This sheet studies \emph{purity} as it appears in the broader chapter on \emph{moduli of sheaves}.  The recurring question is whether a proposed parameter space remembers enough information about families while remaining geometric enough to compute with.
\begin{definition}[Purity viewpoint]
A purity moduli functor assigns to each test scheme $T$ the set, groupoid, or category of families over $T$ satisfying the chosen condition.  Two families are identified only when the moduli problem says that an isomorphism, not merely a bijection on fibers, has been supplied.
\end{definition}
\begin{proposition}[Purity viewpoint]
If a moduli functor for purity is represented by a scheme $M$, then natural transformations from the functor to any functor $G$ are the same as elements of $G(M)$.  In practice this converts universal constructions on families into morphisms out of $M$.
\end{proposition}
\begin{example}[Purity viewpoint]
For purity, take a base $T=\Spec R$ and replace a single object by an $R$-flat family.  Flatness prevents the numerical invariants from jumping on first-order thickenings, so it is the minimal condition under which the fibers can reasonably be compared.
\end{example}
\begin{theorem}[Purity viewpoint]
Suppose the objects in the purity problem have no non-trivial automorphisms and admit effective descent for the chosen topology.  Then a universal family, if it exists, is unique up to unique isomorphism.  Consequently the representing object is unique up to unique isomorphism.
\end{theorem}
\begin{remark}[Purity viewpoint]
The phrase ``the moduli space of purity'' is therefore shorthand.  It may mean a fine moduli scheme, a coarse moduli space, an algebraic stack, a GIT quotient, or a compactification.  The correct meaning is determined by what information is retained about families and automorphisms.
\end{remark}
\begin{proof}[Proof sketch of the theorem/proposition on this sheet]
The argument is the standard Yoneda-and-descent argument specialized to purity.  A family over a test scheme $T$ corresponds, under representability, to a morphism $T\to M$.  Pulling the universal family back along this morphism recovers the given family.  If two morphisms give isomorphic pullbacks, the functorial equality on all test schemes forces the morphisms to agree.  The hypotheses on automorphisms ensure that the equality is not weakened to an equivalence relation in a groupoid.
\end{proof}
\begin{example}[A test calculation for purity]
Let $T=\Spec k[s]$ and let the invariant of interest be denoted $\nu$.  A toy family $X\to T$ is chosen so that $\nu(X_t)=\nu_0$ for $s\ne0$ and the special fiber at $s=0$ records the boundary behavior.  In a concrete course one may take $X$ to be a family of divisors, modules, curves, bundles, or representations, depending on the chapter.  The moduli-theoretic lesson is that the morphism $T\to M$ detects both the general member and the limiting point.
\end{example}
\begin{exercise}
Find one automorphism group occurring in the example and decide whether it obstructs the existence of a fine moduli scheme.
\end{exercise}
\begin{exercise}
Formulate a small example in which the naive set of isomorphism classes is not Hausdorff, not separated, or not functorial.  Then repair the example by adding stability, rigidification, or a stack structure.
\end{exercise}
\clearpage
\section{Sheet 53: Moduli of sheaves: Gieseker stability}
This sheet studies \emph{Gieseker stability} as it appears in the broader chapter on \emph{moduli of sheaves}.  The recurring question is whether a proposed parameter space remembers enough information about families while remaining geometric enough to compute with.
\begin{definition}[Gieseker stability viewpoint]
A Gieseker stability moduli functor assigns to each test scheme $T$ the set, groupoid, or category of families over $T$ satisfying the chosen condition.  Two families are identified only when the moduli problem says that an isomorphism, not merely a bijection on fibers, has been supplied.
\end{definition}
\begin{proposition}[Gieseker stability viewpoint]
If a moduli functor for Gieseker stability is represented by a scheme $M$, then natural transformations from the functor to any functor $G$ are the same as elements of $G(M)$.  In practice this converts universal constructions on families into morphisms out of $M$.
\end{proposition}
\begin{example}[Gieseker stability viewpoint]
For Gieseker stability, take a base $T=\Spec R$ and replace a single object by an $R$-flat family.  Flatness prevents the numerical invariants from jumping on first-order thickenings, so it is the minimal condition under which the fibers can reasonably be compared.
\end{example}
\begin{theorem}[Gieseker stability viewpoint]
Suppose the objects in the Gieseker stability problem have no non-trivial automorphisms and admit effective descent for the chosen topology.  Then a universal family, if it exists, is unique up to unique isomorphism.  Consequently the representing object is unique up to unique isomorphism.
\end{theorem}
\begin{remark}[Gieseker stability viewpoint]
The phrase ``the moduli space of Gieseker stability'' is therefore shorthand.  It may mean a fine moduli scheme, a coarse moduli space, an algebraic stack, a GIT quotient, or a compactification.  The correct meaning is determined by what information is retained about families and automorphisms.
\end{remark}
\begin{proof}[Proof sketch of the theorem/proposition on this sheet]
The argument is the standard Yoneda-and-descent argument specialized to Gieseker stability.  A family over a test scheme $T$ corresponds, under representability, to a morphism $T\to M$.  Pulling the universal family back along this morphism recovers the given family.  If two morphisms give isomorphic pullbacks, the functorial equality on all test schemes forces the morphisms to agree.  The hypotheses on automorphisms ensure that the equality is not weakened to an equivalence relation in a groupoid.
\end{proof}
\begin{example}[A test calculation for Gieseker stability]
Let $T=\Spec k[s]$ and let the invariant of interest be denoted $\nu$.  A toy family $X\to T$ is chosen so that $\nu(X_t)=\nu_0$ for $s\ne0$ and the special fiber at $s=0$ records the boundary behavior.  In a concrete course one may take $X$ to be a family of divisors, modules, curves, bundles, or representations, depending on the chapter.  The moduli-theoretic lesson is that the morphism $T\to M$ detects both the general member and the limiting point.
\end{example}
\begin{exercise}
Check a base-change square explicitly: pull the displayed family back along $T\prime\to T$ and verify which invariant is preserved.
\end{exercise}
\begin{exercise}
Formulate a small example in which the naive set of isomorphism classes is not Hausdorff, not separated, or not functorial.  Then repair the example by adding stability, rigidification, or a stack structure.
\end{exercise}
\clearpage
\section{Sheet 54: Moduli of sheaves: Quot schemes}
This sheet studies \emph{Quot schemes} as it appears in the broader chapter on \emph{moduli of sheaves}.  The recurring question is whether a proposed parameter space remembers enough information about families while remaining geometric enough to compute with.
\begin{definition}[Quot schemes viewpoint]
A Quot schemes moduli functor assigns to each test scheme $T$ the set, groupoid, or category of families over $T$ satisfying the chosen condition.  Two families are identified only when the moduli problem says that an isomorphism, not merely a bijection on fibers, has been supplied.
\end{definition}
\begin{proposition}[Quot schemes viewpoint]
If a moduli functor for Quot schemes is represented by a scheme $M$, then natural transformations from the functor to any functor $G$ are the same as elements of $G(M)$.  In practice this converts universal constructions on families into morphisms out of $M$.
\end{proposition}
\begin{example}[Quot schemes viewpoint]
For Quot schemes, take a base $T=\Spec R$ and replace a single object by an $R$-flat family.  Flatness prevents the numerical invariants from jumping on first-order thickenings, so it is the minimal condition under which the fibers can reasonably be compared.
\end{example}
\begin{theorem}[Quot schemes viewpoint]
Suppose the objects in the Quot schemes problem have no non-trivial automorphisms and admit effective descent for the chosen topology.  Then a universal family, if it exists, is unique up to unique isomorphism.  Consequently the representing object is unique up to unique isomorphism.
\end{theorem}
\begin{remark}[Quot schemes viewpoint]
The phrase ``the moduli space of Quot schemes'' is therefore shorthand.  It may mean a fine moduli scheme, a coarse moduli space, an algebraic stack, a GIT quotient, or a compactification.  The correct meaning is determined by what information is retained about families and automorphisms.
\end{remark}
\begin{proof}[Proof sketch of the theorem/proposition on this sheet]
The argument is the standard Yoneda-and-descent argument specialized to Quot schemes.  A family over a test scheme $T$ corresponds, under representability, to a morphism $T\to M$.  Pulling the universal family back along this morphism recovers the given family.  If two morphisms give isomorphic pullbacks, the functorial equality on all test schemes forces the morphisms to agree.  The hypotheses on automorphisms ensure that the equality is not weakened to an equivalence relation in a groupoid.
\end{proof}
\begin{example}[A test calculation for Quot schemes]
Let $T=\Spec k[s]$ and let the invariant of interest be denoted $\nu$.  A toy family $X\to T$ is chosen so that $\nu(X_t)=\nu_0$ for $s\ne0$ and the special fiber at $s=0$ records the boundary behavior.  In a concrete course one may take $X$ to be a family of divisors, modules, curves, bundles, or representations, depending on the chapter.  The moduli-theoretic lesson is that the morphism $T\to M$ detects both the general member and the limiting point.
\end{example}
\begin{exercise}
Construct a degeneration over $\Spec k[[t]]$ whose generic fiber satisfies the condition but whose special fiber forces a compactification.
\end{exercise}
\begin{exercise}
Formulate a small example in which the naive set of isomorphism classes is not Hausdorff, not separated, or not functorial.  Then repair the example by adding stability, rigidification, or a stack structure.
\end{exercise}
\clearpage
\section{Sheet 55: Moduli of sheaves: determinant line bundles}
This sheet studies \emph{determinant line bundles} as it appears in the broader chapter on \emph{moduli of sheaves}.  The recurring question is whether a proposed parameter space remembers enough information about families while remaining geometric enough to compute with.
\begin{definition}[Determinant line bundles viewpoint]
A determinant line bundles moduli functor assigns to each test scheme $T$ the set, groupoid, or category of families over $T$ satisfying the chosen condition.  Two families are identified only when the moduli problem says that an isomorphism, not merely a bijection on fibers, has been supplied.
\end{definition}
\begin{proposition}[Determinant line bundles viewpoint]
If a moduli functor for determinant line bundles is represented by a scheme $M$, then natural transformations from the functor to any functor $G$ are the same as elements of $G(M)$.  In practice this converts universal constructions on families into morphisms out of $M$.
\end{proposition}
\begin{example}[Determinant line bundles viewpoint]
For determinant line bundles, take a base $T=\Spec R$ and replace a single object by an $R$-flat family.  Flatness prevents the numerical invariants from jumping on first-order thickenings, so it is the minimal condition under which the fibers can reasonably be compared.
\end{example}
\begin{theorem}[Determinant line bundles viewpoint]
Suppose the objects in the determinant line bundles problem have no non-trivial automorphisms and admit effective descent for the chosen topology.  Then a universal family, if it exists, is unique up to unique isomorphism.  Consequently the representing object is unique up to unique isomorphism.
\end{theorem}
\begin{remark}[Determinant line bundles viewpoint]
The phrase ``the moduli space of determinant line bundles'' is therefore shorthand.  It may mean a fine moduli scheme, a coarse moduli space, an algebraic stack, a GIT quotient, or a compactification.  The correct meaning is determined by what information is retained about families and automorphisms.
\end{remark}
\begin{proof}[Proof sketch of the theorem/proposition on this sheet]
The argument is the standard Yoneda-and-descent argument specialized to determinant line bundles.  A family over a test scheme $T$ corresponds, under representability, to a morphism $T\to M$.  Pulling the universal family back along this morphism recovers the given family.  If two morphisms give isomorphic pullbacks, the functorial equality on all test schemes forces the morphisms to agree.  The hypotheses on automorphisms ensure that the equality is not weakened to an equivalence relation in a groupoid.
\end{proof}
\begin{example}[A test calculation for determinant line bundles]
Let $T=\Spec k[s]$ and let the invariant of interest be denoted $\nu$.  A toy family $X\to T$ is chosen so that $\nu(X_t)=\nu_0$ for $s\ne0$ and the special fiber at $s=0$ records the boundary behavior.  In a concrete course one may take $X$ to be a family of divisors, modules, curves, bundles, or representations, depending on the chapter.  The moduli-theoretic lesson is that the morphism $T\to M$ detects both the general member and the limiting point.
\end{example}
\begin{exercise}
Compare the functor with the corresponding groupoid-valued functor.  Which information is lost after passing to isomorphism classes?
\end{exercise}
\begin{exercise}
Formulate a small example in which the naive set of isomorphism classes is not Hausdorff, not separated, or not functorial.  Then repair the example by adding stability, rigidification, or a stack structure.
\end{exercise}
\clearpage
\section{Sheet 56: Derived and obstruction-theoretic ideas: cotangent complexes}
This sheet studies \emph{cotangent complexes} as it appears in the broader chapter on \emph{derived and obstruction-theoretic ideas}.  The recurring question is whether a proposed parameter space remembers enough information about families while remaining geometric enough to compute with.
\begin{definition}[Cotangent complexes viewpoint]
A cotangent complexes moduli functor assigns to each test scheme $T$ the set, groupoid, or category of families over $T$ satisfying the chosen condition.  Two families are identified only when the moduli problem says that an isomorphism, not merely a bijection on fibers, has been supplied.
\end{definition}
\begin{proposition}[Cotangent complexes viewpoint]
If a moduli functor for cotangent complexes is represented by a scheme $M$, then natural transformations from the functor to any functor $G$ are the same as elements of $G(M)$.  In practice this converts universal constructions on families into morphisms out of $M$.
\end{proposition}
\begin{example}[Cotangent complexes viewpoint]
For cotangent complexes, take a base $T=\Spec R$ and replace a single object by an $R$-flat family.  Flatness prevents the numerical invariants from jumping on first-order thickenings, so it is the minimal condition under which the fibers can reasonably be compared.
\end{example}
\begin{theorem}[Cotangent complexes viewpoint]
Suppose the objects in the cotangent complexes problem have no non-trivial automorphisms and admit effective descent for the chosen topology.  Then a universal family, if it exists, is unique up to unique isomorphism.  Consequently the representing object is unique up to unique isomorphism.
\end{theorem}
\begin{remark}[Cotangent complexes viewpoint]
The phrase ``the moduli space of cotangent complexes'' is therefore shorthand.  It may mean a fine moduli scheme, a coarse moduli space, an algebraic stack, a GIT quotient, or a compactification.  The correct meaning is determined by what information is retained about families and automorphisms.
\end{remark}
\begin{proof}[Proof sketch of the theorem/proposition on this sheet]
The argument is the standard Yoneda-and-descent argument specialized to cotangent complexes.  A family over a test scheme $T$ corresponds, under representability, to a morphism $T\to M$.  Pulling the universal family back along this morphism recovers the given family.  If two morphisms give isomorphic pullbacks, the functorial equality on all test schemes forces the morphisms to agree.  The hypotheses on automorphisms ensure that the equality is not weakened to an equivalence relation in a groupoid.
\end{proof}
\begin{example}[A test calculation for cotangent complexes]
Let $T=\Spec k[s]$ and let the invariant of interest be denoted $\nu$.  A toy family $X\to T$ is chosen so that $\nu(X_t)=\nu_0$ for $s\ne0$ and the special fiber at $s=0$ records the boundary behavior.  In a concrete course one may take $X$ to be a family of divisors, modules, curves, bundles, or representations, depending on the chapter.  The moduli-theoretic lesson is that the morphism $T\to M$ detects both the general member and the limiting point.
\end{example}
\begin{exercise}
Write the functor of points for the moduli problem on this sheet and identify its value on $\Spec k$ and on $\Spec k[\epsilon]/(\epsilon^2)$.
\end{exercise}
\begin{exercise}
Formulate a small example in which the naive set of isomorphism classes is not Hausdorff, not separated, or not functorial.  Then repair the example by adding stability, rigidification, or a stack structure.
\end{exercise}
\clearpage
\section{Sheet 57: Derived and obstruction-theoretic ideas: perfect obstruction theories}
This sheet studies \emph{perfect obstruction theories} as it appears in the broader chapter on \emph{derived and obstruction-theoretic ideas}.  The recurring question is whether a proposed parameter space remembers enough information about families while remaining geometric enough to compute with.
\begin{definition}[Perfect obstruction theories viewpoint]
A perfect obstruction theories moduli functor assigns to each test scheme $T$ the set, groupoid, or category of families over $T$ satisfying the chosen condition.  Two families are identified only when the moduli problem says that an isomorphism, not merely a bijection on fibers, has been supplied.
\end{definition}
\begin{proposition}[Perfect obstruction theories viewpoint]
If a moduli functor for perfect obstruction theories is represented by a scheme $M$, then natural transformations from the functor to any functor $G$ are the same as elements of $G(M)$.  In practice this converts universal constructions on families into morphisms out of $M$.
\end{proposition}
\begin{example}[Perfect obstruction theories viewpoint]
For perfect obstruction theories, take a base $T=\Spec R$ and replace a single object by an $R$-flat family.  Flatness prevents the numerical invariants from jumping on first-order thickenings, so it is the minimal condition under which the fibers can reasonably be compared.
\end{example}
\begin{theorem}[Perfect obstruction theories viewpoint]
Suppose the objects in the perfect obstruction theories problem have no non-trivial automorphisms and admit effective descent for the chosen topology.  Then a universal family, if it exists, is unique up to unique isomorphism.  Consequently the representing object is unique up to unique isomorphism.
\end{theorem}
\begin{remark}[Perfect obstruction theories viewpoint]
The phrase ``the moduli space of perfect obstruction theories'' is therefore shorthand.  It may mean a fine moduli scheme, a coarse moduli space, an algebraic stack, a GIT quotient, or a compactification.  The correct meaning is determined by what information is retained about families and automorphisms.
\end{remark}
\begin{proof}[Proof sketch of the theorem/proposition on this sheet]
The argument is the standard Yoneda-and-descent argument specialized to perfect obstruction theories.  A family over a test scheme $T$ corresponds, under representability, to a morphism $T\to M$.  Pulling the universal family back along this morphism recovers the given family.  If two morphisms give isomorphic pullbacks, the functorial equality on all test schemes forces the morphisms to agree.  The hypotheses on automorphisms ensure that the equality is not weakened to an equivalence relation in a groupoid.
\end{proof}
\begin{example}[A test calculation for perfect obstruction theories]
Let $T=\Spec k[s]$ and let the invariant of interest be denoted $\nu$.  A toy family $X\to T$ is chosen so that $\nu(X_t)=\nu_0$ for $s\ne0$ and the special fiber at $s=0$ records the boundary behavior.  In a concrete course one may take $X$ to be a family of divisors, modules, curves, bundles, or representations, depending on the chapter.  The moduli-theoretic lesson is that the morphism $T\to M$ detects both the general member and the limiting point.
\end{example}
\begin{exercise}
Find one automorphism group occurring in the example and decide whether it obstructs the existence of a fine moduli scheme.
\end{exercise}
\begin{exercise}
Formulate a small example in which the naive set of isomorphism classes is not Hausdorff, not separated, or not functorial.  Then repair the example by adding stability, rigidification, or a stack structure.
\end{exercise}
\clearpage
\section{Sheet 58: Derived and obstruction-theoretic ideas: derived intersections}
This sheet studies \emph{derived intersections} as it appears in the broader chapter on \emph{derived and obstruction-theoretic ideas}.  The recurring question is whether a proposed parameter space remembers enough information about families while remaining geometric enough to compute with.
\begin{definition}[Derived intersections viewpoint]
A derived intersections moduli functor assigns to each test scheme $T$ the set, groupoid, or category of families over $T$ satisfying the chosen condition.  Two families are identified only when the moduli problem says that an isomorphism, not merely a bijection on fibers, has been supplied.
\end{definition}
\begin{proposition}[Derived intersections viewpoint]
If a moduli functor for derived intersections is represented by a scheme $M$, then natural transformations from the functor to any functor $G$ are the same as elements of $G(M)$.  In practice this converts universal constructions on families into morphisms out of $M$.
\end{proposition}
\begin{example}[Derived intersections viewpoint]
For derived intersections, take a base $T=\Spec R$ and replace a single object by an $R$-flat family.  Flatness prevents the numerical invariants from jumping on first-order thickenings, so it is the minimal condition under which the fibers can reasonably be compared.
\end{example}
\begin{theorem}[Derived intersections viewpoint]
Suppose the objects in the derived intersections problem have no non-trivial automorphisms and admit effective descent for the chosen topology.  Then a universal family, if it exists, is unique up to unique isomorphism.  Consequently the representing object is unique up to unique isomorphism.
\end{theorem}
\begin{remark}[Derived intersections viewpoint]
The phrase ``the moduli space of derived intersections'' is therefore shorthand.  It may mean a fine moduli scheme, a coarse moduli space, an algebraic stack, a GIT quotient, or a compactification.  The correct meaning is determined by what information is retained about families and automorphisms.
\end{remark}
\begin{proof}[Proof sketch of the theorem/proposition on this sheet]
The argument is the standard Yoneda-and-descent argument specialized to derived intersections.  A family over a test scheme $T$ corresponds, under representability, to a morphism $T\to M$.  Pulling the universal family back along this morphism recovers the given family.  If two morphisms give isomorphic pullbacks, the functorial equality on all test schemes forces the morphisms to agree.  The hypotheses on automorphisms ensure that the equality is not weakened to an equivalence relation in a groupoid.
\end{proof}
\begin{example}[A test calculation for derived intersections]
Let $T=\Spec k[s]$ and let the invariant of interest be denoted $\nu$.  A toy family $X\to T$ is chosen so that $\nu(X_t)=\nu_0$ for $s\ne0$ and the special fiber at $s=0$ records the boundary behavior.  In a concrete course one may take $X$ to be a family of divisors, modules, curves, bundles, or representations, depending on the chapter.  The moduli-theoretic lesson is that the morphism $T\to M$ detects both the general member and the limiting point.
\end{example}
\begin{exercise}
Check a base-change square explicitly: pull the displayed family back along $T\prime\to T$ and verify which invariant is preserved.
\end{exercise}
\begin{exercise}
Formulate a small example in which the naive set of isomorphism classes is not Hausdorff, not separated, or not functorial.  Then repair the example by adding stability, rigidification, or a stack structure.
\end{exercise}
\clearpage
\section{Sheet 59: Derived and obstruction-theoretic ideas: virtual classes}
This sheet studies \emph{virtual classes} as it appears in the broader chapter on \emph{derived and obstruction-theoretic ideas}.  The recurring question is whether a proposed parameter space remembers enough information about families while remaining geometric enough to compute with.
\begin{definition}[Virtual classes viewpoint]
A virtual classes moduli functor assigns to each test scheme $T$ the set, groupoid, or category of families over $T$ satisfying the chosen condition.  Two families are identified only when the moduli problem says that an isomorphism, not merely a bijection on fibers, has been supplied.
\end{definition}
\begin{proposition}[Virtual classes viewpoint]
If a moduli functor for virtual classes is represented by a scheme $M$, then natural transformations from the functor to any functor $G$ are the same as elements of $G(M)$.  In practice this converts universal constructions on families into morphisms out of $M$.
\end{proposition}
\begin{example}[Virtual classes viewpoint]
For virtual classes, take a base $T=\Spec R$ and replace a single object by an $R$-flat family.  Flatness prevents the numerical invariants from jumping on first-order thickenings, so it is the minimal condition under which the fibers can reasonably be compared.
\end{example}
\begin{theorem}[Virtual classes viewpoint]
Suppose the objects in the virtual classes problem have no non-trivial automorphisms and admit effective descent for the chosen topology.  Then a universal family, if it exists, is unique up to unique isomorphism.  Consequently the representing object is unique up to unique isomorphism.
\end{theorem}
\begin{remark}[Virtual classes viewpoint]
The phrase ``the moduli space of virtual classes'' is therefore shorthand.  It may mean a fine moduli scheme, a coarse moduli space, an algebraic stack, a GIT quotient, or a compactification.  The correct meaning is determined by what information is retained about families and automorphisms.
\end{remark}
\begin{proof}[Proof sketch of the theorem/proposition on this sheet]
The argument is the standard Yoneda-and-descent argument specialized to virtual classes.  A family over a test scheme $T$ corresponds, under representability, to a morphism $T\to M$.  Pulling the universal family back along this morphism recovers the given family.  If two morphisms give isomorphic pullbacks, the functorial equality on all test schemes forces the morphisms to agree.  The hypotheses on automorphisms ensure that the equality is not weakened to an equivalence relation in a groupoid.
\end{proof}
\begin{example}[A test calculation for virtual classes]
Let $T=\Spec k[s]$ and let the invariant of interest be denoted $\nu$.  A toy family $X\to T$ is chosen so that $\nu(X_t)=\nu_0$ for $s\ne0$ and the special fiber at $s=0$ records the boundary behavior.  In a concrete course one may take $X$ to be a family of divisors, modules, curves, bundles, or representations, depending on the chapter.  The moduli-theoretic lesson is that the morphism $T\to M$ detects both the general member and the limiting point.
\end{example}
\begin{exercise}
Construct a degeneration over $\Spec k[[t]]$ whose generic fiber satisfies the condition but whose special fiber forces a compactification.
\end{exercise}
\begin{exercise}
Formulate a small example in which the naive set of isomorphism classes is not Hausdorff, not separated, or not functorial.  Then repair the example by adding stability, rigidification, or a stack structure.
\end{exercise}
\clearpage
\section{Sheet 60: Derived and obstruction-theoretic ideas: orientation data}
This sheet studies \emph{orientation data} as it appears in the broader chapter on \emph{derived and obstruction-theoretic ideas}.  The recurring question is whether a proposed parameter space remembers enough information about families while remaining geometric enough to compute with.
\begin{definition}[Orientation data viewpoint]
A orientation data moduli functor assigns to each test scheme $T$ the set, groupoid, or category of families over $T$ satisfying the chosen condition.  Two families are identified only when the moduli problem says that an isomorphism, not merely a bijection on fibers, has been supplied.
\end{definition}
\begin{proposition}[Orientation data viewpoint]
If a moduli functor for orientation data is represented by a scheme $M$, then natural transformations from the functor to any functor $G$ are the same as elements of $G(M)$.  In practice this converts universal constructions on families into morphisms out of $M$.
\end{proposition}
\begin{example}[Orientation data viewpoint]
For orientation data, take a base $T=\Spec R$ and replace a single object by an $R$-flat family.  Flatness prevents the numerical invariants from jumping on first-order thickenings, so it is the minimal condition under which the fibers can reasonably be compared.
\end{example}
\begin{theorem}[Orientation data viewpoint]
Suppose the objects in the orientation data problem have no non-trivial automorphisms and admit effective descent for the chosen topology.  Then a universal family, if it exists, is unique up to unique isomorphism.  Consequently the representing object is unique up to unique isomorphism.
\end{theorem}
\begin{remark}[Orientation data viewpoint]
The phrase ``the moduli space of orientation data'' is therefore shorthand.  It may mean a fine moduli scheme, a coarse moduli space, an algebraic stack, a GIT quotient, or a compactification.  The correct meaning is determined by what information is retained about families and automorphisms.
\end{remark}
\begin{proof}[Proof sketch of the theorem/proposition on this sheet]
The argument is the standard Yoneda-and-descent argument specialized to orientation data.  A family over a test scheme $T$ corresponds, under representability, to a morphism $T\to M$.  Pulling the universal family back along this morphism recovers the given family.  If two morphisms give isomorphic pullbacks, the functorial equality on all test schemes forces the morphisms to agree.  The hypotheses on automorphisms ensure that the equality is not weakened to an equivalence relation in a groupoid.
\end{proof}
\begin{example}[A test calculation for orientation data]
Let $T=\Spec k[s]$ and let the invariant of interest be denoted $\nu$.  A toy family $X\to T$ is chosen so that $\nu(X_t)=\nu_0$ for $s\ne0$ and the special fiber at $s=0$ records the boundary behavior.  In a concrete course one may take $X$ to be a family of divisors, modules, curves, bundles, or representations, depending on the chapter.  The moduli-theoretic lesson is that the morphism $T\to M$ detects both the general member and the limiting point.
\end{example}
\begin{exercise}
Compare the functor with the corresponding groupoid-valued functor.  Which information is lost after passing to isomorphism classes?
\end{exercise}
\begin{exercise}
Formulate a small example in which the naive set of isomorphism classes is not Hausdorff, not separated, or not functorial.  Then repair the example by adding stability, rigidification, or a stack structure.
\end{exercise}
\clearpage
\section{Sheet 61: Examples from low dimensions: binary forms}
This sheet studies \emph{binary forms} as it appears in the broader chapter on \emph{examples from low dimensions}.  The recurring question is whether a proposed parameter space remembers enough information about families while remaining geometric enough to compute with.
\begin{definition}[Binary forms viewpoint]
A binary forms moduli functor assigns to each test scheme $T$ the set, groupoid, or category of families over $T$ satisfying the chosen condition.  Two families are identified only when the moduli problem says that an isomorphism, not merely a bijection on fibers, has been supplied.
\end{definition}
\begin{proposition}[Binary forms viewpoint]
If a moduli functor for binary forms is represented by a scheme $M$, then natural transformations from the functor to any functor $G$ are the same as elements of $G(M)$.  In practice this converts universal constructions on families into morphisms out of $M$.
\end{proposition}
\begin{example}[Binary forms viewpoint]
For binary forms, take a base $T=\Spec R$ and replace a single object by an $R$-flat family.  Flatness prevents the numerical invariants from jumping on first-order thickenings, so it is the minimal condition under which the fibers can reasonably be compared.
\end{example}
\begin{theorem}[Binary forms viewpoint]
Suppose the objects in the binary forms problem have no non-trivial automorphisms and admit effective descent for the chosen topology.  Then a universal family, if it exists, is unique up to unique isomorphism.  Consequently the representing object is unique up to unique isomorphism.
\end{theorem}
\begin{remark}[Binary forms viewpoint]
The phrase ``the moduli space of binary forms'' is therefore shorthand.  It may mean a fine moduli scheme, a coarse moduli space, an algebraic stack, a GIT quotient, or a compactification.  The correct meaning is determined by what information is retained about families and automorphisms.
\end{remark}
\begin{proof}[Proof sketch of the theorem/proposition on this sheet]
The argument is the standard Yoneda-and-descent argument specialized to binary forms.  A family over a test scheme $T$ corresponds, under representability, to a morphism $T\to M$.  Pulling the universal family back along this morphism recovers the given family.  If two morphisms give isomorphic pullbacks, the functorial equality on all test schemes forces the morphisms to agree.  The hypotheses on automorphisms ensure that the equality is not weakened to an equivalence relation in a groupoid.
\end{proof}
\begin{example}[A test calculation for binary forms]
Let $T=\Spec k[s]$ and let the invariant of interest be denoted $\nu$.  A toy family $X\to T$ is chosen so that $\nu(X_t)=\nu_0$ for $s\ne0$ and the special fiber at $s=0$ records the boundary behavior.  In a concrete course one may take $X$ to be a family of divisors, modules, curves, bundles, or representations, depending on the chapter.  The moduli-theoretic lesson is that the morphism $T\to M$ detects both the general member and the limiting point.
\end{example}
\begin{exercise}
Write the functor of points for the moduli problem on this sheet and identify its value on $\Spec k$ and on $\Spec k[\epsilon]/(\epsilon^2)$.
\end{exercise}
\begin{exercise}
Formulate a small example in which the naive set of isomorphism classes is not Hausdorff, not separated, or not functorial.  Then repair the example by adding stability, rigidification, or a stack structure.
\end{exercise}
\clearpage
\section{Sheet 62: Examples from low dimensions: plane conics}
This sheet studies \emph{plane conics} as it appears in the broader chapter on \emph{examples from low dimensions}.  The recurring question is whether a proposed parameter space remembers enough information about families while remaining geometric enough to compute with.
\begin{definition}[Plane conics viewpoint]
A plane conics moduli functor assigns to each test scheme $T$ the set, groupoid, or category of families over $T$ satisfying the chosen condition.  Two families are identified only when the moduli problem says that an isomorphism, not merely a bijection on fibers, has been supplied.
\end{definition}
\begin{proposition}[Plane conics viewpoint]
If a moduli functor for plane conics is represented by a scheme $M$, then natural transformations from the functor to any functor $G$ are the same as elements of $G(M)$.  In practice this converts universal constructions on families into morphisms out of $M$.
\end{proposition}
\begin{example}[Plane conics viewpoint]
For plane conics, take a base $T=\Spec R$ and replace a single object by an $R$-flat family.  Flatness prevents the numerical invariants from jumping on first-order thickenings, so it is the minimal condition under which the fibers can reasonably be compared.
\end{example}
\begin{theorem}[Plane conics viewpoint]
Suppose the objects in the plane conics problem have no non-trivial automorphisms and admit effective descent for the chosen topology.  Then a universal family, if it exists, is unique up to unique isomorphism.  Consequently the representing object is unique up to unique isomorphism.
\end{theorem}
\begin{remark}[Plane conics viewpoint]
The phrase ``the moduli space of plane conics'' is therefore shorthand.  It may mean a fine moduli scheme, a coarse moduli space, an algebraic stack, a GIT quotient, or a compactification.  The correct meaning is determined by what information is retained about families and automorphisms.
\end{remark}
\begin{proof}[Proof sketch of the theorem/proposition on this sheet]
The argument is the standard Yoneda-and-descent argument specialized to plane conics.  A family over a test scheme $T$ corresponds, under representability, to a morphism $T\to M$.  Pulling the universal family back along this morphism recovers the given family.  If two morphisms give isomorphic pullbacks, the functorial equality on all test schemes forces the morphisms to agree.  The hypotheses on automorphisms ensure that the equality is not weakened to an equivalence relation in a groupoid.
\end{proof}
\begin{example}[A test calculation for plane conics]
Let $T=\Spec k[s]$ and let the invariant of interest be denoted $\nu$.  A toy family $X\to T$ is chosen so that $\nu(X_t)=\nu_0$ for $s\ne0$ and the special fiber at $s=0$ records the boundary behavior.  In a concrete course one may take $X$ to be a family of divisors, modules, curves, bundles, or representations, depending on the chapter.  The moduli-theoretic lesson is that the morphism $T\to M$ detects both the general member and the limiting point.
\end{example}
\begin{exercise}
Find one automorphism group occurring in the example and decide whether it obstructs the existence of a fine moduli scheme.
\end{exercise}
\begin{exercise}
Formulate a small example in which the naive set of isomorphism classes is not Hausdorff, not separated, or not functorial.  Then repair the example by adding stability, rigidification, or a stack structure.
\end{exercise}
\clearpage
\section{Sheet 63: Examples from low dimensions: elliptic curves}
This sheet studies \emph{elliptic curves} as it appears in the broader chapter on \emph{examples from low dimensions}.  The recurring question is whether a proposed parameter space remembers enough information about families while remaining geometric enough to compute with.
\begin{definition}[Elliptic curves viewpoint]
A elliptic curves moduli functor assigns to each test scheme $T$ the set, groupoid, or category of families over $T$ satisfying the chosen condition.  Two families are identified only when the moduli problem says that an isomorphism, not merely a bijection on fibers, has been supplied.
\end{definition}
\begin{proposition}[Elliptic curves viewpoint]
If a moduli functor for elliptic curves is represented by a scheme $M$, then natural transformations from the functor to any functor $G$ are the same as elements of $G(M)$.  In practice this converts universal constructions on families into morphisms out of $M$.
\end{proposition}
\begin{example}[Elliptic curves viewpoint]
For elliptic curves, take a base $T=\Spec R$ and replace a single object by an $R$-flat family.  Flatness prevents the numerical invariants from jumping on first-order thickenings, so it is the minimal condition under which the fibers can reasonably be compared.
\end{example}
\begin{theorem}[Elliptic curves viewpoint]
Suppose the objects in the elliptic curves problem have no non-trivial automorphisms and admit effective descent for the chosen topology.  Then a universal family, if it exists, is unique up to unique isomorphism.  Consequently the representing object is unique up to unique isomorphism.
\end{theorem}
\begin{remark}[Elliptic curves viewpoint]
The phrase ``the moduli space of elliptic curves'' is therefore shorthand.  It may mean a fine moduli scheme, a coarse moduli space, an algebraic stack, a GIT quotient, or a compactification.  The correct meaning is determined by what information is retained about families and automorphisms.
\end{remark}
\begin{proof}[Proof sketch of the theorem/proposition on this sheet]
The argument is the standard Yoneda-and-descent argument specialized to elliptic curves.  A family over a test scheme $T$ corresponds, under representability, to a morphism $T\to M$.  Pulling the universal family back along this morphism recovers the given family.  If two morphisms give isomorphic pullbacks, the functorial equality on all test schemes forces the morphisms to agree.  The hypotheses on automorphisms ensure that the equality is not weakened to an equivalence relation in a groupoid.
\end{proof}
\begin{example}[A test calculation for elliptic curves]
Let $T=\Spec k[s]$ and let the invariant of interest be denoted $\nu$.  A toy family $X\to T$ is chosen so that $\nu(X_t)=\nu_0$ for $s\ne0$ and the special fiber at $s=0$ records the boundary behavior.  In a concrete course one may take $X$ to be a family of divisors, modules, curves, bundles, or representations, depending on the chapter.  The moduli-theoretic lesson is that the morphism $T\to M$ detects both the general member and the limiting point.
\end{example}
\begin{exercise}
Check a base-change square explicitly: pull the displayed family back along $T\prime\to T$ and verify which invariant is preserved.
\end{exercise}
\begin{exercise}
Formulate a small example in which the naive set of isomorphism classes is not Hausdorff, not separated, or not functorial.  Then repair the example by adding stability, rigidification, or a stack structure.
\end{exercise}
\clearpage
\section{Sheet 64: Examples from low dimensions: K3 surfaces}
This sheet studies \emph{K3 surfaces} as it appears in the broader chapter on \emph{examples from low dimensions}.  The recurring question is whether a proposed parameter space remembers enough information about families while remaining geometric enough to compute with.
\begin{definition}[K3 surfaces viewpoint]
A K3 surfaces moduli functor assigns to each test scheme $T$ the set, groupoid, or category of families over $T$ satisfying the chosen condition.  Two families are identified only when the moduli problem says that an isomorphism, not merely a bijection on fibers, has been supplied.
\end{definition}
\begin{proposition}[K3 surfaces viewpoint]
If a moduli functor for K3 surfaces is represented by a scheme $M$, then natural transformations from the functor to any functor $G$ are the same as elements of $G(M)$.  In practice this converts universal constructions on families into morphisms out of $M$.
\end{proposition}
\begin{example}[K3 surfaces viewpoint]
For K3 surfaces, take a base $T=\Spec R$ and replace a single object by an $R$-flat family.  Flatness prevents the numerical invariants from jumping on first-order thickenings, so it is the minimal condition under which the fibers can reasonably be compared.
\end{example}
\begin{theorem}[K3 surfaces viewpoint]
Suppose the objects in the K3 surfaces problem have no non-trivial automorphisms and admit effective descent for the chosen topology.  Then a universal family, if it exists, is unique up to unique isomorphism.  Consequently the representing object is unique up to unique isomorphism.
\end{theorem}
\begin{remark}[K3 surfaces viewpoint]
The phrase ``the moduli space of K3 surfaces'' is therefore shorthand.  It may mean a fine moduli scheme, a coarse moduli space, an algebraic stack, a GIT quotient, or a compactification.  The correct meaning is determined by what information is retained about families and automorphisms.
\end{remark}
\begin{proof}[Proof sketch of the theorem/proposition on this sheet]
The argument is the standard Yoneda-and-descent argument specialized to K3 surfaces.  A family over a test scheme $T$ corresponds, under representability, to a morphism $T\to M$.  Pulling the universal family back along this morphism recovers the given family.  If two morphisms give isomorphic pullbacks, the functorial equality on all test schemes forces the morphisms to agree.  The hypotheses on automorphisms ensure that the equality is not weakened to an equivalence relation in a groupoid.
\end{proof}
\begin{example}[A test calculation for K3 surfaces]
Let $T=\Spec k[s]$ and let the invariant of interest be denoted $\nu$.  A toy family $X\to T$ is chosen so that $\nu(X_t)=\nu_0$ for $s\ne0$ and the special fiber at $s=0$ records the boundary behavior.  In a concrete course one may take $X$ to be a family of divisors, modules, curves, bundles, or representations, depending on the chapter.  The moduli-theoretic lesson is that the morphism $T\to M$ detects both the general member and the limiting point.
\end{example}
\begin{exercise}
Construct a degeneration over $\Spec k[[t]]$ whose generic fiber satisfies the condition but whose special fiber forces a compactification.
\end{exercise}
\begin{exercise}
Formulate a small example in which the naive set of isomorphism classes is not Hausdorff, not separated, or not functorial.  Then repair the example by adding stability, rigidification, or a stack structure.
\end{exercise}
\clearpage
\section{Sheet 65: Examples from low dimensions: quiver representations}
This sheet studies \emph{quiver representations} as it appears in the broader chapter on \emph{examples from low dimensions}.  The recurring question is whether a proposed parameter space remembers enough information about families while remaining geometric enough to compute with.
\begin{definition}[Quiver representations viewpoint]
A quiver representations moduli functor assigns to each test scheme $T$ the set, groupoid, or category of families over $T$ satisfying the chosen condition.  Two families are identified only when the moduli problem says that an isomorphism, not merely a bijection on fibers, has been supplied.
\end{definition}
\begin{proposition}[Quiver representations viewpoint]
If a moduli functor for quiver representations is represented by a scheme $M$, then natural transformations from the functor to any functor $G$ are the same as elements of $G(M)$.  In practice this converts universal constructions on families into morphisms out of $M$.
\end{proposition}
\begin{example}[Quiver representations viewpoint]
For quiver representations, take a base $T=\Spec R$ and replace a single object by an $R$-flat family.  Flatness prevents the numerical invariants from jumping on first-order thickenings, so it is the minimal condition under which the fibers can reasonably be compared.
\end{example}
\begin{theorem}[Quiver representations viewpoint]
Suppose the objects in the quiver representations problem have no non-trivial automorphisms and admit effective descent for the chosen topology.  Then a universal family, if it exists, is unique up to unique isomorphism.  Consequently the representing object is unique up to unique isomorphism.
\end{theorem}
\begin{remark}[Quiver representations viewpoint]
The phrase ``the moduli space of quiver representations'' is therefore shorthand.  It may mean a fine moduli scheme, a coarse moduli space, an algebraic stack, a GIT quotient, or a compactification.  The correct meaning is determined by what information is retained about families and automorphisms.
\end{remark}
\begin{proof}[Proof sketch of the theorem/proposition on this sheet]
The argument is the standard Yoneda-and-descent argument specialized to quiver representations.  A family over a test scheme $T$ corresponds, under representability, to a morphism $T\to M$.  Pulling the universal family back along this morphism recovers the given family.  If two morphisms give isomorphic pullbacks, the functorial equality on all test schemes forces the morphisms to agree.  The hypotheses on automorphisms ensure that the equality is not weakened to an equivalence relation in a groupoid.
\end{proof}
\begin{example}[A test calculation for quiver representations]
Let $T=\Spec k[s]$ and let the invariant of interest be denoted $\nu$.  A toy family $X\to T$ is chosen so that $\nu(X_t)=\nu_0$ for $s\ne0$ and the special fiber at $s=0$ records the boundary behavior.  In a concrete course one may take $X$ to be a family of divisors, modules, curves, bundles, or representations, depending on the chapter.  The moduli-theoretic lesson is that the morphism $T\to M$ detects both the general member and the limiting point.
\end{example}
\begin{exercise}
Compare the functor with the corresponding groupoid-valued functor.  Which information is lost after passing to isomorphism classes?
\end{exercise}
\begin{exercise}
Formulate a small example in which the naive set of isomorphism classes is not Hausdorff, not separated, or not functorial.  Then repair the example by adding stability, rigidification, or a stack structure.
\end{exercise}
\clearpage
\section{Sheet 66: Properness and separatedness: valuative criteria}
This sheet studies \emph{valuative criteria} as it appears in the broader chapter on \emph{properness and separatedness}.  The recurring question is whether a proposed parameter space remembers enough information about families while remaining geometric enough to compute with.
\begin{definition}[Valuative criteria viewpoint]
A valuative criteria moduli functor assigns to each test scheme $T$ the set, groupoid, or category of families over $T$ satisfying the chosen condition.  Two families are identified only when the moduli problem says that an isomorphism, not merely a bijection on fibers, has been supplied.
\end{definition}
\begin{proposition}[Valuative criteria viewpoint]
If a moduli functor for valuative criteria is represented by a scheme $M$, then natural transformations from the functor to any functor $G$ are the same as elements of $G(M)$.  In practice this converts universal constructions on families into morphisms out of $M$.
\end{proposition}
\begin{example}[Valuative criteria viewpoint]
For valuative criteria, take a base $T=\Spec R$ and replace a single object by an $R$-flat family.  Flatness prevents the numerical invariants from jumping on first-order thickenings, so it is the minimal condition under which the fibers can reasonably be compared.
\end{example}
\begin{theorem}[Valuative criteria viewpoint]
Suppose the objects in the valuative criteria problem have no non-trivial automorphisms and admit effective descent for the chosen topology.  Then a universal family, if it exists, is unique up to unique isomorphism.  Consequently the representing object is unique up to unique isomorphism.
\end{theorem}
\begin{remark}[Valuative criteria viewpoint]
The phrase ``the moduli space of valuative criteria'' is therefore shorthand.  It may mean a fine moduli scheme, a coarse moduli space, an algebraic stack, a GIT quotient, or a compactification.  The correct meaning is determined by what information is retained about families and automorphisms.
\end{remark}
\begin{proof}[Proof sketch of the theorem/proposition on this sheet]
The argument is the standard Yoneda-and-descent argument specialized to valuative criteria.  A family over a test scheme $T$ corresponds, under representability, to a morphism $T\to M$.  Pulling the universal family back along this morphism recovers the given family.  If two morphisms give isomorphic pullbacks, the functorial equality on all test schemes forces the morphisms to agree.  The hypotheses on automorphisms ensure that the equality is not weakened to an equivalence relation in a groupoid.
\end{proof}
\begin{example}[A test calculation for valuative criteria]
Let $T=\Spec k[s]$ and let the invariant of interest be denoted $\nu$.  A toy family $X\to T$ is chosen so that $\nu(X_t)=\nu_0$ for $s\ne0$ and the special fiber at $s=0$ records the boundary behavior.  In a concrete course one may take $X$ to be a family of divisors, modules, curves, bundles, or representations, depending on the chapter.  The moduli-theoretic lesson is that the morphism $T\to M$ detects both the general member and the limiting point.
\end{example}
\begin{exercise}
Write the functor of points for the moduli problem on this sheet and identify its value on $\Spec k$ and on $\Spec k[\epsilon]/(\epsilon^2)$.
\end{exercise}
\begin{exercise}
Formulate a small example in which the naive set of isomorphism classes is not Hausdorff, not separated, or not functorial.  Then repair the example by adding stability, rigidification, or a stack structure.
\end{exercise}
\clearpage
\section{Sheet 67: Properness and separatedness: stable reduction}
This sheet studies \emph{stable reduction} as it appears in the broader chapter on \emph{properness and separatedness}.  The recurring question is whether a proposed parameter space remembers enough information about families while remaining geometric enough to compute with.
\begin{definition}[Stable reduction viewpoint]
A stable reduction moduli functor assigns to each test scheme $T$ the set, groupoid, or category of families over $T$ satisfying the chosen condition.  Two families are identified only when the moduli problem says that an isomorphism, not merely a bijection on fibers, has been supplied.
\end{definition}
\begin{proposition}[Stable reduction viewpoint]
If a moduli functor for stable reduction is represented by a scheme $M$, then natural transformations from the functor to any functor $G$ are the same as elements of $G(M)$.  In practice this converts universal constructions on families into morphisms out of $M$.
\end{proposition}
\begin{example}[Stable reduction viewpoint]
For stable reduction, take a base $T=\Spec R$ and replace a single object by an $R$-flat family.  Flatness prevents the numerical invariants from jumping on first-order thickenings, so it is the minimal condition under which the fibers can reasonably be compared.
\end{example}
\begin{theorem}[Stable reduction viewpoint]
Suppose the objects in the stable reduction problem have no non-trivial automorphisms and admit effective descent for the chosen topology.  Then a universal family, if it exists, is unique up to unique isomorphism.  Consequently the representing object is unique up to unique isomorphism.
\end{theorem}
\begin{remark}[Stable reduction viewpoint]
The phrase ``the moduli space of stable reduction'' is therefore shorthand.  It may mean a fine moduli scheme, a coarse moduli space, an algebraic stack, a GIT quotient, or a compactification.  The correct meaning is determined by what information is retained about families and automorphisms.
\end{remark}
\begin{proof}[Proof sketch of the theorem/proposition on this sheet]
The argument is the standard Yoneda-and-descent argument specialized to stable reduction.  A family over a test scheme $T$ corresponds, under representability, to a morphism $T\to M$.  Pulling the universal family back along this morphism recovers the given family.  If two morphisms give isomorphic pullbacks, the functorial equality on all test schemes forces the morphisms to agree.  The hypotheses on automorphisms ensure that the equality is not weakened to an equivalence relation in a groupoid.
\end{proof}
\begin{example}[A test calculation for stable reduction]
Let $T=\Spec k[s]$ and let the invariant of interest be denoted $\nu$.  A toy family $X\to T$ is chosen so that $\nu(X_t)=\nu_0$ for $s\ne0$ and the special fiber at $s=0$ records the boundary behavior.  In a concrete course one may take $X$ to be a family of divisors, modules, curves, bundles, or representations, depending on the chapter.  The moduli-theoretic lesson is that the morphism $T\to M$ detects both the general member and the limiting point.
\end{example}
\begin{exercise}
Find one automorphism group occurring in the example and decide whether it obstructs the existence of a fine moduli scheme.
\end{exercise}
\begin{exercise}
Formulate a small example in which the naive set of isomorphism classes is not Hausdorff, not separated, or not functorial.  Then repair the example by adding stability, rigidification, or a stack structure.
\end{exercise}
\clearpage
\section{Sheet 68: Properness and separatedness: Langton semistable reduction}
This sheet studies \emph{Langton semistable reduction} as it appears in the broader chapter on \emph{properness and separatedness}.  The recurring question is whether a proposed parameter space remembers enough information about families while remaining geometric enough to compute with.
\begin{definition}[Langton semistable reduction viewpoint]
A Langton semistable reduction moduli functor assigns to each test scheme $T$ the set, groupoid, or category of families over $T$ satisfying the chosen condition.  Two families are identified only when the moduli problem says that an isomorphism, not merely a bijection on fibers, has been supplied.
\end{definition}
\begin{proposition}[Langton semistable reduction viewpoint]
If a moduli functor for Langton semistable reduction is represented by a scheme $M$, then natural transformations from the functor to any functor $G$ are the same as elements of $G(M)$.  In practice this converts universal constructions on families into morphisms out of $M$.
\end{proposition}
\begin{example}[Langton semistable reduction viewpoint]
For Langton semistable reduction, take a base $T=\Spec R$ and replace a single object by an $R$-flat family.  Flatness prevents the numerical invariants from jumping on first-order thickenings, so it is the minimal condition under which the fibers can reasonably be compared.
\end{example}
\begin{theorem}[Langton semistable reduction viewpoint]
Suppose the objects in the Langton semistable reduction problem have no non-trivial automorphisms and admit effective descent for the chosen topology.  Then a universal family, if it exists, is unique up to unique isomorphism.  Consequently the representing object is unique up to unique isomorphism.
\end{theorem}
\begin{remark}[Langton semistable reduction viewpoint]
The phrase ``the moduli space of Langton semistable reduction'' is therefore shorthand.  It may mean a fine moduli scheme, a coarse moduli space, an algebraic stack, a GIT quotient, or a compactification.  The correct meaning is determined by what information is retained about families and automorphisms.
\end{remark}
\begin{proof}[Proof sketch of the theorem/proposition on this sheet]
The argument is the standard Yoneda-and-descent argument specialized to Langton semistable reduction.  A family over a test scheme $T$ corresponds, under representability, to a morphism $T\to M$.  Pulling the universal family back along this morphism recovers the given family.  If two morphisms give isomorphic pullbacks, the functorial equality on all test schemes forces the morphisms to agree.  The hypotheses on automorphisms ensure that the equality is not weakened to an equivalence relation in a groupoid.
\end{proof}
\begin{example}[A test calculation for Langton semistable reduction]
Let $T=\Spec k[s]$ and let the invariant of interest be denoted $\nu$.  A toy family $X\to T$ is chosen so that $\nu(X_t)=\nu_0$ for $s\ne0$ and the special fiber at $s=0$ records the boundary behavior.  In a concrete course one may take $X$ to be a family of divisors, modules, curves, bundles, or representations, depending on the chapter.  The moduli-theoretic lesson is that the morphism $T\to M$ detects both the general member and the limiting point.
\end{example}
\begin{exercise}
Check a base-change square explicitly: pull the displayed family back along $T\prime\to T$ and verify which invariant is preserved.
\end{exercise}
\begin{exercise}
Formulate a small example in which the naive set of isomorphism classes is not Hausdorff, not separated, or not functorial.  Then repair the example by adding stability, rigidification, or a stack structure.
\end{exercise}
\clearpage
\section{Sheet 69: Properness and separatedness: separated diagonal}
This sheet studies \emph{separated diagonal} as it appears in the broader chapter on \emph{properness and separatedness}.  The recurring question is whether a proposed parameter space remembers enough information about families while remaining geometric enough to compute with.
\begin{definition}[Separated diagonal viewpoint]
A separated diagonal moduli functor assigns to each test scheme $T$ the set, groupoid, or category of families over $T$ satisfying the chosen condition.  Two families are identified only when the moduli problem says that an isomorphism, not merely a bijection on fibers, has been supplied.
\end{definition}
\begin{proposition}[Separated diagonal viewpoint]
If a moduli functor for separated diagonal is represented by a scheme $M$, then natural transformations from the functor to any functor $G$ are the same as elements of $G(M)$.  In practice this converts universal constructions on families into morphisms out of $M$.
\end{proposition}
\begin{example}[Separated diagonal viewpoint]
For separated diagonal, take a base $T=\Spec R$ and replace a single object by an $R$-flat family.  Flatness prevents the numerical invariants from jumping on first-order thickenings, so it is the minimal condition under which the fibers can reasonably be compared.
\end{example}
\begin{theorem}[Separated diagonal viewpoint]
Suppose the objects in the separated diagonal problem have no non-trivial automorphisms and admit effective descent for the chosen topology.  Then a universal family, if it exists, is unique up to unique isomorphism.  Consequently the representing object is unique up to unique isomorphism.
\end{theorem}
\begin{remark}[Separated diagonal viewpoint]
The phrase ``the moduli space of separated diagonal'' is therefore shorthand.  It may mean a fine moduli scheme, a coarse moduli space, an algebraic stack, a GIT quotient, or a compactification.  The correct meaning is determined by what information is retained about families and automorphisms.
\end{remark}
\begin{proof}[Proof sketch of the theorem/proposition on this sheet]
The argument is the standard Yoneda-and-descent argument specialized to separated diagonal.  A family over a test scheme $T$ corresponds, under representability, to a morphism $T\to M$.  Pulling the universal family back along this morphism recovers the given family.  If two morphisms give isomorphic pullbacks, the functorial equality on all test schemes forces the morphisms to agree.  The hypotheses on automorphisms ensure that the equality is not weakened to an equivalence relation in a groupoid.
\end{proof}
\begin{example}[A test calculation for separated diagonal]
Let $T=\Spec k[s]$ and let the invariant of interest be denoted $\nu$.  A toy family $X\to T$ is chosen so that $\nu(X_t)=\nu_0$ for $s\ne0$ and the special fiber at $s=0$ records the boundary behavior.  In a concrete course one may take $X$ to be a family of divisors, modules, curves, bundles, or representations, depending on the chapter.  The moduli-theoretic lesson is that the morphism $T\to M$ detects both the general member and the limiting point.
\end{example}
\begin{exercise}
Construct a degeneration over $\Spec k[[t]]$ whose generic fiber satisfies the condition but whose special fiber forces a compactification.
\end{exercise}
\begin{exercise}
Formulate a small example in which the naive set of isomorphism classes is not Hausdorff, not separated, or not functorial.  Then repair the example by adding stability, rigidification, or a stack structure.
\end{exercise}
\clearpage
\section{Sheet 70: Properness and separatedness: compactifications}
This sheet studies \emph{compactifications} as it appears in the broader chapter on \emph{properness and separatedness}.  The recurring question is whether a proposed parameter space remembers enough information about families while remaining geometric enough to compute with.
\begin{definition}[Compactifications viewpoint]
A compactifications moduli functor assigns to each test scheme $T$ the set, groupoid, or category of families over $T$ satisfying the chosen condition.  Two families are identified only when the moduli problem says that an isomorphism, not merely a bijection on fibers, has been supplied.
\end{definition}
\begin{proposition}[Compactifications viewpoint]
If a moduli functor for compactifications is represented by a scheme $M$, then natural transformations from the functor to any functor $G$ are the same as elements of $G(M)$.  In practice this converts universal constructions on families into morphisms out of $M$.
\end{proposition}
\begin{example}[Compactifications viewpoint]
For compactifications, take a base $T=\Spec R$ and replace a single object by an $R$-flat family.  Flatness prevents the numerical invariants from jumping on first-order thickenings, so it is the minimal condition under which the fibers can reasonably be compared.
\end{example}
\begin{theorem}[Compactifications viewpoint]
Suppose the objects in the compactifications problem have no non-trivial automorphisms and admit effective descent for the chosen topology.  Then a universal family, if it exists, is unique up to unique isomorphism.  Consequently the representing object is unique up to unique isomorphism.
\end{theorem}
\begin{remark}[Compactifications viewpoint]
The phrase ``the moduli space of compactifications'' is therefore shorthand.  It may mean a fine moduli scheme, a coarse moduli space, an algebraic stack, a GIT quotient, or a compactification.  The correct meaning is determined by what information is retained about families and automorphisms.
\end{remark}
\begin{proof}[Proof sketch of the theorem/proposition on this sheet]
The argument is the standard Yoneda-and-descent argument specialized to compactifications.  A family over a test scheme $T$ corresponds, under representability, to a morphism $T\to M$.  Pulling the universal family back along this morphism recovers the given family.  If two morphisms give isomorphic pullbacks, the functorial equality on all test schemes forces the morphisms to agree.  The hypotheses on automorphisms ensure that the equality is not weakened to an equivalence relation in a groupoid.
\end{proof}
\begin{example}[A test calculation for compactifications]
Let $T=\Spec k[s]$ and let the invariant of interest be denoted $\nu$.  A toy family $X\to T$ is chosen so that $\nu(X_t)=\nu_0$ for $s\ne0$ and the special fiber at $s=0$ records the boundary behavior.  In a concrete course one may take $X$ to be a family of divisors, modules, curves, bundles, or representations, depending on the chapter.  The moduli-theoretic lesson is that the morphism $T\to M$ detects both the general member and the limiting point.
\end{example}
\begin{exercise}
Compare the functor with the corresponding groupoid-valued functor.  Which information is lost after passing to isomorphism classes?
\end{exercise}
\begin{exercise}
Formulate a small example in which the naive set of isomorphism classes is not Hausdorff, not separated, or not functorial.  Then repair the example by adding stability, rigidification, or a stack structure.
\end{exercise}
\clearpage
\section{Sheet 71: Universality revisited: versal families}
This sheet studies \emph{versal families} as it appears in the broader chapter on \emph{universality revisited}.  The recurring question is whether a proposed parameter space remembers enough information about families while remaining geometric enough to compute with.
\begin{definition}[Versal families viewpoint]
A versal families moduli functor assigns to each test scheme $T$ the set, groupoid, or category of families over $T$ satisfying the chosen condition.  Two families are identified only when the moduli problem says that an isomorphism, not merely a bijection on fibers, has been supplied.
\end{definition}
\begin{proposition}[Versal families viewpoint]
If a moduli functor for versal families is represented by a scheme $M$, then natural transformations from the functor to any functor $G$ are the same as elements of $G(M)$.  In practice this converts universal constructions on families into morphisms out of $M$.
\end{proposition}
\begin{example}[Versal families viewpoint]
For versal families, take a base $T=\Spec R$ and replace a single object by an $R$-flat family.  Flatness prevents the numerical invariants from jumping on first-order thickenings, so it is the minimal condition under which the fibers can reasonably be compared.
\end{example}
\begin{theorem}[Versal families viewpoint]
Suppose the objects in the versal families problem have no non-trivial automorphisms and admit effective descent for the chosen topology.  Then a universal family, if it exists, is unique up to unique isomorphism.  Consequently the representing object is unique up to unique isomorphism.
\end{theorem}
\begin{remark}[Versal families viewpoint]
The phrase ``the moduli space of versal families'' is therefore shorthand.  It may mean a fine moduli scheme, a coarse moduli space, an algebraic stack, a GIT quotient, or a compactification.  The correct meaning is determined by what information is retained about families and automorphisms.
\end{remark}
\begin{proof}[Proof sketch of the theorem/proposition on this sheet]
The argument is the standard Yoneda-and-descent argument specialized to versal families.  A family over a test scheme $T$ corresponds, under representability, to a morphism $T\to M$.  Pulling the universal family back along this morphism recovers the given family.  If two morphisms give isomorphic pullbacks, the functorial equality on all test schemes forces the morphisms to agree.  The hypotheses on automorphisms ensure that the equality is not weakened to an equivalence relation in a groupoid.
\end{proof}
\begin{example}[A test calculation for versal families]
Let $T=\Spec k[s]$ and let the invariant of interest be denoted $\nu$.  A toy family $X\to T$ is chosen so that $\nu(X_t)=\nu_0$ for $s\ne0$ and the special fiber at $s=0$ records the boundary behavior.  In a concrete course one may take $X$ to be a family of divisors, modules, curves, bundles, or representations, depending on the chapter.  The moduli-theoretic lesson is that the morphism $T\to M$ detects both the general member and the limiting point.
\end{example}
\begin{exercise}
Write the functor of points for the moduli problem on this sheet and identify its value on $\Spec k$ and on $\Spec k[\epsilon]/(\epsilon^2)$.
\end{exercise}
\begin{exercise}
Formulate a small example in which the naive set of isomorphism classes is not Hausdorff, not separated, or not functorial.  Then repair the example by adding stability, rigidification, or a stack structure.
\end{exercise}
\clearpage
\section{Sheet 72: Universality revisited: miniversal deformations}
This sheet studies \emph{miniversal deformations} as it appears in the broader chapter on \emph{universality revisited}.  The recurring question is whether a proposed parameter space remembers enough information about families while remaining geometric enough to compute with.
\begin{definition}[Miniversal deformations viewpoint]
A miniversal deformations moduli functor assigns to each test scheme $T$ the set, groupoid, or category of families over $T$ satisfying the chosen condition.  Two families are identified only when the moduli problem says that an isomorphism, not merely a bijection on fibers, has been supplied.
\end{definition}
\begin{proposition}[Miniversal deformations viewpoint]
If a moduli functor for miniversal deformations is represented by a scheme $M$, then natural transformations from the functor to any functor $G$ are the same as elements of $G(M)$.  In practice this converts universal constructions on families into morphisms out of $M$.
\end{proposition}
\begin{example}[Miniversal deformations viewpoint]
For miniversal deformations, take a base $T=\Spec R$ and replace a single object by an $R$-flat family.  Flatness prevents the numerical invariants from jumping on first-order thickenings, so it is the minimal condition under which the fibers can reasonably be compared.
\end{example}
\begin{theorem}[Miniversal deformations viewpoint]
Suppose the objects in the miniversal deformations problem have no non-trivial automorphisms and admit effective descent for the chosen topology.  Then a universal family, if it exists, is unique up to unique isomorphism.  Consequently the representing object is unique up to unique isomorphism.
\end{theorem}
\begin{remark}[Miniversal deformations viewpoint]
The phrase ``the moduli space of miniversal deformations'' is therefore shorthand.  It may mean a fine moduli scheme, a coarse moduli space, an algebraic stack, a GIT quotient, or a compactification.  The correct meaning is determined by what information is retained about families and automorphisms.
\end{remark}
\begin{proof}[Proof sketch of the theorem/proposition on this sheet]
The argument is the standard Yoneda-and-descent argument specialized to miniversal deformations.  A family over a test scheme $T$ corresponds, under representability, to a morphism $T\to M$.  Pulling the universal family back along this morphism recovers the given family.  If two morphisms give isomorphic pullbacks, the functorial equality on all test schemes forces the morphisms to agree.  The hypotheses on automorphisms ensure that the equality is not weakened to an equivalence relation in a groupoid.
\end{proof}
\begin{example}[A test calculation for miniversal deformations]
Let $T=\Spec k[s]$ and let the invariant of interest be denoted $\nu$.  A toy family $X\to T$ is chosen so that $\nu(X_t)=\nu_0$ for $s\ne0$ and the special fiber at $s=0$ records the boundary behavior.  In a concrete course one may take $X$ to be a family of divisors, modules, curves, bundles, or representations, depending on the chapter.  The moduli-theoretic lesson is that the morphism $T\to M$ detects both the general member and the limiting point.
\end{example}
\begin{exercise}
Find one automorphism group occurring in the example and decide whether it obstructs the existence of a fine moduli scheme.
\end{exercise}
\begin{exercise}
Formulate a small example in which the naive set of isomorphism classes is not Hausdorff, not separated, or not functorial.  Then repair the example by adding stability, rigidification, or a stack structure.
\end{exercise}
\clearpage
\section{Sheet 73: Universality revisited: automorphisms}
This sheet studies \emph{automorphisms} as it appears in the broader chapter on \emph{universality revisited}.  The recurring question is whether a proposed parameter space remembers enough information about families while remaining geometric enough to compute with.
\begin{definition}[Automorphisms viewpoint]
A automorphisms moduli functor assigns to each test scheme $T$ the set, groupoid, or category of families over $T$ satisfying the chosen condition.  Two families are identified only when the moduli problem says that an isomorphism, not merely a bijection on fibers, has been supplied.
\end{definition}
\begin{proposition}[Automorphisms viewpoint]
If a moduli functor for automorphisms is represented by a scheme $M$, then natural transformations from the functor to any functor $G$ are the same as elements of $G(M)$.  In practice this converts universal constructions on families into morphisms out of $M$.
\end{proposition}
\begin{example}[Automorphisms viewpoint]
For automorphisms, take a base $T=\Spec R$ and replace a single object by an $R$-flat family.  Flatness prevents the numerical invariants from jumping on first-order thickenings, so it is the minimal condition under which the fibers can reasonably be compared.
\end{example}
\begin{theorem}[Automorphisms viewpoint]
Suppose the objects in the automorphisms problem have no non-trivial automorphisms and admit effective descent for the chosen topology.  Then a universal family, if it exists, is unique up to unique isomorphism.  Consequently the representing object is unique up to unique isomorphism.
\end{theorem}
\begin{remark}[Automorphisms viewpoint]
The phrase ``the moduli space of automorphisms'' is therefore shorthand.  It may mean a fine moduli scheme, a coarse moduli space, an algebraic stack, a GIT quotient, or a compactification.  The correct meaning is determined by what information is retained about families and automorphisms.
\end{remark}
\begin{proof}[Proof sketch of the theorem/proposition on this sheet]
The argument is the standard Yoneda-and-descent argument specialized to automorphisms.  A family over a test scheme $T$ corresponds, under representability, to a morphism $T\to M$.  Pulling the universal family back along this morphism recovers the given family.  If two morphisms give isomorphic pullbacks, the functorial equality on all test schemes forces the morphisms to agree.  The hypotheses on automorphisms ensure that the equality is not weakened to an equivalence relation in a groupoid.
\end{proof}
\begin{example}[A test calculation for automorphisms]
Let $T=\Spec k[s]$ and let the invariant of interest be denoted $\nu$.  A toy family $X\to T$ is chosen so that $\nu(X_t)=\nu_0$ for $s\ne0$ and the special fiber at $s=0$ records the boundary behavior.  In a concrete course one may take $X$ to be a family of divisors, modules, curves, bundles, or representations, depending on the chapter.  The moduli-theoretic lesson is that the morphism $T\to M$ detects both the general member and the limiting point.
\end{example}
\begin{exercise}
Check a base-change square explicitly: pull the displayed family back along $T\prime\to T$ and verify which invariant is preserved.
\end{exercise}
\begin{exercise}
Formulate a small example in which the naive set of isomorphism classes is not Hausdorff, not separated, or not functorial.  Then repair the example by adding stability, rigidification, or a stack structure.
\end{exercise}
\clearpage
\section{Sheet 74: Universality revisited: rigidification}
This sheet studies \emph{rigidification} as it appears in the broader chapter on \emph{universality revisited}.  The recurring question is whether a proposed parameter space remembers enough information about families while remaining geometric enough to compute with.
\begin{definition}[Rigidification viewpoint]
A rigidification moduli functor assigns to each test scheme $T$ the set, groupoid, or category of families over $T$ satisfying the chosen condition.  Two families are identified only when the moduli problem says that an isomorphism, not merely a bijection on fibers, has been supplied.
\end{definition}
\begin{proposition}[Rigidification viewpoint]
If a moduli functor for rigidification is represented by a scheme $M$, then natural transformations from the functor to any functor $G$ are the same as elements of $G(M)$.  In practice this converts universal constructions on families into morphisms out of $M$.
\end{proposition}
\begin{example}[Rigidification viewpoint]
For rigidification, take a base $T=\Spec R$ and replace a single object by an $R$-flat family.  Flatness prevents the numerical invariants from jumping on first-order thickenings, so it is the minimal condition under which the fibers can reasonably be compared.
\end{example}
\begin{theorem}[Rigidification viewpoint]
Suppose the objects in the rigidification problem have no non-trivial automorphisms and admit effective descent for the chosen topology.  Then a universal family, if it exists, is unique up to unique isomorphism.  Consequently the representing object is unique up to unique isomorphism.
\end{theorem}
\begin{remark}[Rigidification viewpoint]
The phrase ``the moduli space of rigidification'' is therefore shorthand.  It may mean a fine moduli scheme, a coarse moduli space, an algebraic stack, a GIT quotient, or a compactification.  The correct meaning is determined by what information is retained about families and automorphisms.
\end{remark}
\begin{proof}[Proof sketch of the theorem/proposition on this sheet]
The argument is the standard Yoneda-and-descent argument specialized to rigidification.  A family over a test scheme $T$ corresponds, under representability, to a morphism $T\to M$.  Pulling the universal family back along this morphism recovers the given family.  If two morphisms give isomorphic pullbacks, the functorial equality on all test schemes forces the morphisms to agree.  The hypotheses on automorphisms ensure that the equality is not weakened to an equivalence relation in a groupoid.
\end{proof}
\begin{example}[A test calculation for rigidification]
Let $T=\Spec k[s]$ and let the invariant of interest be denoted $\nu$.  A toy family $X\to T$ is chosen so that $\nu(X_t)=\nu_0$ for $s\ne0$ and the special fiber at $s=0$ records the boundary behavior.  In a concrete course one may take $X$ to be a family of divisors, modules, curves, bundles, or representations, depending on the chapter.  The moduli-theoretic lesson is that the morphism $T\to M$ detects both the general member and the limiting point.
\end{example}
\begin{exercise}
Construct a degeneration over $\Spec k[[t]]$ whose generic fiber satisfies the condition but whose special fiber forces a compactification.
\end{exercise}
\begin{exercise}
Formulate a small example in which the naive set of isomorphism classes is not Hausdorff, not separated, or not functorial.  Then repair the example by adding stability, rigidification, or a stack structure.
\end{exercise}
\clearpage
\section{Sheet 75: Universality revisited: period maps}
This sheet studies \emph{period maps} as it appears in the broader chapter on \emph{universality revisited}.  The recurring question is whether a proposed parameter space remembers enough information about families while remaining geometric enough to compute with.
\begin{definition}[Period maps viewpoint]
A period maps moduli functor assigns to each test scheme $T$ the set, groupoid, or category of families over $T$ satisfying the chosen condition.  Two families are identified only when the moduli problem says that an isomorphism, not merely a bijection on fibers, has been supplied.
\end{definition}
\begin{proposition}[Period maps viewpoint]
If a moduli functor for period maps is represented by a scheme $M$, then natural transformations from the functor to any functor $G$ are the same as elements of $G(M)$.  In practice this converts universal constructions on families into morphisms out of $M$.
\end{proposition}
\begin{example}[Period maps viewpoint]
For period maps, take a base $T=\Spec R$ and replace a single object by an $R$-flat family.  Flatness prevents the numerical invariants from jumping on first-order thickenings, so it is the minimal condition under which the fibers can reasonably be compared.
\end{example}
\begin{theorem}[Period maps viewpoint]
Suppose the objects in the period maps problem have no non-trivial automorphisms and admit effective descent for the chosen topology.  Then a universal family, if it exists, is unique up to unique isomorphism.  Consequently the representing object is unique up to unique isomorphism.
\end{theorem}
\begin{remark}[Period maps viewpoint]
The phrase ``the moduli space of period maps'' is therefore shorthand.  It may mean a fine moduli scheme, a coarse moduli space, an algebraic stack, a GIT quotient, or a compactification.  The correct meaning is determined by what information is retained about families and automorphisms.
\end{remark}
\begin{proof}[Proof sketch of the theorem/proposition on this sheet]
The argument is the standard Yoneda-and-descent argument specialized to period maps.  A family over a test scheme $T$ corresponds, under representability, to a morphism $T\to M$.  Pulling the universal family back along this morphism recovers the given family.  If two morphisms give isomorphic pullbacks, the functorial equality on all test schemes forces the morphisms to agree.  The hypotheses on automorphisms ensure that the equality is not weakened to an equivalence relation in a groupoid.
\end{proof}
\begin{example}[A test calculation for period maps]
Let $T=\Spec k[s]$ and let the invariant of interest be denoted $\nu$.  A toy family $X\to T$ is chosen so that $\nu(X_t)=\nu_0$ for $s\ne0$ and the special fiber at $s=0$ records the boundary behavior.  In a concrete course one may take $X$ to be a family of divisors, modules, curves, bundles, or representations, depending on the chapter.  The moduli-theoretic lesson is that the morphism $T\to M$ detects both the general member and the limiting point.
\end{example}
\begin{exercise}
Compare the functor with the corresponding groupoid-valued functor.  Which information is lost after passing to isomorphism classes?
\end{exercise}
\begin{exercise}
Formulate a small example in which the naive set of isomorphism classes is not Hausdorff, not separated, or not functorial.  Then repair the example by adding stability, rigidification, or a stack structure.
\end{exercise}
\clearpage
\section{Sheet 76: Arithmetic moduli: integral models}
This sheet studies \emph{integral models} as it appears in the broader chapter on \emph{arithmetic moduli}.  The recurring question is whether a proposed parameter space remembers enough information about families while remaining geometric enough to compute with.
\begin{definition}[Integral models viewpoint]
A integral models moduli functor assigns to each test scheme $T$ the set, groupoid, or category of families over $T$ satisfying the chosen condition.  Two families are identified only when the moduli problem says that an isomorphism, not merely a bijection on fibers, has been supplied.
\end{definition}
\begin{proposition}[Integral models viewpoint]
If a moduli functor for integral models is represented by a scheme $M$, then natural transformations from the functor to any functor $G$ are the same as elements of $G(M)$.  In practice this converts universal constructions on families into morphisms out of $M$.
\end{proposition}
\begin{example}[Integral models viewpoint]
For integral models, take a base $T=\Spec R$ and replace a single object by an $R$-flat family.  Flatness prevents the numerical invariants from jumping on first-order thickenings, so it is the minimal condition under which the fibers can reasonably be compared.
\end{example}
\begin{theorem}[Integral models viewpoint]
Suppose the objects in the integral models problem have no non-trivial automorphisms and admit effective descent for the chosen topology.  Then a universal family, if it exists, is unique up to unique isomorphism.  Consequently the representing object is unique up to unique isomorphism.
\end{theorem}
\begin{remark}[Integral models viewpoint]
The phrase ``the moduli space of integral models'' is therefore shorthand.  It may mean a fine moduli scheme, a coarse moduli space, an algebraic stack, a GIT quotient, or a compactification.  The correct meaning is determined by what information is retained about families and automorphisms.
\end{remark}
\begin{proof}[Proof sketch of the theorem/proposition on this sheet]
The argument is the standard Yoneda-and-descent argument specialized to integral models.  A family over a test scheme $T$ corresponds, under representability, to a morphism $T\to M$.  Pulling the universal family back along this morphism recovers the given family.  If two morphisms give isomorphic pullbacks, the functorial equality on all test schemes forces the morphisms to agree.  The hypotheses on automorphisms ensure that the equality is not weakened to an equivalence relation in a groupoid.
\end{proof}
\begin{example}[A test calculation for integral models]
Let $T=\Spec k[s]$ and let the invariant of interest be denoted $\nu$.  A toy family $X\to T$ is chosen so that $\nu(X_t)=\nu_0$ for $s\ne0$ and the special fiber at $s=0$ records the boundary behavior.  In a concrete course one may take $X$ to be a family of divisors, modules, curves, bundles, or representations, depending on the chapter.  The moduli-theoretic lesson is that the morphism $T\to M$ detects both the general member and the limiting point.
\end{example}
\begin{exercise}
Write the functor of points for the moduli problem on this sheet and identify its value on $\Spec k$ and on $\Spec k[\epsilon]/(\epsilon^2)$.
\end{exercise}
\begin{exercise}
Formulate a small example in which the naive set of isomorphism classes is not Hausdorff, not separated, or not functorial.  Then repair the example by adding stability, rigidification, or a stack structure.
\end{exercise}
\clearpage
\section{Sheet 77: Arithmetic moduli: level structures}
This sheet studies \emph{level structures} as it appears in the broader chapter on \emph{arithmetic moduli}.  The recurring question is whether a proposed parameter space remembers enough information about families while remaining geometric enough to compute with.
\begin{definition}[Level structures viewpoint]
A level structures moduli functor assigns to each test scheme $T$ the set, groupoid, or category of families over $T$ satisfying the chosen condition.  Two families are identified only when the moduli problem says that an isomorphism, not merely a bijection on fibers, has been supplied.
\end{definition}
\begin{proposition}[Level structures viewpoint]
If a moduli functor for level structures is represented by a scheme $M$, then natural transformations from the functor to any functor $G$ are the same as elements of $G(M)$.  In practice this converts universal constructions on families into morphisms out of $M$.
\end{proposition}
\begin{example}[Level structures viewpoint]
For level structures, take a base $T=\Spec R$ and replace a single object by an $R$-flat family.  Flatness prevents the numerical invariants from jumping on first-order thickenings, so it is the minimal condition under which the fibers can reasonably be compared.
\end{example}
\begin{theorem}[Level structures viewpoint]
Suppose the objects in the level structures problem have no non-trivial automorphisms and admit effective descent for the chosen topology.  Then a universal family, if it exists, is unique up to unique isomorphism.  Consequently the representing object is unique up to unique isomorphism.
\end{theorem}
\begin{remark}[Level structures viewpoint]
The phrase ``the moduli space of level structures'' is therefore shorthand.  It may mean a fine moduli scheme, a coarse moduli space, an algebraic stack, a GIT quotient, or a compactification.  The correct meaning is determined by what information is retained about families and automorphisms.
\end{remark}
\begin{proof}[Proof sketch of the theorem/proposition on this sheet]
The argument is the standard Yoneda-and-descent argument specialized to level structures.  A family over a test scheme $T$ corresponds, under representability, to a morphism $T\to M$.  Pulling the universal family back along this morphism recovers the given family.  If two morphisms give isomorphic pullbacks, the functorial equality on all test schemes forces the morphisms to agree.  The hypotheses on automorphisms ensure that the equality is not weakened to an equivalence relation in a groupoid.
\end{proof}
\begin{example}[A test calculation for level structures]
Let $T=\Spec k[s]$ and let the invariant of interest be denoted $\nu$.  A toy family $X\to T$ is chosen so that $\nu(X_t)=\nu_0$ for $s\ne0$ and the special fiber at $s=0$ records the boundary behavior.  In a concrete course one may take $X$ to be a family of divisors, modules, curves, bundles, or representations, depending on the chapter.  The moduli-theoretic lesson is that the morphism $T\to M$ detects both the general member and the limiting point.
\end{example}
\begin{exercise}
Find one automorphism group occurring in the example and decide whether it obstructs the existence of a fine moduli scheme.
\end{exercise}
\begin{exercise}
Formulate a small example in which the naive set of isomorphism classes is not Hausdorff, not separated, or not functorial.  Then repair the example by adding stability, rigidification, or a stack structure.
\end{exercise}
\clearpage
\section{Sheet 78: Arithmetic moduli: modular curves}
This sheet studies \emph{modular curves} as it appears in the broader chapter on \emph{arithmetic moduli}.  The recurring question is whether a proposed parameter space remembers enough information about families while remaining geometric enough to compute with.
\begin{definition}[Modular curves viewpoint]
A modular curves moduli functor assigns to each test scheme $T$ the set, groupoid, or category of families over $T$ satisfying the chosen condition.  Two families are identified only when the moduli problem says that an isomorphism, not merely a bijection on fibers, has been supplied.
\end{definition}
\begin{proposition}[Modular curves viewpoint]
If a moduli functor for modular curves is represented by a scheme $M$, then natural transformations from the functor to any functor $G$ are the same as elements of $G(M)$.  In practice this converts universal constructions on families into morphisms out of $M$.
\end{proposition}
\begin{example}[Modular curves viewpoint]
For modular curves, take a base $T=\Spec R$ and replace a single object by an $R$-flat family.  Flatness prevents the numerical invariants from jumping on first-order thickenings, so it is the minimal condition under which the fibers can reasonably be compared.
\end{example}
\begin{theorem}[Modular curves viewpoint]
Suppose the objects in the modular curves problem have no non-trivial automorphisms and admit effective descent for the chosen topology.  Then a universal family, if it exists, is unique up to unique isomorphism.  Consequently the representing object is unique up to unique isomorphism.
\end{theorem}
\begin{remark}[Modular curves viewpoint]
The phrase ``the moduli space of modular curves'' is therefore shorthand.  It may mean a fine moduli scheme, a coarse moduli space, an algebraic stack, a GIT quotient, or a compactification.  The correct meaning is determined by what information is retained about families and automorphisms.
\end{remark}
\begin{proof}[Proof sketch of the theorem/proposition on this sheet]
The argument is the standard Yoneda-and-descent argument specialized to modular curves.  A family over a test scheme $T$ corresponds, under representability, to a morphism $T\to M$.  Pulling the universal family back along this morphism recovers the given family.  If two morphisms give isomorphic pullbacks, the functorial equality on all test schemes forces the morphisms to agree.  The hypotheses on automorphisms ensure that the equality is not weakened to an equivalence relation in a groupoid.
\end{proof}
\begin{example}[A test calculation for modular curves]
Let $T=\Spec k[s]$ and let the invariant of interest be denoted $\nu$.  A toy family $X\to T$ is chosen so that $\nu(X_t)=\nu_0$ for $s\ne0$ and the special fiber at $s=0$ records the boundary behavior.  In a concrete course one may take $X$ to be a family of divisors, modules, curves, bundles, or representations, depending on the chapter.  The moduli-theoretic lesson is that the morphism $T\to M$ detects both the general member and the limiting point.
\end{example}
\begin{exercise}
Check a base-change square explicitly: pull the displayed family back along $T\prime\to T$ and verify which invariant is preserved.
\end{exercise}
\begin{exercise}
Formulate a small example in which the naive set of isomorphism classes is not Hausdorff, not separated, or not functorial.  Then repair the example by adding stability, rigidification, or a stack structure.
\end{exercise}
\clearpage
\section{Sheet 79: Arithmetic moduli: Hecke correspondences}
This sheet studies \emph{Hecke correspondences} as it appears in the broader chapter on \emph{arithmetic moduli}.  The recurring question is whether a proposed parameter space remembers enough information about families while remaining geometric enough to compute with.
\begin{definition}[Hecke correspondences viewpoint]
A Hecke correspondences moduli functor assigns to each test scheme $T$ the set, groupoid, or category of families over $T$ satisfying the chosen condition.  Two families are identified only when the moduli problem says that an isomorphism, not merely a bijection on fibers, has been supplied.
\end{definition}
\begin{proposition}[Hecke correspondences viewpoint]
If a moduli functor for Hecke correspondences is represented by a scheme $M$, then natural transformations from the functor to any functor $G$ are the same as elements of $G(M)$.  In practice this converts universal constructions on families into morphisms out of $M$.
\end{proposition}
\begin{example}[Hecke correspondences viewpoint]
For Hecke correspondences, take a base $T=\Spec R$ and replace a single object by an $R$-flat family.  Flatness prevents the numerical invariants from jumping on first-order thickenings, so it is the minimal condition under which the fibers can reasonably be compared.
\end{example}
\begin{theorem}[Hecke correspondences viewpoint]
Suppose the objects in the Hecke correspondences problem have no non-trivial automorphisms and admit effective descent for the chosen topology.  Then a universal family, if it exists, is unique up to unique isomorphism.  Consequently the representing object is unique up to unique isomorphism.
\end{theorem}
\begin{remark}[Hecke correspondences viewpoint]
The phrase ``the moduli space of Hecke correspondences'' is therefore shorthand.  It may mean a fine moduli scheme, a coarse moduli space, an algebraic stack, a GIT quotient, or a compactification.  The correct meaning is determined by what information is retained about families and automorphisms.
\end{remark}
\begin{proof}[Proof sketch of the theorem/proposition on this sheet]
The argument is the standard Yoneda-and-descent argument specialized to Hecke correspondences.  A family over a test scheme $T$ corresponds, under representability, to a morphism $T\to M$.  Pulling the universal family back along this morphism recovers the given family.  If two morphisms give isomorphic pullbacks, the functorial equality on all test schemes forces the morphisms to agree.  The hypotheses on automorphisms ensure that the equality is not weakened to an equivalence relation in a groupoid.
\end{proof}
\begin{example}[A test calculation for Hecke correspondences]
Let $T=\Spec k[s]$ and let the invariant of interest be denoted $\nu$.  A toy family $X\to T$ is chosen so that $\nu(X_t)=\nu_0$ for $s\ne0$ and the special fiber at $s=0$ records the boundary behavior.  In a concrete course one may take $X$ to be a family of divisors, modules, curves, bundles, or representations, depending on the chapter.  The moduli-theoretic lesson is that the morphism $T\to M$ detects both the general member and the limiting point.
\end{example}
\begin{exercise}
Construct a degeneration over $\Spec k[[t]]$ whose generic fiber satisfies the condition but whose special fiber forces a compactification.
\end{exercise}
\begin{exercise}
Formulate a small example in which the naive set of isomorphism classes is not Hausdorff, not separated, or not functorial.  Then repair the example by adding stability, rigidification, or a stack structure.
\end{exercise}
\clearpage
\section{Sheet 80: Arithmetic moduli: heights}
This sheet studies \emph{heights} as it appears in the broader chapter on \emph{arithmetic moduli}.  The recurring question is whether a proposed parameter space remembers enough information about families while remaining geometric enough to compute with.
\begin{definition}[Heights viewpoint]
A heights moduli functor assigns to each test scheme $T$ the set, groupoid, or category of families over $T$ satisfying the chosen condition.  Two families are identified only when the moduli problem says that an isomorphism, not merely a bijection on fibers, has been supplied.
\end{definition}
\begin{proposition}[Heights viewpoint]
If a moduli functor for heights is represented by a scheme $M$, then natural transformations from the functor to any functor $G$ are the same as elements of $G(M)$.  In practice this converts universal constructions on families into morphisms out of $M$.
\end{proposition}
\begin{example}[Heights viewpoint]
For heights, take a base $T=\Spec R$ and replace a single object by an $R$-flat family.  Flatness prevents the numerical invariants from jumping on first-order thickenings, so it is the minimal condition under which the fibers can reasonably be compared.
\end{example}
\begin{theorem}[Heights viewpoint]
Suppose the objects in the heights problem have no non-trivial automorphisms and admit effective descent for the chosen topology.  Then a universal family, if it exists, is unique up to unique isomorphism.  Consequently the representing object is unique up to unique isomorphism.
\end{theorem}
\begin{remark}[Heights viewpoint]
The phrase ``the moduli space of heights'' is therefore shorthand.  It may mean a fine moduli scheme, a coarse moduli space, an algebraic stack, a GIT quotient, or a compactification.  The correct meaning is determined by what information is retained about families and automorphisms.
\end{remark}
\begin{proof}[Proof sketch of the theorem/proposition on this sheet]
The argument is the standard Yoneda-and-descent argument specialized to heights.  A family over a test scheme $T$ corresponds, under representability, to a morphism $T\to M$.  Pulling the universal family back along this morphism recovers the given family.  If two morphisms give isomorphic pullbacks, the functorial equality on all test schemes forces the morphisms to agree.  The hypotheses on automorphisms ensure that the equality is not weakened to an equivalence relation in a groupoid.
\end{proof}
\begin{example}[A test calculation for heights]
Let $T=\Spec k[s]$ and let the invariant of interest be denoted $\nu$.  A toy family $X\to T$ is chosen so that $\nu(X_t)=\nu_0$ for $s\ne0$ and the special fiber at $s=0$ records the boundary behavior.  In a concrete course one may take $X$ to be a family of divisors, modules, curves, bundles, or representations, depending on the chapter.  The moduli-theoretic lesson is that the morphism $T\to M$ detects both the general member and the limiting point.
\end{example}
\begin{exercise}
Compare the functor with the corresponding groupoid-valued functor.  Which information is lost after passing to isomorphism classes?
\end{exercise}
\begin{exercise}
Formulate a small example in which the naive set of isomorphism classes is not Hausdorff, not separated, or not functorial.  Then repair the example by adding stability, rigidification, or a stack structure.
\end{exercise}
\clearpage
\section{Sheet 81: Tropical and logarithmic shadows: dual complexes}
This sheet studies \emph{dual complexes} as it appears in the broader chapter on \emph{tropical and logarithmic shadows}.  The recurring question is whether a proposed parameter space remembers enough information about families while remaining geometric enough to compute with.
\begin{definition}[Dual complexes viewpoint]
A dual complexes moduli functor assigns to each test scheme $T$ the set, groupoid, or category of families over $T$ satisfying the chosen condition.  Two families are identified only when the moduli problem says that an isomorphism, not merely a bijection on fibers, has been supplied.
\end{definition}
\begin{proposition}[Dual complexes viewpoint]
If a moduli functor for dual complexes is represented by a scheme $M$, then natural transformations from the functor to any functor $G$ are the same as elements of $G(M)$.  In practice this converts universal constructions on families into morphisms out of $M$.
\end{proposition}
\begin{example}[Dual complexes viewpoint]
For dual complexes, take a base $T=\Spec R$ and replace a single object by an $R$-flat family.  Flatness prevents the numerical invariants from jumping on first-order thickenings, so it is the minimal condition under which the fibers can reasonably be compared.
\end{example}
\begin{theorem}[Dual complexes viewpoint]
Suppose the objects in the dual complexes problem have no non-trivial automorphisms and admit effective descent for the chosen topology.  Then a universal family, if it exists, is unique up to unique isomorphism.  Consequently the representing object is unique up to unique isomorphism.
\end{theorem}
\begin{remark}[Dual complexes viewpoint]
The phrase ``the moduli space of dual complexes'' is therefore shorthand.  It may mean a fine moduli scheme, a coarse moduli space, an algebraic stack, a GIT quotient, or a compactification.  The correct meaning is determined by what information is retained about families and automorphisms.
\end{remark}
\begin{proof}[Proof sketch of the theorem/proposition on this sheet]
The argument is the standard Yoneda-and-descent argument specialized to dual complexes.  A family over a test scheme $T$ corresponds, under representability, to a morphism $T\to M$.  Pulling the universal family back along this morphism recovers the given family.  If two morphisms give isomorphic pullbacks, the functorial equality on all test schemes forces the morphisms to agree.  The hypotheses on automorphisms ensure that the equality is not weakened to an equivalence relation in a groupoid.
\end{proof}
\begin{example}[A test calculation for dual complexes]
Let $T=\Spec k[s]$ and let the invariant of interest be denoted $\nu$.  A toy family $X\to T$ is chosen so that $\nu(X_t)=\nu_0$ for $s\ne0$ and the special fiber at $s=0$ records the boundary behavior.  In a concrete course one may take $X$ to be a family of divisors, modules, curves, bundles, or representations, depending on the chapter.  The moduli-theoretic lesson is that the morphism $T\to M$ detects both the general member and the limiting point.
\end{example}
\begin{exercise}
Write the functor of points for the moduli problem on this sheet and identify its value on $\Spec k$ and on $\Spec k[\epsilon]/(\epsilon^2)$.
\end{exercise}
\begin{exercise}
Formulate a small example in which the naive set of isomorphism classes is not Hausdorff, not separated, or not functorial.  Then repair the example by adding stability, rigidification, or a stack structure.
\end{exercise}
\clearpage
\section{Sheet 82: Tropical and logarithmic shadows: log smoothness}
This sheet studies \emph{log smoothness} as it appears in the broader chapter on \emph{tropical and logarithmic shadows}.  The recurring question is whether a proposed parameter space remembers enough information about families while remaining geometric enough to compute with.
\begin{definition}[Log smoothness viewpoint]
A log smoothness moduli functor assigns to each test scheme $T$ the set, groupoid, or category of families over $T$ satisfying the chosen condition.  Two families are identified only when the moduli problem says that an isomorphism, not merely a bijection on fibers, has been supplied.
\end{definition}
\begin{proposition}[Log smoothness viewpoint]
If a moduli functor for log smoothness is represented by a scheme $M$, then natural transformations from the functor to any functor $G$ are the same as elements of $G(M)$.  In practice this converts universal constructions on families into morphisms out of $M$.
\end{proposition}
\begin{example}[Log smoothness viewpoint]
For log smoothness, take a base $T=\Spec R$ and replace a single object by an $R$-flat family.  Flatness prevents the numerical invariants from jumping on first-order thickenings, so it is the minimal condition under which the fibers can reasonably be compared.
\end{example}
\begin{theorem}[Log smoothness viewpoint]
Suppose the objects in the log smoothness problem have no non-trivial automorphisms and admit effective descent for the chosen topology.  Then a universal family, if it exists, is unique up to unique isomorphism.  Consequently the representing object is unique up to unique isomorphism.
\end{theorem}
\begin{remark}[Log smoothness viewpoint]
The phrase ``the moduli space of log smoothness'' is therefore shorthand.  It may mean a fine moduli scheme, a coarse moduli space, an algebraic stack, a GIT quotient, or a compactification.  The correct meaning is determined by what information is retained about families and automorphisms.
\end{remark}
\begin{proof}[Proof sketch of the theorem/proposition on this sheet]
The argument is the standard Yoneda-and-descent argument specialized to log smoothness.  A family over a test scheme $T$ corresponds, under representability, to a morphism $T\to M$.  Pulling the universal family back along this morphism recovers the given family.  If two morphisms give isomorphic pullbacks, the functorial equality on all test schemes forces the morphisms to agree.  The hypotheses on automorphisms ensure that the equality is not weakened to an equivalence relation in a groupoid.
\end{proof}
\begin{example}[A test calculation for log smoothness]
Let $T=\Spec k[s]$ and let the invariant of interest be denoted $\nu$.  A toy family $X\to T$ is chosen so that $\nu(X_t)=\nu_0$ for $s\ne0$ and the special fiber at $s=0$ records the boundary behavior.  In a concrete course one may take $X$ to be a family of divisors, modules, curves, bundles, or representations, depending on the chapter.  The moduli-theoretic lesson is that the morphism $T\to M$ detects both the general member and the limiting point.
\end{example}
\begin{exercise}
Find one automorphism group occurring in the example and decide whether it obstructs the existence of a fine moduli scheme.
\end{exercise}
\begin{exercise}
Formulate a small example in which the naive set of isomorphism classes is not Hausdorff, not separated, or not functorial.  Then repair the example by adding stability, rigidification, or a stack structure.
\end{exercise}
\clearpage
\section{Sheet 83: Tropical and logarithmic shadows: tropical curves}
This sheet studies \emph{tropical curves} as it appears in the broader chapter on \emph{tropical and logarithmic shadows}.  The recurring question is whether a proposed parameter space remembers enough information about families while remaining geometric enough to compute with.
\begin{definition}[Tropical curves viewpoint]
A tropical curves moduli functor assigns to each test scheme $T$ the set, groupoid, or category of families over $T$ satisfying the chosen condition.  Two families are identified only when the moduli problem says that an isomorphism, not merely a bijection on fibers, has been supplied.
\end{definition}
\begin{proposition}[Tropical curves viewpoint]
If a moduli functor for tropical curves is represented by a scheme $M$, then natural transformations from the functor to any functor $G$ are the same as elements of $G(M)$.  In practice this converts universal constructions on families into morphisms out of $M$.
\end{proposition}
\begin{example}[Tropical curves viewpoint]
For tropical curves, take a base $T=\Spec R$ and replace a single object by an $R$-flat family.  Flatness prevents the numerical invariants from jumping on first-order thickenings, so it is the minimal condition under which the fibers can reasonably be compared.
\end{example}
\begin{theorem}[Tropical curves viewpoint]
Suppose the objects in the tropical curves problem have no non-trivial automorphisms and admit effective descent for the chosen topology.  Then a universal family, if it exists, is unique up to unique isomorphism.  Consequently the representing object is unique up to unique isomorphism.
\end{theorem}
\begin{remark}[Tropical curves viewpoint]
The phrase ``the moduli space of tropical curves'' is therefore shorthand.  It may mean a fine moduli scheme, a coarse moduli space, an algebraic stack, a GIT quotient, or a compactification.  The correct meaning is determined by what information is retained about families and automorphisms.
\end{remark}
\begin{proof}[Proof sketch of the theorem/proposition on this sheet]
The argument is the standard Yoneda-and-descent argument specialized to tropical curves.  A family over a test scheme $T$ corresponds, under representability, to a morphism $T\to M$.  Pulling the universal family back along this morphism recovers the given family.  If two morphisms give isomorphic pullbacks, the functorial equality on all test schemes forces the morphisms to agree.  The hypotheses on automorphisms ensure that the equality is not weakened to an equivalence relation in a groupoid.
\end{proof}
\begin{example}[A test calculation for tropical curves]
Let $T=\Spec k[s]$ and let the invariant of interest be denoted $\nu$.  A toy family $X\to T$ is chosen so that $\nu(X_t)=\nu_0$ for $s\ne0$ and the special fiber at $s=0$ records the boundary behavior.  In a concrete course one may take $X$ to be a family of divisors, modules, curves, bundles, or representations, depending on the chapter.  The moduli-theoretic lesson is that the morphism $T\to M$ detects both the general member and the limiting point.
\end{example}
\begin{exercise}
Check a base-change square explicitly: pull the displayed family back along $T\prime\to T$ and verify which invariant is preserved.
\end{exercise}
\begin{exercise}
Formulate a small example in which the naive set of isomorphism classes is not Hausdorff, not separated, or not functorial.  Then repair the example by adding stability, rigidification, or a stack structure.
\end{exercise}
\clearpage
\section{Sheet 84: Tropical and logarithmic shadows: expanded degenerations}
This sheet studies \emph{expanded degenerations} as it appears in the broader chapter on \emph{tropical and logarithmic shadows}.  The recurring question is whether a proposed parameter space remembers enough information about families while remaining geometric enough to compute with.
\begin{definition}[Expanded degenerations viewpoint]
A expanded degenerations moduli functor assigns to each test scheme $T$ the set, groupoid, or category of families over $T$ satisfying the chosen condition.  Two families are identified only when the moduli problem says that an isomorphism, not merely a bijection on fibers, has been supplied.
\end{definition}
\begin{proposition}[Expanded degenerations viewpoint]
If a moduli functor for expanded degenerations is represented by a scheme $M$, then natural transformations from the functor to any functor $G$ are the same as elements of $G(M)$.  In practice this converts universal constructions on families into morphisms out of $M$.
\end{proposition}
\begin{example}[Expanded degenerations viewpoint]
For expanded degenerations, take a base $T=\Spec R$ and replace a single object by an $R$-flat family.  Flatness prevents the numerical invariants from jumping on first-order thickenings, so it is the minimal condition under which the fibers can reasonably be compared.
\end{example}
\begin{theorem}[Expanded degenerations viewpoint]
Suppose the objects in the expanded degenerations problem have no non-trivial automorphisms and admit effective descent for the chosen topology.  Then a universal family, if it exists, is unique up to unique isomorphism.  Consequently the representing object is unique up to unique isomorphism.
\end{theorem}
\begin{remark}[Expanded degenerations viewpoint]
The phrase ``the moduli space of expanded degenerations'' is therefore shorthand.  It may mean a fine moduli scheme, a coarse moduli space, an algebraic stack, a GIT quotient, or a compactification.  The correct meaning is determined by what information is retained about families and automorphisms.
\end{remark}
\begin{proof}[Proof sketch of the theorem/proposition on this sheet]
The argument is the standard Yoneda-and-descent argument specialized to expanded degenerations.  A family over a test scheme $T$ corresponds, under representability, to a morphism $T\to M$.  Pulling the universal family back along this morphism recovers the given family.  If two morphisms give isomorphic pullbacks, the functorial equality on all test schemes forces the morphisms to agree.  The hypotheses on automorphisms ensure that the equality is not weakened to an equivalence relation in a groupoid.
\end{proof}
\begin{example}[A test calculation for expanded degenerations]
Let $T=\Spec k[s]$ and let the invariant of interest be denoted $\nu$.  A toy family $X\to T$ is chosen so that $\nu(X_t)=\nu_0$ for $s\ne0$ and the special fiber at $s=0$ records the boundary behavior.  In a concrete course one may take $X$ to be a family of divisors, modules, curves, bundles, or representations, depending on the chapter.  The moduli-theoretic lesson is that the morphism $T\to M$ detects both the general member and the limiting point.
\end{example}
\begin{exercise}
Construct a degeneration over $\Spec k[[t]]$ whose generic fiber satisfies the condition but whose special fiber forces a compactification.
\end{exercise}
\begin{exercise}
Formulate a small example in which the naive set of isomorphism classes is not Hausdorff, not separated, or not functorial.  Then repair the example by adding stability, rigidification, or a stack structure.
\end{exercise}
\clearpage
\section{Sheet 85: Tropical and logarithmic shadows: skeletons}
This sheet studies \emph{skeletons} as it appears in the broader chapter on \emph{tropical and logarithmic shadows}.  The recurring question is whether a proposed parameter space remembers enough information about families while remaining geometric enough to compute with.
\begin{definition}[Skeletons viewpoint]
A skeletons moduli functor assigns to each test scheme $T$ the set, groupoid, or category of families over $T$ satisfying the chosen condition.  Two families are identified only when the moduli problem says that an isomorphism, not merely a bijection on fibers, has been supplied.
\end{definition}
\begin{proposition}[Skeletons viewpoint]
If a moduli functor for skeletons is represented by a scheme $M$, then natural transformations from the functor to any functor $G$ are the same as elements of $G(M)$.  In practice this converts universal constructions on families into morphisms out of $M$.
\end{proposition}
\begin{example}[Skeletons viewpoint]
For skeletons, take a base $T=\Spec R$ and replace a single object by an $R$-flat family.  Flatness prevents the numerical invariants from jumping on first-order thickenings, so it is the minimal condition under which the fibers can reasonably be compared.
\end{example}
\begin{theorem}[Skeletons viewpoint]
Suppose the objects in the skeletons problem have no non-trivial automorphisms and admit effective descent for the chosen topology.  Then a universal family, if it exists, is unique up to unique isomorphism.  Consequently the representing object is unique up to unique isomorphism.
\end{theorem}
\begin{remark}[Skeletons viewpoint]
The phrase ``the moduli space of skeletons'' is therefore shorthand.  It may mean a fine moduli scheme, a coarse moduli space, an algebraic stack, a GIT quotient, or a compactification.  The correct meaning is determined by what information is retained about families and automorphisms.
\end{remark}
\begin{proof}[Proof sketch of the theorem/proposition on this sheet]
The argument is the standard Yoneda-and-descent argument specialized to skeletons.  A family over a test scheme $T$ corresponds, under representability, to a morphism $T\to M$.  Pulling the universal family back along this morphism recovers the given family.  If two morphisms give isomorphic pullbacks, the functorial equality on all test schemes forces the morphisms to agree.  The hypotheses on automorphisms ensure that the equality is not weakened to an equivalence relation in a groupoid.
\end{proof}
\begin{example}[A test calculation for skeletons]
Let $T=\Spec k[s]$ and let the invariant of interest be denoted $\nu$.  A toy family $X\to T$ is chosen so that $\nu(X_t)=\nu_0$ for $s\ne0$ and the special fiber at $s=0$ records the boundary behavior.  In a concrete course one may take $X$ to be a family of divisors, modules, curves, bundles, or representations, depending on the chapter.  The moduli-theoretic lesson is that the morphism $T\to M$ detects both the general member and the limiting point.
\end{example}
\begin{exercise}
Compare the functor with the corresponding groupoid-valued functor.  Which information is lost after passing to isomorphism classes?
\end{exercise}
\begin{exercise}
Formulate a small example in which the naive set of isomorphism classes is not Hausdorff, not separated, or not functorial.  Then repair the example by adding stability, rigidification, or a stack structure.
\end{exercise}
\clearpage
\section{Sheet 86: Computational examples: Groebner strata}
This sheet studies \emph{Groebner strata} as it appears in the broader chapter on \emph{computational examples}.  The recurring question is whether a proposed parameter space remembers enough information about families while remaining geometric enough to compute with.
\begin{definition}[Groebner strata viewpoint]
A Groebner strata moduli functor assigns to each test scheme $T$ the set, groupoid, or category of families over $T$ satisfying the chosen condition.  Two families are identified only when the moduli problem says that an isomorphism, not merely a bijection on fibers, has been supplied.
\end{definition}
\begin{proposition}[Groebner strata viewpoint]
If a moduli functor for Groebner strata is represented by a scheme $M$, then natural transformations from the functor to any functor $G$ are the same as elements of $G(M)$.  In practice this converts universal constructions on families into morphisms out of $M$.
\end{proposition}
\begin{example}[Groebner strata viewpoint]
For Groebner strata, take a base $T=\Spec R$ and replace a single object by an $R$-flat family.  Flatness prevents the numerical invariants from jumping on first-order thickenings, so it is the minimal condition under which the fibers can reasonably be compared.
\end{example}
\begin{theorem}[Groebner strata viewpoint]
Suppose the objects in the Groebner strata problem have no non-trivial automorphisms and admit effective descent for the chosen topology.  Then a universal family, if it exists, is unique up to unique isomorphism.  Consequently the representing object is unique up to unique isomorphism.
\end{theorem}
\begin{remark}[Groebner strata viewpoint]
The phrase ``the moduli space of Groebner strata'' is therefore shorthand.  It may mean a fine moduli scheme, a coarse moduli space, an algebraic stack, a GIT quotient, or a compactification.  The correct meaning is determined by what information is retained about families and automorphisms.
\end{remark}
\begin{proof}[Proof sketch of the theorem/proposition on this sheet]
The argument is the standard Yoneda-and-descent argument specialized to Groebner strata.  A family over a test scheme $T$ corresponds, under representability, to a morphism $T\to M$.  Pulling the universal family back along this morphism recovers the given family.  If two morphisms give isomorphic pullbacks, the functorial equality on all test schemes forces the morphisms to agree.  The hypotheses on automorphisms ensure that the equality is not weakened to an equivalence relation in a groupoid.
\end{proof}
\begin{example}[A test calculation for Groebner strata]
Let $T=\Spec k[s]$ and let the invariant of interest be denoted $\nu$.  A toy family $X\to T$ is chosen so that $\nu(X_t)=\nu_0$ for $s\ne0$ and the special fiber at $s=0$ records the boundary behavior.  In a concrete course one may take $X$ to be a family of divisors, modules, curves, bundles, or representations, depending on the chapter.  The moduli-theoretic lesson is that the morphism $T\to M$ detects both the general member and the limiting point.
\end{example}
\begin{exercise}
Write the functor of points for the moduli problem on this sheet and identify its value on $\Spec k$ and on $\Spec k[\epsilon]/(\epsilon^2)$.
\end{exercise}
\begin{exercise}
Formulate a small example in which the naive set of isomorphism classes is not Hausdorff, not separated, or not functorial.  Then repair the example by adding stability, rigidification, or a stack structure.
\end{exercise}
\clearpage
\section{Sheet 87: Computational examples: explicit Quot charts}
This sheet studies \emph{explicit Quot charts} as it appears in the broader chapter on \emph{computational examples}.  The recurring question is whether a proposed parameter space remembers enough information about families while remaining geometric enough to compute with.
\begin{definition}[Explicit quot charts viewpoint]
A explicit Quot charts moduli functor assigns to each test scheme $T$ the set, groupoid, or category of families over $T$ satisfying the chosen condition.  Two families are identified only when the moduli problem says that an isomorphism, not merely a bijection on fibers, has been supplied.
\end{definition}
\begin{proposition}[Explicit quot charts viewpoint]
If a moduli functor for explicit Quot charts is represented by a scheme $M$, then natural transformations from the functor to any functor $G$ are the same as elements of $G(M)$.  In practice this converts universal constructions on families into morphisms out of $M$.
\end{proposition}
\begin{example}[Explicit quot charts viewpoint]
For explicit Quot charts, take a base $T=\Spec R$ and replace a single object by an $R$-flat family.  Flatness prevents the numerical invariants from jumping on first-order thickenings, so it is the minimal condition under which the fibers can reasonably be compared.
\end{example}
\begin{theorem}[Explicit quot charts viewpoint]
Suppose the objects in the explicit Quot charts problem have no non-trivial automorphisms and admit effective descent for the chosen topology.  Then a universal family, if it exists, is unique up to unique isomorphism.  Consequently the representing object is unique up to unique isomorphism.
\end{theorem}
\begin{remark}[Explicit quot charts viewpoint]
The phrase ``the moduli space of explicit Quot charts'' is therefore shorthand.  It may mean a fine moduli scheme, a coarse moduli space, an algebraic stack, a GIT quotient, or a compactification.  The correct meaning is determined by what information is retained about families and automorphisms.
\end{remark}
\begin{proof}[Proof sketch of the theorem/proposition on this sheet]
The argument is the standard Yoneda-and-descent argument specialized to explicit Quot charts.  A family over a test scheme $T$ corresponds, under representability, to a morphism $T\to M$.  Pulling the universal family back along this morphism recovers the given family.  If two morphisms give isomorphic pullbacks, the functorial equality on all test schemes forces the morphisms to agree.  The hypotheses on automorphisms ensure that the equality is not weakened to an equivalence relation in a groupoid.
\end{proof}
\begin{example}[A test calculation for explicit Quot charts]
Let $T=\Spec k[s]$ and let the invariant of interest be denoted $\nu$.  A toy family $X\to T$ is chosen so that $\nu(X_t)=\nu_0$ for $s\ne0$ and the special fiber at $s=0$ records the boundary behavior.  In a concrete course one may take $X$ to be a family of divisors, modules, curves, bundles, or representations, depending on the chapter.  The moduli-theoretic lesson is that the morphism $T\to M$ detects both the general member and the limiting point.
\end{example}
\begin{exercise}
Find one automorphism group occurring in the example and decide whether it obstructs the existence of a fine moduli scheme.
\end{exercise}
\begin{exercise}
Formulate a small example in which the naive set of isomorphism classes is not Hausdorff, not separated, or not functorial.  Then repair the example by adding stability, rigidification, or a stack structure.
\end{exercise}
\clearpage
\section{Sheet 88: Computational examples: quiver GIT}
This sheet studies \emph{quiver GIT} as it appears in the broader chapter on \emph{computational examples}.  The recurring question is whether a proposed parameter space remembers enough information about families while remaining geometric enough to compute with.
\begin{definition}[Quiver git viewpoint]
A quiver GIT moduli functor assigns to each test scheme $T$ the set, groupoid, or category of families over $T$ satisfying the chosen condition.  Two families are identified only when the moduli problem says that an isomorphism, not merely a bijection on fibers, has been supplied.
\end{definition}
\begin{proposition}[Quiver git viewpoint]
If a moduli functor for quiver GIT is represented by a scheme $M$, then natural transformations from the functor to any functor $G$ are the same as elements of $G(M)$.  In practice this converts universal constructions on families into morphisms out of $M$.
\end{proposition}
\begin{example}[Quiver git viewpoint]
For quiver GIT, take a base $T=\Spec R$ and replace a single object by an $R$-flat family.  Flatness prevents the numerical invariants from jumping on first-order thickenings, so it is the minimal condition under which the fibers can reasonably be compared.
\end{example}
\begin{theorem}[Quiver git viewpoint]
Suppose the objects in the quiver GIT problem have no non-trivial automorphisms and admit effective descent for the chosen topology.  Then a universal family, if it exists, is unique up to unique isomorphism.  Consequently the representing object is unique up to unique isomorphism.
\end{theorem}
\begin{remark}[Quiver git viewpoint]
The phrase ``the moduli space of quiver GIT'' is therefore shorthand.  It may mean a fine moduli scheme, a coarse moduli space, an algebraic stack, a GIT quotient, or a compactification.  The correct meaning is determined by what information is retained about families and automorphisms.
\end{remark}
\begin{proof}[Proof sketch of the theorem/proposition on this sheet]
The argument is the standard Yoneda-and-descent argument specialized to quiver GIT.  A family over a test scheme $T$ corresponds, under representability, to a morphism $T\to M$.  Pulling the universal family back along this morphism recovers the given family.  If two morphisms give isomorphic pullbacks, the functorial equality on all test schemes forces the morphisms to agree.  The hypotheses on automorphisms ensure that the equality is not weakened to an equivalence relation in a groupoid.
\end{proof}
\begin{example}[A test calculation for quiver GIT]
Let $T=\Spec k[s]$ and let the invariant of interest be denoted $\nu$.  A toy family $X\to T$ is chosen so that $\nu(X_t)=\nu_0$ for $s\ne0$ and the special fiber at $s=0$ records the boundary behavior.  In a concrete course one may take $X$ to be a family of divisors, modules, curves, bundles, or representations, depending on the chapter.  The moduli-theoretic lesson is that the morphism $T\to M$ detects both the general member and the limiting point.
\end{example}
\begin{exercise}
Check a base-change square explicitly: pull the displayed family back along $T\prime\to T$ and verify which invariant is preserved.
\end{exercise}
\begin{exercise}
Formulate a small example in which the naive set of isomorphism classes is not Hausdorff, not separated, or not functorial.  Then repair the example by adding stability, rigidification, or a stack structure.
\end{exercise}
\clearpage
\section{Sheet 89: Computational examples: wall crossing examples}
This sheet studies \emph{wall crossing examples} as it appears in the broader chapter on \emph{computational examples}.  The recurring question is whether a proposed parameter space remembers enough information about families while remaining geometric enough to compute with.
\begin{definition}[Wall crossing examples viewpoint]
A wall crossing examples moduli functor assigns to each test scheme $T$ the set, groupoid, or category of families over $T$ satisfying the chosen condition.  Two families are identified only when the moduli problem says that an isomorphism, not merely a bijection on fibers, has been supplied.
\end{definition}
\begin{proposition}[Wall crossing examples viewpoint]
If a moduli functor for wall crossing examples is represented by a scheme $M$, then natural transformations from the functor to any functor $G$ are the same as elements of $G(M)$.  In practice this converts universal constructions on families into morphisms out of $M$.
\end{proposition}
\begin{example}[Wall crossing examples viewpoint]
For wall crossing examples, take a base $T=\Spec R$ and replace a single object by an $R$-flat family.  Flatness prevents the numerical invariants from jumping on first-order thickenings, so it is the minimal condition under which the fibers can reasonably be compared.
\end{example}
\begin{theorem}[Wall crossing examples viewpoint]
Suppose the objects in the wall crossing examples problem have no non-trivial automorphisms and admit effective descent for the chosen topology.  Then a universal family, if it exists, is unique up to unique isomorphism.  Consequently the representing object is unique up to unique isomorphism.
\end{theorem}
\begin{remark}[Wall crossing examples viewpoint]
The phrase ``the moduli space of wall crossing examples'' is therefore shorthand.  It may mean a fine moduli scheme, a coarse moduli space, an algebraic stack, a GIT quotient, or a compactification.  The correct meaning is determined by what information is retained about families and automorphisms.
\end{remark}
\begin{proof}[Proof sketch of the theorem/proposition on this sheet]
The argument is the standard Yoneda-and-descent argument specialized to wall crossing examples.  A family over a test scheme $T$ corresponds, under representability, to a morphism $T\to M$.  Pulling the universal family back along this morphism recovers the given family.  If two morphisms give isomorphic pullbacks, the functorial equality on all test schemes forces the morphisms to agree.  The hypotheses on automorphisms ensure that the equality is not weakened to an equivalence relation in a groupoid.
\end{proof}
\begin{example}[A test calculation for wall crossing examples]
Let $T=\Spec k[s]$ and let the invariant of interest be denoted $\nu$.  A toy family $X\to T$ is chosen so that $\nu(X_t)=\nu_0$ for $s\ne0$ and the special fiber at $s=0$ records the boundary behavior.  In a concrete course one may take $X$ to be a family of divisors, modules, curves, bundles, or representations, depending on the chapter.  The moduli-theoretic lesson is that the morphism $T\to M$ detects both the general member and the limiting point.
\end{example}
\begin{exercise}
Construct a degeneration over $\Spec k[[t]]$ whose generic fiber satisfies the condition but whose special fiber forces a compactification.
\end{exercise}
\begin{exercise}
Formulate a small example in which the naive set of isomorphism classes is not Hausdorff, not separated, or not functorial.  Then repair the example by adding stability, rigidification, or a stack structure.
\end{exercise}
\clearpage
\section{Sheet 90: Computational examples: Macaulay-style calculations}
This sheet studies \emph{Macaulay-style calculations} as it appears in the broader chapter on \emph{computational examples}.  The recurring question is whether a proposed parameter space remembers enough information about families while remaining geometric enough to compute with.
\begin{definition}[Macaulay-style calculations viewpoint]
A Macaulay-style calculations moduli functor assigns to each test scheme $T$ the set, groupoid, or category of families over $T$ satisfying the chosen condition.  Two families are identified only when the moduli problem says that an isomorphism, not merely a bijection on fibers, has been supplied.
\end{definition}
\begin{proposition}[Macaulay-style calculations viewpoint]
If a moduli functor for Macaulay-style calculations is represented by a scheme $M$, then natural transformations from the functor to any functor $G$ are the same as elements of $G(M)$.  In practice this converts universal constructions on families into morphisms out of $M$.
\end{proposition}
\begin{example}[Macaulay-style calculations viewpoint]
For Macaulay-style calculations, take a base $T=\Spec R$ and replace a single object by an $R$-flat family.  Flatness prevents the numerical invariants from jumping on first-order thickenings, so it is the minimal condition under which the fibers can reasonably be compared.
\end{example}
\begin{theorem}[Macaulay-style calculations viewpoint]
Suppose the objects in the Macaulay-style calculations problem have no non-trivial automorphisms and admit effective descent for the chosen topology.  Then a universal family, if it exists, is unique up to unique isomorphism.  Consequently the representing object is unique up to unique isomorphism.
\end{theorem}
\begin{remark}[Macaulay-style calculations viewpoint]
The phrase ``the moduli space of Macaulay-style calculations'' is therefore shorthand.  It may mean a fine moduli scheme, a coarse moduli space, an algebraic stack, a GIT quotient, or a compactification.  The correct meaning is determined by what information is retained about families and automorphisms.
\end{remark}
\begin{proof}[Proof sketch of the theorem/proposition on this sheet]
The argument is the standard Yoneda-and-descent argument specialized to Macaulay-style calculations.  A family over a test scheme $T$ corresponds, under representability, to a morphism $T\to M$.  Pulling the universal family back along this morphism recovers the given family.  If two morphisms give isomorphic pullbacks, the functorial equality on all test schemes forces the morphisms to agree.  The hypotheses on automorphisms ensure that the equality is not weakened to an equivalence relation in a groupoid.
\end{proof}
\begin{example}[A test calculation for Macaulay-style calculations]
Let $T=\Spec k[s]$ and let the invariant of interest be denoted $\nu$.  A toy family $X\to T$ is chosen so that $\nu(X_t)=\nu_0$ for $s\ne0$ and the special fiber at $s=0$ records the boundary behavior.  In a concrete course one may take $X$ to be a family of divisors, modules, curves, bundles, or representations, depending on the chapter.  The moduli-theoretic lesson is that the morphism $T\to M$ detects both the general member and the limiting point.
\end{example}
\begin{exercise}
Compare the functor with the corresponding groupoid-valued functor.  Which information is lost after passing to isomorphism classes?
\end{exercise}
\begin{exercise}
Formulate a small example in which the naive set of isomorphism classes is not Hausdorff, not separated, or not functorial.  Then repair the example by adding stability, rigidification, or a stack structure.
\end{exercise}
\clearpage
\section{Sheet 91: Exercises and projects: proof-writing drills}
This sheet studies \emph{proof-writing drills} as it appears in the broader chapter on \emph{exercises and projects}.  The recurring question is whether a proposed parameter space remembers enough information about families while remaining geometric enough to compute with.
\begin{definition}[Proof-writing drills viewpoint]
A proof-writing drills moduli functor assigns to each test scheme $T$ the set, groupoid, or category of families over $T$ satisfying the chosen condition.  Two families are identified only when the moduli problem says that an isomorphism, not merely a bijection on fibers, has been supplied.
\end{definition}
\begin{proposition}[Proof-writing drills viewpoint]
If a moduli functor for proof-writing drills is represented by a scheme $M$, then natural transformations from the functor to any functor $G$ are the same as elements of $G(M)$.  In practice this converts universal constructions on families into morphisms out of $M$.
\end{proposition}
\begin{example}[Proof-writing drills viewpoint]
For proof-writing drills, take a base $T=\Spec R$ and replace a single object by an $R$-flat family.  Flatness prevents the numerical invariants from jumping on first-order thickenings, so it is the minimal condition under which the fibers can reasonably be compared.
\end{example}
\begin{theorem}[Proof-writing drills viewpoint]
Suppose the objects in the proof-writing drills problem have no non-trivial automorphisms and admit effective descent for the chosen topology.  Then a universal family, if it exists, is unique up to unique isomorphism.  Consequently the representing object is unique up to unique isomorphism.
\end{theorem}
\begin{remark}[Proof-writing drills viewpoint]
The phrase ``the moduli space of proof-writing drills'' is therefore shorthand.  It may mean a fine moduli scheme, a coarse moduli space, an algebraic stack, a GIT quotient, or a compactification.  The correct meaning is determined by what information is retained about families and automorphisms.
\end{remark}
\begin{proof}[Proof sketch of the theorem/proposition on this sheet]
The argument is the standard Yoneda-and-descent argument specialized to proof-writing drills.  A family over a test scheme $T$ corresponds, under representability, to a morphism $T\to M$.  Pulling the universal family back along this morphism recovers the given family.  If two morphisms give isomorphic pullbacks, the functorial equality on all test schemes forces the morphisms to agree.  The hypotheses on automorphisms ensure that the equality is not weakened to an equivalence relation in a groupoid.
\end{proof}
\begin{example}[A test calculation for proof-writing drills]
Let $T=\Spec k[s]$ and let the invariant of interest be denoted $\nu$.  A toy family $X\to T$ is chosen so that $\nu(X_t)=\nu_0$ for $s\ne0$ and the special fiber at $s=0$ records the boundary behavior.  In a concrete course one may take $X$ to be a family of divisors, modules, curves, bundles, or representations, depending on the chapter.  The moduli-theoretic lesson is that the morphism $T\to M$ detects both the general member and the limiting point.
\end{example}
\begin{exercise}
Write the functor of points for the moduli problem on this sheet and identify its value on $\Spec k$ and on $\Spec k[\epsilon]/(\epsilon^2)$.
\end{exercise}
\begin{exercise}
Formulate a small example in which the naive set of isomorphism classes is not Hausdorff, not separated, or not functorial.  Then repair the example by adding stability, rigidification, or a stack structure.
\end{exercise}
\clearpage
\section{Sheet 92: Exercises and projects: examples to compute}
This sheet studies \emph{examples to compute} as it appears in the broader chapter on \emph{exercises and projects}.  The recurring question is whether a proposed parameter space remembers enough information about families while remaining geometric enough to compute with.
\begin{definition}[Examples to compute viewpoint]
A examples to compute moduli functor assigns to each test scheme $T$ the set, groupoid, or category of families over $T$ satisfying the chosen condition.  Two families are identified only when the moduli problem says that an isomorphism, not merely a bijection on fibers, has been supplied.
\end{definition}
\begin{proposition}[Examples to compute viewpoint]
If a moduli functor for examples to compute is represented by a scheme $M$, then natural transformations from the functor to any functor $G$ are the same as elements of $G(M)$.  In practice this converts universal constructions on families into morphisms out of $M$.
\end{proposition}
\begin{example}[Examples to compute viewpoint]
For examples to compute, take a base $T=\Spec R$ and replace a single object by an $R$-flat family.  Flatness prevents the numerical invariants from jumping on first-order thickenings, so it is the minimal condition under which the fibers can reasonably be compared.
\end{example}
\begin{theorem}[Examples to compute viewpoint]
Suppose the objects in the examples to compute problem have no non-trivial automorphisms and admit effective descent for the chosen topology.  Then a universal family, if it exists, is unique up to unique isomorphism.  Consequently the representing object is unique up to unique isomorphism.
\end{theorem}
\begin{remark}[Examples to compute viewpoint]
The phrase ``the moduli space of examples to compute'' is therefore shorthand.  It may mean a fine moduli scheme, a coarse moduli space, an algebraic stack, a GIT quotient, or a compactification.  The correct meaning is determined by what information is retained about families and automorphisms.
\end{remark}
\begin{proof}[Proof sketch of the theorem/proposition on this sheet]
The argument is the standard Yoneda-and-descent argument specialized to examples to compute.  A family over a test scheme $T$ corresponds, under representability, to a morphism $T\to M$.  Pulling the universal family back along this morphism recovers the given family.  If two morphisms give isomorphic pullbacks, the functorial equality on all test schemes forces the morphisms to agree.  The hypotheses on automorphisms ensure that the equality is not weakened to an equivalence relation in a groupoid.
\end{proof}
\begin{example}[A test calculation for examples to compute]
Let $T=\Spec k[s]$ and let the invariant of interest be denoted $\nu$.  A toy family $X\to T$ is chosen so that $\nu(X_t)=\nu_0$ for $s\ne0$ and the special fiber at $s=0$ records the boundary behavior.  In a concrete course one may take $X$ to be a family of divisors, modules, curves, bundles, or representations, depending on the chapter.  The moduli-theoretic lesson is that the morphism $T\to M$ detects both the general member and the limiting point.
\end{example}
\begin{exercise}
Find one automorphism group occurring in the example and decide whether it obstructs the existence of a fine moduli scheme.
\end{exercise}
\begin{exercise}
Formulate a small example in which the naive set of isomorphism classes is not Hausdorff, not separated, or not functorial.  Then repair the example by adding stability, rigidification, or a stack structure.
\end{exercise}
\clearpage
\section{Sheet 93: Exercises and projects: counterexamples}
This sheet studies \emph{counterexamples} as it appears in the broader chapter on \emph{exercises and projects}.  The recurring question is whether a proposed parameter space remembers enough information about families while remaining geometric enough to compute with.
\begin{definition}[Counterexamples viewpoint]
A counterexamples moduli functor assigns to each test scheme $T$ the set, groupoid, or category of families over $T$ satisfying the chosen condition.  Two families are identified only when the moduli problem says that an isomorphism, not merely a bijection on fibers, has been supplied.
\end{definition}
\begin{proposition}[Counterexamples viewpoint]
If a moduli functor for counterexamples is represented by a scheme $M$, then natural transformations from the functor to any functor $G$ are the same as elements of $G(M)$.  In practice this converts universal constructions on families into morphisms out of $M$.
\end{proposition}
\begin{example}[Counterexamples viewpoint]
For counterexamples, take a base $T=\Spec R$ and replace a single object by an $R$-flat family.  Flatness prevents the numerical invariants from jumping on first-order thickenings, so it is the minimal condition under which the fibers can reasonably be compared.
\end{example}
\begin{theorem}[Counterexamples viewpoint]
Suppose the objects in the counterexamples problem have no non-trivial automorphisms and admit effective descent for the chosen topology.  Then a universal family, if it exists, is unique up to unique isomorphism.  Consequently the representing object is unique up to unique isomorphism.
\end{theorem}
\begin{remark}[Counterexamples viewpoint]
The phrase ``the moduli space of counterexamples'' is therefore shorthand.  It may mean a fine moduli scheme, a coarse moduli space, an algebraic stack, a GIT quotient, or a compactification.  The correct meaning is determined by what information is retained about families and automorphisms.
\end{remark}
\begin{proof}[Proof sketch of the theorem/proposition on this sheet]
The argument is the standard Yoneda-and-descent argument specialized to counterexamples.  A family over a test scheme $T$ corresponds, under representability, to a morphism $T\to M$.  Pulling the universal family back along this morphism recovers the given family.  If two morphisms give isomorphic pullbacks, the functorial equality on all test schemes forces the morphisms to agree.  The hypotheses on automorphisms ensure that the equality is not weakened to an equivalence relation in a groupoid.
\end{proof}
\begin{example}[A test calculation for counterexamples]
Let $T=\Spec k[s]$ and let the invariant of interest be denoted $\nu$.  A toy family $X\to T$ is chosen so that $\nu(X_t)=\nu_0$ for $s\ne0$ and the special fiber at $s=0$ records the boundary behavior.  In a concrete course one may take $X$ to be a family of divisors, modules, curves, bundles, or representations, depending on the chapter.  The moduli-theoretic lesson is that the morphism $T\to M$ detects both the general member and the limiting point.
\end{example}
\begin{exercise}
Check a base-change square explicitly: pull the displayed family back along $T\prime\to T$ and verify which invariant is preserved.
\end{exercise}
\begin{exercise}
Formulate a small example in which the naive set of isomorphism classes is not Hausdorff, not separated, or not functorial.  Then repair the example by adding stability, rigidification, or a stack structure.
\end{exercise}
\clearpage
\section{Sheet 94: Exercises and projects: reading projects}
This sheet studies \emph{reading projects} as it appears in the broader chapter on \emph{exercises and projects}.  The recurring question is whether a proposed parameter space remembers enough information about families while remaining geometric enough to compute with.
\begin{definition}[Reading projects viewpoint]
A reading projects moduli functor assigns to each test scheme $T$ the set, groupoid, or category of families over $T$ satisfying the chosen condition.  Two families are identified only when the moduli problem says that an isomorphism, not merely a bijection on fibers, has been supplied.
\end{definition}
\begin{proposition}[Reading projects viewpoint]
If a moduli functor for reading projects is represented by a scheme $M$, then natural transformations from the functor to any functor $G$ are the same as elements of $G(M)$.  In practice this converts universal constructions on families into morphisms out of $M$.
\end{proposition}
\begin{example}[Reading projects viewpoint]
For reading projects, take a base $T=\Spec R$ and replace a single object by an $R$-flat family.  Flatness prevents the numerical invariants from jumping on first-order thickenings, so it is the minimal condition under which the fibers can reasonably be compared.
\end{example}
\begin{theorem}[Reading projects viewpoint]
Suppose the objects in the reading projects problem have no non-trivial automorphisms and admit effective descent for the chosen topology.  Then a universal family, if it exists, is unique up to unique isomorphism.  Consequently the representing object is unique up to unique isomorphism.
\end{theorem}
\begin{remark}[Reading projects viewpoint]
The phrase ``the moduli space of reading projects'' is therefore shorthand.  It may mean a fine moduli scheme, a coarse moduli space, an algebraic stack, a GIT quotient, or a compactification.  The correct meaning is determined by what information is retained about families and automorphisms.
\end{remark}
\begin{proof}[Proof sketch of the theorem/proposition on this sheet]
The argument is the standard Yoneda-and-descent argument specialized to reading projects.  A family over a test scheme $T$ corresponds, under representability, to a morphism $T\to M$.  Pulling the universal family back along this morphism recovers the given family.  If two morphisms give isomorphic pullbacks, the functorial equality on all test schemes forces the morphisms to agree.  The hypotheses on automorphisms ensure that the equality is not weakened to an equivalence relation in a groupoid.
\end{proof}
\begin{example}[A test calculation for reading projects]
Let $T=\Spec k[s]$ and let the invariant of interest be denoted $\nu$.  A toy family $X\to T$ is chosen so that $\nu(X_t)=\nu_0$ for $s\ne0$ and the special fiber at $s=0$ records the boundary behavior.  In a concrete course one may take $X$ to be a family of divisors, modules, curves, bundles, or representations, depending on the chapter.  The moduli-theoretic lesson is that the morphism $T\to M$ detects both the general member and the limiting point.
\end{example}
\begin{exercise}
Construct a degeneration over $\Spec k[[t]]$ whose generic fiber satisfies the condition but whose special fiber forces a compactification.
\end{exercise}
\begin{exercise}
Formulate a small example in which the naive set of isomorphism classes is not Hausdorff, not separated, or not functorial.  Then repair the example by adding stability, rigidification, or a stack structure.
\end{exercise}
\clearpage
\section{Sheet 95: Exercises and projects: research prompts}
This sheet studies \emph{research prompts} as it appears in the broader chapter on \emph{exercises and projects}.  The recurring question is whether a proposed parameter space remembers enough information about families while remaining geometric enough to compute with.
\begin{definition}[Research prompts viewpoint]
A research prompts moduli functor assigns to each test scheme $T$ the set, groupoid, or category of families over $T$ satisfying the chosen condition.  Two families are identified only when the moduli problem says that an isomorphism, not merely a bijection on fibers, has been supplied.
\end{definition}
\begin{proposition}[Research prompts viewpoint]
If a moduli functor for research prompts is represented by a scheme $M$, then natural transformations from the functor to any functor $G$ are the same as elements of $G(M)$.  In practice this converts universal constructions on families into morphisms out of $M$.
\end{proposition}
\begin{example}[Research prompts viewpoint]
For research prompts, take a base $T=\Spec R$ and replace a single object by an $R$-flat family.  Flatness prevents the numerical invariants from jumping on first-order thickenings, so it is the minimal condition under which the fibers can reasonably be compared.
\end{example}
\begin{theorem}[Research prompts viewpoint]
Suppose the objects in the research prompts problem have no non-trivial automorphisms and admit effective descent for the chosen topology.  Then a universal family, if it exists, is unique up to unique isomorphism.  Consequently the representing object is unique up to unique isomorphism.
\end{theorem}
\begin{remark}[Research prompts viewpoint]
The phrase ``the moduli space of research prompts'' is therefore shorthand.  It may mean a fine moduli scheme, a coarse moduli space, an algebraic stack, a GIT quotient, or a compactification.  The correct meaning is determined by what information is retained about families and automorphisms.
\end{remark}
\begin{proof}[Proof sketch of the theorem/proposition on this sheet]
The argument is the standard Yoneda-and-descent argument specialized to research prompts.  A family over a test scheme $T$ corresponds, under representability, to a morphism $T\to M$.  Pulling the universal family back along this morphism recovers the given family.  If two morphisms give isomorphic pullbacks, the functorial equality on all test schemes forces the morphisms to agree.  The hypotheses on automorphisms ensure that the equality is not weakened to an equivalence relation in a groupoid.
\end{proof}
\begin{example}[A test calculation for research prompts]
Let $T=\Spec k[s]$ and let the invariant of interest be denoted $\nu$.  A toy family $X\to T$ is chosen so that $\nu(X_t)=\nu_0$ for $s\ne0$ and the special fiber at $s=0$ records the boundary behavior.  In a concrete course one may take $X$ to be a family of divisors, modules, curves, bundles, or representations, depending on the chapter.  The moduli-theoretic lesson is that the morphism $T\to M$ detects both the general member and the limiting point.
\end{example}
\begin{exercise}
Compare the functor with the corresponding groupoid-valued functor.  Which information is lost after passing to isomorphism classes?
\end{exercise}
\begin{exercise}
Formulate a small example in which the naive set of isomorphism classes is not Hausdorff, not separated, or not functorial.  Then repair the example by adding stability, rigidification, or a stack structure.
\end{exercise}
\clearpage
\section{Sheet 96: Epilogue: comparison of moduli problems}
This sheet studies \emph{comparison of moduli problems} as it appears in the broader chapter on \emph{epilogue}.  The recurring question is whether a proposed parameter space remembers enough information about families while remaining geometric enough to compute with.
\begin{definition}[Comparison of moduli problems viewpoint]
A comparison of moduli problems moduli functor assigns to each test scheme $T$ the set, groupoid, or category of families over $T$ satisfying the chosen condition.  Two families are identified only when the moduli problem says that an isomorphism, not merely a bijection on fibers, has been supplied.
\end{definition}
\begin{proposition}[Comparison of moduli problems viewpoint]
If a moduli functor for comparison of moduli problems is represented by a scheme $M$, then natural transformations from the functor to any functor $G$ are the same as elements of $G(M)$.  In practice this converts universal constructions on families into morphisms out of $M$.
\end{proposition}
\begin{example}[Comparison of moduli problems viewpoint]
For comparison of moduli problems, take a base $T=\Spec R$ and replace a single object by an $R$-flat family.  Flatness prevents the numerical invariants from jumping on first-order thickenings, so it is the minimal condition under which the fibers can reasonably be compared.
\end{example}
\begin{theorem}[Comparison of moduli problems viewpoint]
Suppose the objects in the comparison of moduli problems problem have no non-trivial automorphisms and admit effective descent for the chosen topology.  Then a universal family, if it exists, is unique up to unique isomorphism.  Consequently the representing object is unique up to unique isomorphism.
\end{theorem}
\begin{remark}[Comparison of moduli problems viewpoint]
The phrase ``the moduli space of comparison of moduli problems'' is therefore shorthand.  It may mean a fine moduli scheme, a coarse moduli space, an algebraic stack, a GIT quotient, or a compactification.  The correct meaning is determined by what information is retained about families and automorphisms.
\end{remark}
\begin{proof}[Proof sketch of the theorem/proposition on this sheet]
The argument is the standard Yoneda-and-descent argument specialized to comparison of moduli problems.  A family over a test scheme $T$ corresponds, under representability, to a morphism $T\to M$.  Pulling the universal family back along this morphism recovers the given family.  If two morphisms give isomorphic pullbacks, the functorial equality on all test schemes forces the morphisms to agree.  The hypotheses on automorphisms ensure that the equality is not weakened to an equivalence relation in a groupoid.
\end{proof}
\begin{example}[A test calculation for comparison of moduli problems]
Let $T=\Spec k[s]$ and let the invariant of interest be denoted $\nu$.  A toy family $X\to T$ is chosen so that $\nu(X_t)=\nu_0$ for $s\ne0$ and the special fiber at $s=0$ records the boundary behavior.  In a concrete course one may take $X$ to be a family of divisors, modules, curves, bundles, or representations, depending on the chapter.  The moduli-theoretic lesson is that the morphism $T\to M$ detects both the general member and the limiting point.
\end{example}
\begin{exercise}
Write the functor of points for the moduli problem on this sheet and identify its value on $\Spec k$ and on $\Spec k[\epsilon]/(\epsilon^2)$.
\end{exercise}
\begin{exercise}
Formulate a small example in which the naive set of isomorphism classes is not Hausdorff, not separated, or not functorial.  Then repair the example by adding stability, rigidification, or a stack structure.
\end{exercise}
\clearpage
\section{Sheet 97: Epilogue: dictionary of notions}
This sheet studies \emph{dictionary of notions} as it appears in the broader chapter on \emph{epilogue}.  The recurring question is whether a proposed parameter space remembers enough information about families while remaining geometric enough to compute with.
\begin{definition}[Dictionary of notions viewpoint]
A dictionary of notions moduli functor assigns to each test scheme $T$ the set, groupoid, or category of families over $T$ satisfying the chosen condition.  Two families are identified only when the moduli problem says that an isomorphism, not merely a bijection on fibers, has been supplied.
\end{definition}
\begin{proposition}[Dictionary of notions viewpoint]
If a moduli functor for dictionary of notions is represented by a scheme $M$, then natural transformations from the functor to any functor $G$ are the same as elements of $G(M)$.  In practice this converts universal constructions on families into morphisms out of $M$.
\end{proposition}
\begin{example}[Dictionary of notions viewpoint]
For dictionary of notions, take a base $T=\Spec R$ and replace a single object by an $R$-flat family.  Flatness prevents the numerical invariants from jumping on first-order thickenings, so it is the minimal condition under which the fibers can reasonably be compared.
\end{example}
\begin{theorem}[Dictionary of notions viewpoint]
Suppose the objects in the dictionary of notions problem have no non-trivial automorphisms and admit effective descent for the chosen topology.  Then a universal family, if it exists, is unique up to unique isomorphism.  Consequently the representing object is unique up to unique isomorphism.
\end{theorem}
\begin{remark}[Dictionary of notions viewpoint]
The phrase ``the moduli space of dictionary of notions'' is therefore shorthand.  It may mean a fine moduli scheme, a coarse moduli space, an algebraic stack, a GIT quotient, or a compactification.  The correct meaning is determined by what information is retained about families and automorphisms.
\end{remark}
\begin{proof}[Proof sketch of the theorem/proposition on this sheet]
The argument is the standard Yoneda-and-descent argument specialized to dictionary of notions.  A family over a test scheme $T$ corresponds, under representability, to a morphism $T\to M$.  Pulling the universal family back along this morphism recovers the given family.  If two morphisms give isomorphic pullbacks, the functorial equality on all test schemes forces the morphisms to agree.  The hypotheses on automorphisms ensure that the equality is not weakened to an equivalence relation in a groupoid.
\end{proof}
\begin{example}[A test calculation for dictionary of notions]
Let $T=\Spec k[s]$ and let the invariant of interest be denoted $\nu$.  A toy family $X\to T$ is chosen so that $\nu(X_t)=\nu_0$ for $s\ne0$ and the special fiber at $s=0$ records the boundary behavior.  In a concrete course one may take $X$ to be a family of divisors, modules, curves, bundles, or representations, depending on the chapter.  The moduli-theoretic lesson is that the morphism $T\to M$ detects both the general member and the limiting point.
\end{example}
\begin{exercise}
Find one automorphism group occurring in the example and decide whether it obstructs the existence of a fine moduli scheme.
\end{exercise}
\begin{exercise}
Formulate a small example in which the naive set of isomorphism classes is not Hausdorff, not separated, or not functorial.  Then repair the example by adding stability, rigidification, or a stack structure.
\end{exercise}
\clearpage
\section{Sheet 98: Epilogue: common pitfalls}
This sheet studies \emph{common pitfalls} as it appears in the broader chapter on \emph{epilogue}.  The recurring question is whether a proposed parameter space remembers enough information about families while remaining geometric enough to compute with.
\begin{definition}[Common pitfalls viewpoint]
A common pitfalls moduli functor assigns to each test scheme $T$ the set, groupoid, or category of families over $T$ satisfying the chosen condition.  Two families are identified only when the moduli problem says that an isomorphism, not merely a bijection on fibers, has been supplied.
\end{definition}
\begin{proposition}[Common pitfalls viewpoint]
If a moduli functor for common pitfalls is represented by a scheme $M$, then natural transformations from the functor to any functor $G$ are the same as elements of $G(M)$.  In practice this converts universal constructions on families into morphisms out of $M$.
\end{proposition}
\begin{example}[Common pitfalls viewpoint]
For common pitfalls, take a base $T=\Spec R$ and replace a single object by an $R$-flat family.  Flatness prevents the numerical invariants from jumping on first-order thickenings, so it is the minimal condition under which the fibers can reasonably be compared.
\end{example}
\begin{theorem}[Common pitfalls viewpoint]
Suppose the objects in the common pitfalls problem have no non-trivial automorphisms and admit effective descent for the chosen topology.  Then a universal family, if it exists, is unique up to unique isomorphism.  Consequently the representing object is unique up to unique isomorphism.
\end{theorem}
\begin{remark}[Common pitfalls viewpoint]
The phrase ``the moduli space of common pitfalls'' is therefore shorthand.  It may mean a fine moduli scheme, a coarse moduli space, an algebraic stack, a GIT quotient, or a compactification.  The correct meaning is determined by what information is retained about families and automorphisms.
\end{remark}
\begin{proof}[Proof sketch of the theorem/proposition on this sheet]
The argument is the standard Yoneda-and-descent argument specialized to common pitfalls.  A family over a test scheme $T$ corresponds, under representability, to a morphism $T\to M$.  Pulling the universal family back along this morphism recovers the given family.  If two morphisms give isomorphic pullbacks, the functorial equality on all test schemes forces the morphisms to agree.  The hypotheses on automorphisms ensure that the equality is not weakened to an equivalence relation in a groupoid.
\end{proof}
\begin{example}[A test calculation for common pitfalls]
Let $T=\Spec k[s]$ and let the invariant of interest be denoted $\nu$.  A toy family $X\to T$ is chosen so that $\nu(X_t)=\nu_0$ for $s\ne0$ and the special fiber at $s=0$ records the boundary behavior.  In a concrete course one may take $X$ to be a family of divisors, modules, curves, bundles, or representations, depending on the chapter.  The moduli-theoretic lesson is that the morphism $T\to M$ detects both the general member and the limiting point.
\end{example}
\begin{exercise}
Check a base-change square explicitly: pull the displayed family back along $T\prime\to T$ and verify which invariant is preserved.
\end{exercise}
\begin{exercise}
Formulate a small example in which the naive set of isomorphism classes is not Hausdorff, not separated, or not functorial.  Then repair the example by adding stability, rigidification, or a stack structure.
\end{exercise}
\clearpage
\section{Sheet 99: Epilogue: open-ended directions}
This sheet studies \emph{open-ended directions} as it appears in the broader chapter on \emph{epilogue}.  The recurring question is whether a proposed parameter space remembers enough information about families while remaining geometric enough to compute with.
\begin{definition}[Open-ended directions viewpoint]
A open-ended directions moduli functor assigns to each test scheme $T$ the set, groupoid, or category of families over $T$ satisfying the chosen condition.  Two families are identified only when the moduli problem says that an isomorphism, not merely a bijection on fibers, has been supplied.
\end{definition}
\begin{proposition}[Open-ended directions viewpoint]
If a moduli functor for open-ended directions is represented by a scheme $M$, then natural transformations from the functor to any functor $G$ are the same as elements of $G(M)$.  In practice this converts universal constructions on families into morphisms out of $M$.
\end{proposition}
\begin{example}[Open-ended directions viewpoint]
For open-ended directions, take a base $T=\Spec R$ and replace a single object by an $R$-flat family.  Flatness prevents the numerical invariants from jumping on first-order thickenings, so it is the minimal condition under which the fibers can reasonably be compared.
\end{example}
\begin{theorem}[Open-ended directions viewpoint]
Suppose the objects in the open-ended directions problem have no non-trivial automorphisms and admit effective descent for the chosen topology.  Then a universal family, if it exists, is unique up to unique isomorphism.  Consequently the representing object is unique up to unique isomorphism.
\end{theorem}
\begin{remark}[Open-ended directions viewpoint]
The phrase ``the moduli space of open-ended directions'' is therefore shorthand.  It may mean a fine moduli scheme, a coarse moduli space, an algebraic stack, a GIT quotient, or a compactification.  The correct meaning is determined by what information is retained about families and automorphisms.
\end{remark}
\begin{proof}[Proof sketch of the theorem/proposition on this sheet]
The argument is the standard Yoneda-and-descent argument specialized to open-ended directions.  A family over a test scheme $T$ corresponds, under representability, to a morphism $T\to M$.  Pulling the universal family back along this morphism recovers the given family.  If two morphisms give isomorphic pullbacks, the functorial equality on all test schemes forces the morphisms to agree.  The hypotheses on automorphisms ensure that the equality is not weakened to an equivalence relation in a groupoid.
\end{proof}
\begin{example}[A test calculation for open-ended directions]
Let $T=\Spec k[s]$ and let the invariant of interest be denoted $\nu$.  A toy family $X\to T$ is chosen so that $\nu(X_t)=\nu_0$ for $s\ne0$ and the special fiber at $s=0$ records the boundary behavior.  In a concrete course one may take $X$ to be a family of divisors, modules, curves, bundles, or representations, depending on the chapter.  The moduli-theoretic lesson is that the morphism $T\to M$ detects both the general member and the limiting point.
\end{example}
\begin{exercise}
Construct a degeneration over $\Spec k[[t]]$ whose generic fiber satisfies the condition but whose special fiber forces a compactification.
\end{exercise}
\begin{exercise}
Formulate a small example in which the naive set of isomorphism classes is not Hausdorff, not separated, or not functorial.  Then repair the example by adding stability, rigidification, or a stack structure.
\end{exercise}
\clearpage
\section{Sheet 100: Epilogue: summary table}
This sheet studies \emph{summary table} as it appears in the broader chapter on \emph{epilogue}.  The recurring question is whether a proposed parameter space remembers enough information about families while remaining geometric enough to compute with.
\begin{definition}[Summary table viewpoint]
A summary table moduli functor assigns to each test scheme $T$ the set, groupoid, or category of families over $T$ satisfying the chosen condition.  Two families are identified only when the moduli problem says that an isomorphism, not merely a bijection on fibers, has been supplied.
\end{definition}
\begin{proposition}[Summary table viewpoint]
If a moduli functor for summary table is represented by a scheme $M$, then natural transformations from the functor to any functor $G$ are the same as elements of $G(M)$.  In practice this converts universal constructions on families into morphisms out of $M$.
\end{proposition}
\begin{example}[Summary table viewpoint]
For summary table, take a base $T=\Spec R$ and replace a single object by an $R$-flat family.  Flatness prevents the numerical invariants from jumping on first-order thickenings, so it is the minimal condition under which the fibers can reasonably be compared.
\end{example}
\begin{theorem}[Summary table viewpoint]
Suppose the objects in the summary table problem have no non-trivial automorphisms and admit effective descent for the chosen topology.  Then a universal family, if it exists, is unique up to unique isomorphism.  Consequently the representing object is unique up to unique isomorphism.
\end{theorem}
\begin{remark}[Summary table viewpoint]
The phrase ``the moduli space of summary table'' is therefore shorthand.  It may mean a fine moduli scheme, a coarse moduli space, an algebraic stack, a GIT quotient, or a compactification.  The correct meaning is determined by what information is retained about families and automorphisms.
\end{remark}
\begin{proof}[Proof sketch of the theorem/proposition on this sheet]
The argument is the standard Yoneda-and-descent argument specialized to summary table.  A family over a test scheme $T$ corresponds, under representability, to a morphism $T\to M$.  Pulling the universal family back along this morphism recovers the given family.  If two morphisms give isomorphic pullbacks, the functorial equality on all test schemes forces the morphisms to agree.  The hypotheses on automorphisms ensure that the equality is not weakened to an equivalence relation in a groupoid.
\end{proof}
\begin{example}[A test calculation for summary table]
Let $T=\Spec k[s]$ and let the invariant of interest be denoted $\nu$.  A toy family $X\to T$ is chosen so that $\nu(X_t)=\nu_0$ for $s\ne0$ and the special fiber at $s=0$ records the boundary behavior.  In a concrete course one may take $X$ to be a family of divisors, modules, curves, bundles, or representations, depending on the chapter.  The moduli-theoretic lesson is that the morphism $T\to M$ detects both the general member and the limiting point.
\end{example}
\begin{exercise}
Compare the functor with the corresponding groupoid-valued functor.  Which information is lost after passing to isomorphism classes?
\end{exercise}
\begin{exercise}
Formulate a small example in which the naive set of isomorphism classes is not Hausdorff, not separated, or not functorial.  Then repair the example by adding stability, rigidification, or a stack structure.
\end{exercise}
\clearpage

\appendix
\section{Checklist for a new moduli problem}
\begin{enumerate}[label=\arabic*.]
\item Specify objects over every test scheme, not only over algebraically closed fields.
\item Specify morphisms between families and decide whether the target is a set, groupoid, or higher groupoid.
\item Prove base change compatibility.
\item Identify automorphism groups and tangent spaces.
\item Search for boundedness, openness of stability, separatedness, and properness.
\item Decide whether the final answer should be a fine moduli space, coarse space, stack, derived stack, or compactification.
\end{enumerate}
\section{Glossary}
A \emph{fine} moduli space represents the functor and carries a universal family.  A \emph{coarse} moduli space corepresents maps to schemes and has points corresponding to isomorphism classes over algebraically closed fields.  A \emph{stack} retains automorphisms and descent data.  A \emph{stability condition} removes or controls bad orbits.  A \emph{compactification} adds limiting objects so that properness becomes plausible.
\clearpage

\begin{thebibliography}{9}
\bibitem{GIT} D. Mumford, J. Fogarty, and F. Kirwan, \emph{Geometric Invariant Theory}, third edition.
\bibitem{HM} J. Harris and I. Morrison, \emph{Moduli of Curves}.
\bibitem{HL} D. Huybrechts and M. Lehn, \emph{The Geometry of Moduli Spaces of Sheaves}.
\bibitem{Olsson} M. Olsson, \emph{Algebraic Spaces and Stacks}.
\bibitem{Sernesi} E. Sernesi, \emph{Deformations of Algebraic Schemes}.
\bibitem{Vakil} R. Vakil, \emph{The Rising Sea: Foundations of Algebraic Geometry}, online notes.
\end{thebibliography}
\end{document}
